% NeurIPS 2026 anonymous-submission paper draft.
\documentclass{article}

% Pass options before loading NeurIPS style.
\PassOptionsToPackage{numbers,square,sort&compress}{natbib}
\PassOptionsToPackage{hypertexnames=false}{hyperref}

% Load NeurIPS style.
\usepackage{neurips_2026}

% Force numeric bracket citations if natbib is loaded by the style file.
\makeatletter
\@ifpackageloaded{natbib}{%
  \setcitestyle{numbers,square,sort&compress}%
}{}
\makeatother

% Colors + hyperlink setup.
\usepackage[dvipsnames]{xcolor}
\definecolor{Bleu}{RGB}{0,0,204}

% Load hyperref only if not already loaded.
\makeatletter
\@ifpackageloaded{hyperref}{}{%
  \usepackage{hyperref}%
}
\makeatother

\hypersetup{
  colorlinks=true,
  citecolor=Bleu,
  linkcolor=Bleu,
  urlcolor=Bleu,
  hypertexnames=false
}

% Packages. Do not load geometry or natbib separately; NeurIPS handles them.
\usepackage{amsmath,amssymb,amsthm}
\usepackage{mathtools}
\usepackage{bbm}
\usepackage{graphicx}
\usepackage{booktabs}
\usepackage{microtype}
\usepackage{enumitem}
\usepackage{subcaption}
\usepackage{changepage}
\usepackage{multirow}
\usepackage{setspace}
 
\usepackage{tabularx}
% Algorithms.
\usepackage{algorithm}
\usepackage{algorithmic}

% Macros.
\DeclareMathOperator*{\argmin}{arg\,min}
\DeclareMathOperator{\Var}{Var}
\newcommand{\E}{\mathbb{E}}
\newcommand{\slightspacing}{\setstretch{1.175}}
\newcommand{\ALGCOMMENT}[1]{\item[]\texttt{\#~#1}}
\newcommand{\BCINPUT}{\STATE \textbf{Input: }}
\newcommand{\BCOUTPUT}{\STATE \textbf{Output: }}

% Theorems.
\theoremstyle{plain}
\newtheorem{theorem}{Theorem}
\newtheorem{lemma}{Lemma}
\newtheorem{proposition}{Proposition}

\theoremstyle{definition}
\newtheorem{definition}{Definition}

\theoremstyle{remark}
\newtheorem{remark}{Remark}

\title{Bellman Calibration for $V$-Learning in Offline Reinforcement Learning}

\author{%
  Lars van der Laan \\
  Department of Statistics, University of Washington \\
  \texttt{lvdlaan@uw.edu} \\
  \And
  Nathan Kallus \\
  Netflix and Cornell University
}

\begin{document}

\maketitle

\begin{abstract}
Reliable long-horizon value prediction is difficult in offline reinforcement
learning because fitted value methods combine bootstrapping, function
approximation, and distribution shift, while standard guarantees often require
Bellman completeness or realizability. We introduce \emph{Bellman calibration},
a weak reliability criterion requiring that states assigned similar predicted
values have average Bellman targets that agree with those predictions. This
criterion yields a scalar calibration error for diagnosing systematic numerical
miscalibration, which we estimate from off-policy data using doubly robust
Bellman target estimates. We then propose \emph{Iterated Bellman Calibration}, a
model-agnostic post-hoc procedure that recalibrates any learned value predictor
by fitting a one-dimensional map of its original prediction, with histogram and
isotonic variants. We prove finite-sample guarantees showing that Bellman
calibration error is controlled at one-dimensional nonparametric rates without
Bellman completeness or value-function realizability. Our value-error bounds
separate statistical estimation, finite-iteration, and approximation errors, clarifying when calibration improves value prediction and when
its gains are limited by the information in the original predictor or insufficient coverage.
\end{abstract}
\section{Introduction}

Many applications require predicting the long-term consequences of a decision
policy in a sequential, stochastic environment. We study this problem in a
Markov decision process (MDP), where the goal is to estimate the long-run value
of a target policy~\(\pi\), possibly from data collected under a different
behavior policy. This off-policy value-prediction problem arises throughout RL
and in many application domains: clinicians anticipate long-term outcomes under
proposed treatment rules \citep{van2019calibration}; online platforms estimate
customer lifetime value and retention under alternative recommendation
strategies \citep{theocharous2015ad,maystre2022temporally,xue2025auro}; and
economists study downstream effects of counterfactual programs
\citep{rust1987optimal,cowgill2019economics}. In these settings, systematic
miscalibration can distort policy comparisons, misstate long-run effects, and
lead to poor decisions.

Modern predictors such as neural networks and gradient-boosted trees are often
miscalibrated: their numerical predictions may systematically differ from the
quantities they are intended to estimate because of misspecification, limited
data, or distribution shift
\citep{zadrozny2001obtaining,niculescu2005predicting,bella2010calibration,guo2017calibration}.
For long-horizon value prediction, this concern is amplified by bootstrapping:
fitted value methods regress on targets that themselves depend on previous
value estimates, so finite-iteration error, approximation error, and off-policy
shift can propagate into the final value scale
\citep{gordon1995stable,tsitsiklis1996analysis,munos2008finite,
farahmand2010error,agarwal2021deep}. Existing guarantees typically control
these errors by imposing strong structure, such as Bellman completeness,
realizability, linear approximation, or finite state spaces
\citep{chen2019information,xie2022role,meyn2024projected,xie2021batch}. Min--max
alternatives relax some aspects of fitted value iteration but introduce their
own challenges, including saddle-point optimization, rich critic or dual
classes, and strong coverage requirements
\citep{dai2018sbeed,uehara2020minimax,jin2021bellman,xie2021batch}.

Calibration is a classical reliability property: among units assigned the same
predicted value, the corresponding outcomes should agree with that value on
average. In population terms, a calibrated predictor satisfies
\(f(X) = E[Y \mid f(X)]\). This suggests an RL analogue. Informally, a value
predictor is calibrated if, among states assigned the same predicted value, the
average target-policy Bellman backup matches that value; that is,
\(v(S) = E_\pi\{R + \gamma v(S') \mid v(S)\}\). This is a weak but meaningful
form of Bellman consistency: it constrains only the scalar prediction rather
than the full state-to-value mapping, while still ruling out systematic under-
or overestimation within level sets of the prediction.

Despite its central role in supervised learning, where predictions are
calibrated against observed labels, empirical frequencies, or conditional means
\citep{lichtenstein1977calibration,platt1999probabilistic,zadrozny2001obtaining,
vovk2005algorithmic,gupta2020distribution,gupta2021distribution,
whitehouse2024orthogonal}, calibration remains comparatively underexplored for
long-horizon value prediction in RL. Table~\ref{tab:rl-calibration-motivation}
highlights why this is problematic: value estimates can be highly informative
yet numerically misaligned because of temporal shift, off-policy fitting, finite
Bellman iteration, or approximation and optimization error. Unlike in supervised
learning, however, the calibration target is not directly observed; long-run
returns must be inferred from one-step transitions through Bellman structure.
We therefore study post-hoc Bellman calibration: a restricted recalibration of a
pretrained value predictor that corrects systematic scale errors without a new
high-dimensional Bellman fit, with finite-sample guarantees that do not require
realizability or Bellman completeness.

\begin{table}[t]
\centering
\small
\setlength{\tabcolsep}{4pt}
\renewcommand{\arraystretch}{1.08}
\caption{Why RL value predictions can be informative but poorly calibrated.}
\label{tab:rl-calibration-motivation}
\begin{tabularx}{\linewidth}{p{0.47\linewidth} p{0.45\linewidth}}
\toprule
\textbf{Failure mode} & \textbf{Calibration correction} \\
\midrule

Temporal or distribution shift: value estimates reflect past rewards, policies,
or environments rather than the current target setting
\citep{kumar2020conservative}.
&
Update the value scale using recent or gold-standard trajectories, while
retaining signal from the historical model.
\\
\midrule

Off-policy fitting: a high-capacity value model trained on large
behavior-policy data can be informative but mis-scaled for the target-policy
state--action distribution \citep{fujimoto2019off}.
&
Use the large off-policy sample to learn the value predictor, then recalibrate its
scale on a smaller on-policy or target-policy reference sample.
\\
\midrule

Finite Bellman iteration: computational limits can force early stopping or a
shorter effective horizon, mis-scaling values for the intended return target
\citep{munos2008finite}.
&
Re-align predictions to the intended horizon without rerunning the full Bellman
iteration.
\\
\midrule

Approximation and optimization error: misspecified value classes, restricted
critics, or unstable saddle-point training can yield biased value predictions
\citep{uehara2021finite}.
&
Remap the biased but informative prediction to observed or Bellman-consistent values,
rather than requiring exact specification.
\\

\bottomrule
\end{tabularx}
\end{table}

\textbf{Contributions.}
\textbf{Contributions.}
\textbf{Contributions.}
We connect two largely separate literatures: post-hoc calibration for
supervised prediction and long-horizon value learning in RL. Existing work gives
rich methods for learning value functions, but much less theory for
model-agnostic reliability assessment and post-hoc correction of already learned
value predictors. Our contributions are:


\begin{enumerate}[leftmargin=1.5em]
\item We formalize \textbf{Bellman calibration}, a dynamic analogue of
supervised calibration for infinite-horizon value prediction. Bellman consistency
implies calibration, but calibration is weaker: predictions need only match
average Bellman backups among states assigned similar values.

\item We propose \textbf{\(\ell^2\)-Bellman calibration error} as a scalar
diagnostic for systematic value miscalibration and develop debiased off-policy
estimators of this error. We also prove a \textbf{calibration-refinement} bound
showing that value error decomposes into calibration error and the approximation
error of the best predictor based only on the original model output.

\item We introduce \textbf{post-hoc Bellman calibration} via histogram binning
and isotonic regression. The methods apply to any value predictor and use
doubly robust Bellman targets for off-policy data.

\item We prove \textbf{finite-sample guarantees} for histogram calibration,
controlling both \(\ell^2\)-Bellman calibration error and value error. The
bounds require neither Bellman completeness nor realizability and separate
statistical, auxiliary-model, and iteration errors.
\end{enumerate}


\subsection{Related work}

\textbf{Post-hoc calibration.}
Our work builds on post-hoc calibration for supervised prediction
\citep{zadrozny2001obtaining,zadrozny2002transforming,niculescu2005predicting,
bella2010calibration,guo2017calibration}; see \citet{roth2022uncertain} for an
overview. This literature calibrates static predictions against observed labels,
empirical frequencies, or conditional means. We study long-horizon value
predictions in MDPs, where returns are not directly observed and must be inferred
through Bellman structure. We are also related to causal calibration for static
treatment-effect prediction \citep{van2023causal,whitehouse2024orthogonal}, but
our setting is dynamic and long-horizon, with calibration targets defined by
value functions rather than one-step conditional effects. The closest work is \citep{van2025automaticDRL}, who use isotonic post-hoc
calibration of \(Q\)-functions inside Double RL policy-value estimators
\citep{kallus2020double,kallus2022efficiently}. Their calibration step is used
as an auxiliary-model correction for a global value estimator. We instead study calibrated value prediction directly and give finite-sample
guarantees for both Bellman calibration error and value error. The same framework also applies to \(Q\)-function calibration by replacing 
importance-weighted FVI with unweighted FQE and treating each state-action pair
as the state.
 

\textbf{Relation to Bellman learning and fine-tuning.}
Bellman calibration complements methods that change value learning itself,
including minimax or adversarial Bellman objectives
\citep{uehara2020minimax,uehara2023offline,xie2021batch,di2023pessimistic},
conservative updates \citep{kumar2020conservative,an2021uncertainty}, and
structural assumptions such as linearity, low rank, representations, or
discretization
\citep{shah2020sample,chang2022learning,peng1993convergence,van2006performance}.
It can be applied after any such method because it only needs scalar value
predictions and a reference sample. Unlike offline-to-online transfer,
critic/policy fine-tuning, or safe policy improvement
\citep{rusu2016progressive,nair2020awac,lee2022offline,laroche2019safe},
Bellman calibration fixes the target policy, value model, representation, and
Bellman learner, and fits only a one-dimensional post-hoc map to the target
value scale. Thus it is an output-scale correction rather than RL fine-tuning,
with finite-sample guarantees for calibration and value error. Related work on
calibrated dynamics uncertainty, conformal sets, and distributional RL targets
different objects
\citep{malik2019calibrated,zhang2023conformal,sun2023conformal,
bellemare2017distributional,dabney2018implicit}.

\section{Calibration for Value Functions}
\subsection{Setup and Notation}

We consider an MDP with state space \(\mathcal S\), discrete action space
\(\mathcal A\), initial-state distribution \(\rho\) for \(S_0\), transition kernel
\(P(s'\mid s,a)\), reward function \(r_0(s,a)\), and discount factor
\(\gamma\in[0,1)\). Data are collected under a behavior policy
\(b_0(a\mid s)\). Each calibration sample is an i.i.d. transition
\((S,A,R,S')\), where \((S,A)\sim \rho\times b_0\),
\(S'\sim P(\cdot\mid S,A)\), and
\(R=r_0(S,A)+\varepsilon\) with \(E[\varepsilon\mid S,A]=0\). We observe a calibration sample
\(\mathcal C_n:=\{(S_i,A_i,R_i,S_i')\}_{i=1}^n\).

Let \(\pi(a\mid s)\) be the target policy. The value function under \(\pi\) is $v_0(s)
=
E_{\pi}[\sum_{t=0}^\infty \gamma^t R_t
\,|\, S_0=s],$ where \(E_\pi\) denotes the law induced by
\(A_t\sim\pi(\cdot\mid S_t)\) and
\(S_{t+1}\sim P(\cdot\mid S_t,A_t)\).

For a state-action function \(f\) and state function \(v\), define
\[
(\pi f)(s):=\sum_{a\in\mathcal A}\pi(a\mid s)f(s,a),
\qquad
(b_0 f)(s):=\sum_{a\in\mathcal A}b_0(a\mid s)f(s,a),
\qquad
(Pv)(s,a):=E[v(S')\mid S=s,A=a].
\]
Write
\[
w_\pi:=\pi/b_0,\qquad
r_{0,\pi}:=\pi r_0,\qquad
P_\pi:=\pi P,\qquad
\mathcal T_\pi v := r_{0,\pi}+\gamma P_\pi v
\]
for the policy Bellman operator. The value function is the unique bounded fixed
point \citep{bellman1952theory,bellman1966dynamic}
\begin{equation}
\label{eqn::bellman}
v_0=\mathcal T_\pi v_0 .
\end{equation}
For a measure \(\mu\) on \(\mathcal S\), write
\[
\|f\|_{2,\mu}
:=
\left\{\int f(s)^2\,\mu(ds)\right\}^{1/2}.
\]
Let \(L^2(\rho)\) denote the space of square-integrable state functions under
\(\rho\). Let \(\|f\|_{n,S,A}\) and \(\|g\|_{n,S'}\) denote empirical \(L^2\)
norms over \(\{(S_i,A_i)\}_{i=1}^n\) and \(\{S_i'\}_{i=1}^n\), respectively.
We write \(\lesssim\) for inequalities up to absolute constants and
\([B]:=\{1,\ldots,B\}\).
 

\subsection{Bellman Calibration: A Minimal Requirement}

Let \(\hat v\) be an estimated value function for the target policy \(\pi\),
fit on data independent of the calibration sample \(\mathcal C_n\). Our goal is
not to replace the underlying value-learning procedure, but to impose a weak
post-hoc Bellman consistency requirement on its output. Throughout, we call the
scalar \(\hat v(S)\) the predicted value; in the supervised calibration
literature, this same scalar is often called a prediction score.

A natural starting point is \emph{statewise Bellman consistency}:
\[
\hat v(S)
\approx
E_\pi\!\left[R+\gamma \hat v(S')\mid S\right].
\]
Equivalently, the Bellman residual
\(R+\gamma \hat v(S')-\hat v(S)\) should have conditional mean zero given the
full state. In high-dimensional or continuous state spaces, this is a strong
requirement: checking or enforcing it requires estimating a state-dependent
conditional mean and is closely tied to structural assumptions such as Bellman
completeness \citep{chen2019information}.

We instead require Bellman consistency only after grouping states by their
predicted value:
\[
\hat v(S)
\approx
E_\pi\!\left[R+\gamma \hat v(S')\mid \hat v(S)\right].
\]
Thus, predictions should agree on average with their Bellman targets within
groups of states assigned the same predicted value \(\hat v(S)\). This is the
dynamic analogue of post-hoc calibration in supervised prediction, where
predictions are required to be correct on average conditional on their own
numerical output
\citep{zadrozny2002transforming,niculescu2005predicting,guo2017calibration}.
The benefit is that a high-dimensional Bellman moment is reduced to a
one-dimensional conditional moment, enabling post-hoc correction without
requiring the full fitted Bellman update to be correctly specified.

\textbf{Bellman calibration.}
For any candidate value function \(v\), define the Bellman-calibration map
\begin{equation}
\label{eqn::calibrationmap}
\Gamma_0(v)(s)
:=
E_\pi\!\left[R+\gamma v(S') \mid v(S)=v(s)\right].
\end{equation}
We say that \(v\) is \emph{Bellman calibrated} if
\[
v(S)=\Gamma_0(v)(S)\quad\text{almost surely},
\]
or equivalently,
\[
E_\pi\!\left[R+\gamma v(S')-v(S)\mid v(S)\right]=0
\quad\text{almost surely}.
\]
 

\textbf{A necessary but weak condition.}
The true value function \(v_0\) is Bellman calibrated. The Bellman equation gives
\(v_0(S)=E_\pi[R+\gamma v_0(S')\mid S]\); conditioning on the coarser predicted value
\(v_0(S)\) yields
\[
v_0(S)
=
E_\pi[R+\gamma v_0(S')\mid v_0(S)]
=
\Gamma_0(v_0)(S).
\]
Thus Bellman consistency implies Bellman calibration. The converse need not
hold: because calibration only conditions on the scalar predicted value \(v(S)\),
state-level errors can average out within predicted-value groups
\citep{gupta2020distribution,van2023causal}. Hence Bellman calibration is a
minimal reliability requirement, not a guarantee of value accuracy.

\textbf{Reward-model interpretation.}
Define the reward implied by a value predictor \(v\) as
\[
r_{\pi,v}(s):=v(s)-\gamma P_\pi v(s).
\]
This is the one-step reward under which \(v\) would satisfy the Bellman
equation for policy \(\pi\). Bellman calibration is equivalent to
\[
E_\pi\!\left[R-r_{\pi,v}(S)\mid v(S)\right]=0.
\]
Thus, among states assigned the same predicted long-run value, the reward
implied by \(v\) matches the true one-step reward on average. When
\(\gamma=0\), this reduces to ordinary regression calibration
\citep{roth2022uncertain}.

\subsection{A Calibration--Refinement Bound for Value Error}

We next relate Bellman calibration to value prediction error. The key point is
that calibration is a post-processing operation: it can correct systematic
Bellman residuals among states with the same predicted value, but it cannot
recover information lost by reducing the state \(S\) to \(\hat v(S)\).

For a candidate value function \(v\), define the \(\ell^2\) \textbf{Bellman calibration
error}
\begin{equation}
\label{eqn::calerror}
\mathrm{Cal}_{\ell^2}^2(v)
:=
\|v-\Gamma_0(v)\|_{2,\rho}^2
=
E_\rho\!\left[
  \left\{
    E_\pi[R+\gamma v(S')-v(S)\mid v(S)]
  \right\}^2
\right].
\end{equation}
This criterion is zero exactly when \(v\) is Bellman calibrated, and measures
the remaining Bellman residuals conditional on the predicted value.

The value predictors generated by the original prediction \(\hat v(S)\) form the
class
\[
\mathcal V_{\hat v}
:=
\{g\circ\hat v:g:\mathbb R\to\mathbb R \text{ measurable}\},
\qquad
\Pi_{\hat v}
:=
\text{\(L^2(\rho)\)-projection onto }\mathcal V_{\hat v}.
\]
Define the projected transition by
\[
P_{\pi,\hat v}
:=
\Pi_{\hat v}P_\pi,
\qquad
P_{\pi,\hat v}f(s)
=
E_\pi[f(S')\mid \hat v(S)=\hat v(s)].
\]

Define the discounted projected concentrability by
\[
\mathcal C_{\hat v,\gamma}
:=
\left\{(1-\gamma)\sum_{t=0}^\infty
\gamma^t
\left\|
\frac{d(\rho P_{\pi,\hat v}^t)}{d\rho}
\right\|_\infty^{1/2}
\right\}^2 .
\]

\begin{enumerate}[label=\textbf{A\arabic*}, ref={A\arabic*}, leftmargin=1.5em, series=cond]
\item \label{cond::A1}
\textbf{Projected concentrability.}
\(\mathcal C_{\hat v,\gamma}<\infty\).
\end{enumerate}
When \(\rho\) is stationary under \(P_{\pi}\), it is also stationary under
\(P_{\pi,\hat v}\) by Lemma~\ref{lemma:aggregated-stationary}, so
\(\mathcal C_{\hat v,\gamma}=1\). More generally,
\(\mathcal C_{\hat v,\gamma}\) measures distribution drift from \(\rho\): it is
small when projected rollouts remain comparable to \(\rho\), and large when they
concentrate where \(\rho\) has little mass.

\begin{theorem}[Calibration--refinement bound]
\label{theorem::calrefine}
Under \ref{cond::A1},
\[
\|\hat v-v_0\|_{2,\rho}
\le
\frac{\sqrt{\mathcal C_{\hat v,\gamma}}}{1-\gamma}
\left\{
\mathrm{Cal}_{\ell^2}(\hat v)
+
\min_{\theta:\mathbb R\to\mathbb R}
\|\theta\circ\hat v-v_0\|_{2,\rho}
\right\}.
\]
\end{theorem}

Theorem~\ref{theorem::calrefine} gives a post-processing guarantee: value error
is controlled by Bellman calibration error \(\mathrm{Cal}_{\ell^2}(\hat v)\)
plus the refinement error $\min_{\theta:\mathbb R\to\mathbb R}
\|\theta\circ\hat v-v_0\|_{2,\rho},$ the error of the best value predictor that depends on the state only through
\(\hat v(S)\). The calibration error measures systematic Bellman residuals
within levels of \(\hat v(S)\), whereas the refinement error measures the
remaining value variation not captured by \(\hat v(S)\). Thus, calibration can
remove Bellman miscalibration using only a post-processing map of the original
prediction, but it cannot recover information about \(v_0(S)\) absent from
\(\hat v(S)\). In particular, when calibration error is driven to zero, value
error is controlled, up to the projected concentrability and discount factors,
by the refinement error. This mirrors classical calibration--refinement
decompositions in classification and regression
\citep{murphy1973new,degroot1983comparison,van2023causal,whitehouse2024orthogonal}.
The factor \(\sqrt{\mathcal C_{\hat v,\gamma}}/(1-\gamma)\) plays the same role
as the concentrability and discount penalties in classical FVI/FQE error bounds
\citep{munos2008finite}.

\subsection{Evaluating Bellman Calibration Error with Doubly Robust Bellman Targets}
\label{sec:bellman-calibration-error}

We now focus on estimating \(\mathrm{Cal}_{\ell^2}(\hat v)\), the component of
Theorem~\ref{theorem::calrefine} that can be evaluated from transition data
and improved without refitting the value predictor. This makes Bellman calibration error both
a diagnostic for systematic Bellman residuals within the level sets of a fixed
predictor \(\hat v\) and an objective for post-hoc calibration.

Because data are generated under the behavior policy \(b_0\), we estimate the
statewise Bellman target using a doubly robust Bellman pseudo-outcome, extending
the calibration target estimates of \citep{xu2022calibration,van2023causal} to
dynamic settings. Let \(\widehat w_\pi\) estimate \(w_\pi:=\pi/b_0\). For a
fixed \(v\), define
\[
q_v(s,a)
:=
E\{R+\gamma v(S')\mid S=s,A=a\}.
\]
Let \(\hat q_v\) estimate \(q_v\). For \(O=(S,A,R,S')\), set
\[
\widehat Y_v(O)
:=
(\pi \hat q_v)(S)
+
\widehat w_\pi(A\mid S)
\{R+\gamma v(S')-\hat q_v(S,A)\},
\qquad
\widehat{\mathcal T}_\pi v(s)
:=
E\{\widehat Y_v(O)\mid S=s\}.
\]
The next result gives the conditional bias of this pseudo-outcome.

\begin{theorem}[Doubly robust Bellman target]
\label{theorem::DRpseudo}
For any fixed \(v\),
\[
E\!\left[
\{\widehat{\mathcal T}_\pi v-\mathcal T_\pi v\}(S)
\mid v(S)=t
\right]
=
E\!\left[
  b_0\!\left\{(w_\pi-\widehat w_\pi)(\hat q_v-q_v)\right\}(S)
  \mid v(S)=t
\right].
\]
Consequently, \(E\{\widehat Y_v(O)\mid v(S)\}
=
E\{(\mathcal T_\pi v)(S)\mid v(S)\}\) if either
\(\widehat w_\pi=w_\pi\) or \(\hat q_v=q_v\).
\end{theorem}

This construction includes importance weighting \((\hat q_v=0)\) and
model-based Bellman-target evaluation \((\widehat w_\pi=0)\) as special cases.
The theorem shows that either nuisance estimator suffices to recover the
prediction-conditional Bellman target.

To estimate the proposed calibration error, we use a \textbf{cross-fitted debiased moment estimator}.
Let \(U_i:=v(S_i)\) and \(\widehat Y_i:=\widehat Y_v(O_i)\). Following
\citet{xu2022calibration}, we estimate the squared calibration error by 
\[
\widehat{\mathrm{Cal}}_{\ell^2,\mathrm{db}}^2(v)
:=
\frac1n\sum_{i=1}^n
\{\widehat Y_i-U_i\}
\{\widehat m_{-i}(U_i)-U_i\},
\]
where \(\widehat m_{-i}\) is fit without observation \(i\) to estimate
\(u\mapsto E(\widehat Y_v\mid v(S)=u)\). The population identity behind this
estimator is
\[
E\!\left[
\{\widehat Y_v-v(S)\}
\left\{E(\widehat Y_v\mid v(S))-v(S)\right\}
\right]
=
E\!\left[
\left\{E(\widehat Y_v\mid v(S))-v(S)\right\}^2
\right].
\]
By Theorem \ref{theorem::DRpseudo}, this equals
\(\mathrm{Cal}_{\ell^2}^2(v)\), up to function approximation error.


\section{Post-hoc Bellman Calibration}
\label{sec:posthoc-calibration}

\subsection{Iterated Bellman Calibration}
\label{sec:iterated-bellman-calibration}

Algorithm~\ref{alg::cal-gen} gives a post-hoc FVI-style procedure that
transforms an initial value predictor \(\hat v\) into a calibrated predictor
\(V_{\mathrm{cal}}=V_K=g_K\circ \hat v\). Starting from \(V_0=\hat v\), each
iteration forms doubly robust Bellman targets \(\widehat Y_{V_k}(O_i)\) and
regresses them on the fixed one-dimensional inputs \(\hat v(S_i)\) over a
calibrator class \(\mathcal G\), yielding
\(V_{k+1}=g_{k+1}\circ \hat v\). Equivalently, for
\(\mathcal V_{\hat v}(\mathcal G):=\{g\circ\hat v:g\in\mathcal G\}\), the
algorithm is a projected Bellman update restricted to
\(\mathcal V_{\hat v}(\mathcal G)\). This avoids high-dimensional function
approximation while correcting Bellman miscalibration in the original predictor;
the contraction results below give geometric convergence, so
\(O\{\log(1/\varepsilon)/(1-\gamma)\}\) iterations suffice to reach optimization
error \(\varepsilon\).

\begin{algorithm}[h]
\caption{\textbf{Iterated Bellman Calibration}}
\label{alg::cal-gen}
\begin{algorithmic}[1]
\BCINPUT \(\hat v\), \(\mathcal C_n=\{O_i=(S_i,A_i,R_i,S_i')\}_{i=1}^n\), Pseudo-outcome \(v\mapsto\widehat Y_v\), class \(\mathcal G\), iterations \(K\)
\STATE Initialize \(V_0=\hat v\)
\FOR{\(k=0,\ldots,K-1\)}
  \STATE \(V_{k+1}=g_{k+1}\circ\hat v \quad \text{where} \quad g_{k+1}\in\arg\min_{g\in\mathcal G}\sum_{i=1}^n
  \{\widehat Y_{V_k}(O_i)-g(\hat v(S_i))\}^2\)
\ENDFOR
\BCOUTPUT \(V_K\)
\end{algorithmic}
\end{algorithm}

Heuristically, \(V_K\) becomes calibrated because each iteration approximates
one application of the projected Bellman operator $\mathcal B_{\hat v} f
:=
\Pi_{\hat v} r_{0,\pi} + \gamma P_{\pi,\hat v} f,$ whose fixed point over functions of \(\hat v(S)\) has zero Bellman calibration
error. Exact finite-sample recovery of this projected fixed point would
generally require realizability in \(\mathcal V_{\hat v}(\mathcal G)\) and
Bellman completeness of the fitted update \citep{munos2008finite}. We instead
seek guarantees for post-hoc improvement through Bellman calibration without
imposing these conditions. We therefore instantiate Algorithm~\ref{alg::cal-gen} with one-dimensional
calibrators based on histogram regression \citep{stone1977consistent} and
isotonic regression \citep{barlow1972isotonic}. These are dynamic analogues of
histogram binning
\citep{zadrozny2001obtaining,zadrozny2002transforming,gupta2021distribution}
and isotonic calibration
\citep{zadrozny2002transforming,niculescu2005predicting,van2023causal,
van2025generalized}, and they admit finite-sample Bellman calibration
guarantees without Bellman completeness.

 

\subsection{Histogram Calibration}
\label{sec:hist}

Histogram calibration instantiates Algorithm~\ref{alg::cal-gen} with a
piecewise-constant calibrator. Partition the initial values
\(\{\hat v(S_i)\}_{i=1}^n\) into ordered bins \(I_1,\ldots,I_B\), using either
empirical quantiles (\emph{equal-mass}) or uniform discretization
(\emph{equal-width}). We take
\[
\mathcal G_B
:=
\mathrm{span}\{\mathbf 1_{I_b}:b\in[B]\},
\qquad
\mathcal H_B(\hat v)
:=
\{g\circ\hat v:g\in\mathcal G_B\},
\]
so every \(h\in\mathcal H_B(\hat v)\) is a value predictor that is constant
within bins of the original prediction. The bins may be data-adaptive, but
\(B\) is deterministic and may grow with \(n\).

With this class, each iteration simply averages the fitted Bellman targets
within each bin:
\[
g_{k+1}(t)
=
\frac{1}{n_b}
\sum_{i:\hat v(S_i)\in I_b}
\widehat Y_i^{(k)},
\qquad
t\in I_b,
\]
where \(n_b:=\sum_{i=1}^n\mathbf 1\{\hat v(S_i)\in I_b\}\). Thus, histogram
calibration is fitted Bellman iteration on a finite aggregation of the initial
predicted value. At convergence, the output approximately satisfies the empirical aggregated
Bellman equation within each bin:
\[
V_K(s)
\approx
E_n\!\left[
\widehat Y_{V_K}(O)
\,\middle|\,
\hat v(S)\in I_b
\right],
\qquad \hat v(s)\in I_b .
\]
This is the empirical analogue of Bellman calibration and is the key property
used in the next result.

 \textbf{Calibration error.}
We next give finite-sample, finite-iteration bounds for the Bellman calibration
error \eqref{eqn::calerror}. The analysis uses the following regularity
conditions.

\begin{enumerate}[label=\textbf{C\arabic*}, ref={C\arabic*}, leftmargin=1.5em, series=cond]

\item \label{cond::C1}
\textbf{Boundedness.}
\(R\), \(w_\pi\), \(\widehat w_\pi\), \(\hat v\), the Bellman targets \(q_f\),
and their estimates \(\hat q_f\) are uniformly bounded by \(M<\infty\) for all
\(f\) considered.

\item \label{cond::C1b}
\textbf{Sample splitting.}
The auxiliary estimators \(\widehat w_\pi\) and \(v\mapsto\hat q_v\), and the
initial predictor \(\hat v\) are trained on data independent of
\(\mathcal C_n\).

\item \label{cond::C2}
\textbf{Stability of the Bellman-target regression.}
For all \(f,g\in\mathcal H_B(\hat v)\) with
\(\|f\|_\infty,\|g\|_\infty\le M\),
\[
\|\hat q_f-\hat q_g\|_{n,S,A}
\le
L\,\|f-g\|_{n,S'} .
\]
\end{enumerate}

These conditions are mainly technical. Boundedness can be enforced by clipping
or truncation, and sample splitting can be replaced by cross-fitting. Condition
\ref{cond::C2} controls the complexity introduced by the map
\(v\mapsto\hat q_v\); it holds trivially for the importance-weighted target
\((\hat q_v=0)\), and holds with \(L=1\) for cellwise estimates of
\(E\{R+\gamma f(S')\mid S,A\}\).

\begin{theorem}[Calibration error for histogram binning]
\label{theorem::Calerrorhist}
Assume \ref{cond::C1}--\ref{cond::C2}. Then there is a constant
\(C<\infty\) such that, for any \(K\in\mathbb N\), with probability at least
\(1-\delta\),
\begin{align*}
\mathrm{Cal}_{\ell^2}(V_K)
\le\;&
C\left\{
\sqrt{\frac{B}{n}\log\!\left(\frac{n}{B}\right)}
+
\sqrt{\frac{\log(1/\delta)}{n}}
\right\} \\
&+
\big\|b_0\{(\widehat w_\pi-w_\pi)
(\widehat q_{V_K}-q_{V_K})\}\big\|_{2,\rho}  +
\big\|
\widehat{\mathcal T}_\pi V_K
-
\widehat{\mathcal T}_\pi V_{K-1}
\big\|_{2,\rho}.
\end{align*}
\end{theorem}

The bound separates three errors. The first line is the statistical error of
one-dimensional histogram calibration, matching classical rates when the
calibration target is observed directly
\citep{gupta2020distribution,gupta2021distribution,van2025generalized}. The
second line is the doubly robust auxiliary-model error from
Theorem~\ref{theorem::DRpseudo}. The third line is the finite-iteration error,
which decays geometrically as $\gamma^K$ under the contraction conditions in
Appendix~\ref{appendix::convergence}, as in fitted value iteration
\citep{munos2008finite}.

The leading calibration rate does not accumulate over the horizon: \(\gamma\)
enters only through the finite-iteration and auxiliary-model terms. This reflects that
Bellman calibration is a one-step conditional prediction problem, even though
the value function itself is long-horizon. However, Theorem~\ref{theorem::calrefine}
shows that translating calibration into value accuracy still depends on
discounting and coverage.

\textbf{Value error.}
We next show that post-hoc calibration does not add error beyond the
information loss from restricting to the histogram calibrator class, while it can
remove Bellman miscalibration in the original predictor.

Let \(\Pi_B\) denote the \(L^2(\rho)\)-projection onto
\(\mathcal H_B(\hat v)\), i.e.,
\[
\Pi_B q:=\arg\min_{h\in\mathcal H_B(\hat v)}\|q-h\|_{2,\rho},
\qquad
\mathcal B_{\pi,B}:=\Pi_B\mathcal T_\pi,
\qquad
V_B^\star=\mathcal B_{\pi,B}V_B^\star .
\]
Write \(P_{\pi,B}:=\Pi_B P_\pi\), so that, for
\(\hat b(s)=b\) when \(\hat v(s)\in I_b\),
\[
P_{\pi,B}f(s)
=
E_\pi[f(S')\mid \hat v(S)\in I_{\hat b(s)}].
\]
Finally, define the aggregated discounted concentrability by
\[
\mathcal C_{B,\gamma}
:=
\left\{(1-\gamma)\sum_{t=0}^\infty
\gamma^t
\left\|
\frac{d(\rho P_{\pi,B}^t)}{d\rho}
\right\|_\infty^{1/2}
\right\}^2 .
\]

\begin{enumerate}[label=\textbf{C\arabic*}, ref={C\arabic*}, leftmargin=1.5em, resume=cond]
\item \label{cond::Ccoverage}
\textbf{Aggregated projected concentrability.}
\(\mathcal C_{B,\gamma} < \infty\).
\end{enumerate}
 

\begin{theorem}[Value error for histogram calibration]
\label{theorem::regretHist}
Assume \ref{cond::C1}--\ref{cond::Ccoverage}. Then there exists a constant
\(C<\infty\) such that, for any \(K\in\mathbb N\), with probability at least
\(1-\delta\),
{
\begin{align*}
\|V_K-v_0\|_{2,\rho}
\le\;&
\gamma^K\|\hat v-V_B^\star\|_{\infty}
+
\frac{\sqrt{\mathcal C_{B,\gamma}}}{1-\gamma}
\min_{h\in\mathcal H_B(\hat v)}\|h-v_0\|_{2,\rho}
\\
&+
C\frac{\sqrt{\mathcal C_{B,\gamma}}}{1-\gamma}
\Bigg[
  \sqrt{\frac{B}{n}\log\!\left(\frac{n}{B}\right)}
  +
  \sqrt{\frac{\log(K/\delta)}{n}}
 +
\max_{0\le j\le K}
\big\|
b_0\{
(\widehat w_\pi-w_\pi)
(\widehat q_{V_j}-q_{V_j})
\}
\big\|_{2,\rho}
\Bigg].
\end{align*}
}
\end{theorem}
\textbf{Proof idea.}
The result avoids realizability and Bellman completeness for the original value
class because completeness holds after aggregation. For any
\(f\in\mathcal H_B(\hat v)\), the aggregated transition \(P_{\pi,B}f\) is
binwise constant in \(\hat v\), so \(P_{\pi,B}f\in\mathcal H_B(\hat v)\)
and hence \(\mathcal B_{\pi,B}f\in\mathcal H_B(\hat v)\). The proof uses
sup-norm contraction to obtain the aggregated fixed point, then propagates
\(\rho\)-norm one-step regression errors through a discounted resolvent
controlled by \(\mathcal C_{B,\gamma}\). It also bounds
\(V_B^\star-v_0\) by a histogram analogue of
Theorem~\ref{theorem::calrefine}. Unlike prior discretization arguments for
Bellman consistency
\citep{whitt1978approximations,peng1993convergence,van2006performance,
chen2019information,xie2021batch}, we discretize \(\hat v(S)\) to
target Bellman calibration.
 
\textbf{Discussion.}
Theorems~\ref{theorem::Calerrorhist} and~\ref{theorem::regretHist} separate
calibration from value accuracy. With known importance weights and negligible
iteration error, histogram calibration drives Bellman calibration error to zero
at the one-dimensional rate \(\sqrt{(B/n)\log(n/B)}\). Value error contains the
same stochastic term, but also includes the prediction-only approximation error $(\sqrt{\mathcal C_{B,\gamma}})/(1-\gamma)
\min_{h\in\mathcal H_B(\hat v)}\|h-v_0\|_{2,\rho},$ the cost of restricting the calibrated predictor to binned transformations of
the original prediction. Thus \(B\) controls the usual bias--variance tradeoff:
coarse bins may discard information in \(\hat v\), while fine bins increase
variance. The coverage factor \(\mathcal C_{B,\gamma}\) is the histogram
analogue of FVI concentrability.

It suffices to choose \(B\) by predictive cross-validation. The prediction-only approximation term
satisfies
\[
\min_{h\in\mathcal H_B(\hat v)}\|h-v_0\|_{2,\rho}
\le
\|\hat v-v_0\|_{2,\rho}
+
\min_{h\in\mathcal H_B(\hat v)}\|h-\hat v\|_{2,\rho}.
\]
Thus binning is no worse than the original predictor, up to identity
approximation on the range of \(\hat v(S)\). For equal-mass bins, the histogram
bias--variance balance suggests \(B\asymp n^{1/3}\), yielding
\citep{gyorfi2002distribution}
\[
\|V_K-v_0\|_{2,\rho}
\le
\frac{\sqrt{\mathcal C_{B,\gamma}}}{1-\gamma}
\|\hat v-v_0\|_{2,\rho}
+
O_p\{(\log n/n)^{1/3}\},
\]
with Bellman calibration error vanishing at the same rate.


 


\subsection{Isotonic Calibration}

Histogram calibration requires choosing the number of bins \(B\). As a
low-tuning alternative, we consider \emph{isotonic iterated Bellman
calibration}, obtained by taking
\(\mathcal G=\mathcal G_{\mathrm{iso}}\) in
Algorithm~\ref{alg::cal-gen}, where \(\mathcal G_{\mathrm{iso}}\) is the class
of nondecreasing functions on \(\mathbb R\). Each update isotonic-regresses the
fitted Bellman targets on the initial values of \(\hat v(S_i)\), yielding a
monotone calibrator. This can be computed in near-linear time by the
pool-adjacent-violators algorithm (PAVA) \citep{best1990active}.

Like histogram calibration, isotonic calibration averages Bellman targets over
groups of nearby predicted values, but the partition is chosen adaptively by the
monotonicity constraint \citep{barlow1972isotonic}. This provides
regularization without selecting \(B\). Since the identity map is monotone,
isotonic calibration can preserve an already reliable predictor while correcting
systematic monotone distortions.

\begin{enumerate}[label=\textbf{C\arabic*}, ref={C\arabic*}, leftmargin=1.5em, resume=cond]
\item \label{cond::C3}
\textbf{Finite variation.}
There exists \(V<\infty\) such that, almost surely, the function $t
\mapsto
E\!\left[
(\mathcal T_\pi V_K)(S)
\mid
\hat v(S)=t,\mathcal C_n
\right]$
has total variation at most \(V\).
\end{enumerate}

Condition~\ref{cond::C3} is standard in isotonic calibration
\citep{van2023causal} and controls the complexity of the one-dimensional
calibration curve. For histogram binning, the analogous complexity is governed
by the number of bins \(B\).

\begin{theorem}[Calibration error for isotonic calibration]
\label{theorem::Calerroriso}
Assume \ref{cond::C1}, \ref{cond::C1b}, \ref{cond::C2}, and \ref{cond::C3}
with \(\mathcal G=\mathcal G_{\mathrm{iso}}\). Then there exists a constant
\(C<\infty\) such that, for any \(K\in\mathbb N\), with probability at least
\(1-\delta\),
{\small
\begin{align*}
\mathrm{Cal}_{\ell^2}(V_K)
\le\;&
C\left\{
n^{-1/3}
+
\sqrt{\frac{\log(1/\delta)}{n}}
\right\}  +
\big\|
b_0\{
(\widehat w_\pi-w_\pi)
(\widehat q_{V_K}-q_{V_K})
\}
\big\|_{2,\rho}  +
\big\|
\widehat{\mathcal T}_\pi V_K
-
\widehat{\mathcal T}_\pi V_{K-1}
\big\|_{2,\rho}.
\end{align*}
}
\end{theorem}
 
The isotonic analysis parallels the histogram case, but uses the adaptive
pooling induced by the pool-adjacent-violators algorithm. We do not assume the
calibration target is monotone; monotonicity is used only as a regularizer for
pooling nearby predicted values. The resulting \(n^{-1/3}\) rate matches the
classical isotonic regression rate \citep{chatterjee2013improved} and the rate
of histogram binning with \(B\asymp n^{1/3}\), but avoids choosing \(B\) in
advance.

 
Our histogram proof does not yield value-error bounds for fully iterated
isotonic calibration without Bellman completeness, because the PAVA partition can
change across Bellman iterations. Appendix~\ref{appendix::hybrid} gives a hybrid
procedure that first learns the partition by isotonic regression and then
applies the histogram analysis conditionally on that fixed partition.

\section{Experimental investigation}
\label{sec:experiments}

We evaluate Bellman calibration as a post-hoc correction for fitted off-policy
value predictors. Each base learner estimates a \(Q\)-function and forms the
target-policy value predictor \(\widehat V\); calibration then fits a
one-dimensional map \(g\), using importance-weighted fitted value iteration, and
reports \(g\circ \widehat V\). We use linear FQE and random-feature FQE as base
learners, with linear, histogram, and isotonic calibrators. The experiments
isolate the failure mode in Table~\ref{tab:rl-calibration-motivation}:
\(\widehat V\) is informative but mis-scaled. For non-temporal-shift settings,
we use 5-fold cross-calibration \citep{van2023causal}; for temporal shift, we
train FQE on a large old-regime batch, fit \(g\) on a small current-regime
batch, and evaluate on current-regime states. All main results use 100
independent replications and independent evaluation data. Values below one
indicate improvement. See Appendix~\ref{app:experiments} for details.

\begin{figure}[t]
  \centering
    \centering
    \includegraphics[width=\linewidth]{figures/calibration_story_compact.pdf}
    \vspace{-0.5em}
    \subcaption{Synthetic failure modes.}
    \label{fig:calibration-story}
\caption{\textbf{Post-hoc value calibration improves mis-scaled value
predictions.}
Value MSE and Bellman calibration error ratios across synthetic failure modes,
relative to the matched uncalibrated baseline.}
  \label{fig:main-calibration-results}
\end{figure}

Figure~\ref{fig:calibration-story} and
Table~\ref{tab:main-simulation-results} in Appendix~\ref{app:tablemain} show
consistent gains across four failure modes. Under affine misspecification,
linear calibration reduces true-\(V\) MSE by \(25\%\) and Bellman calibration
error by \(37\%\). With finite Bellman-iteration error, random-feature FQE
predictions remain rank-informative but under-scaled; linear and isotonic
calibration reduce true-\(V\) MSE by \(22\%\) and Bellman calibration error by
\(26\)--\(34\%\). Under temporal reward shift, recalibrating an old-regime FQE
model on a small current-regime batch reduces current-regime MSE by \(37\%\) and
Bellman calibration error by \(28\%\), outperforming a same-size retrain. Richer
calibrators help under monotone misspecification: isotonic and histogram
calibration nearly halve Bellman calibration error, while linear calibration
has only mild effect. Overall, linear calibration corrects affine scale errors,
whereas isotonic calibration handles more general monotone distortions with
no tuning.
 
\textbf{Conclusion.}
Bellman calibration provides an estimable, post-hoc reliability criterion for
long-horizon value prediction: predictions should agree, on average, with their
target-policy Bellman backups among states assigned similar values. Iterated
Bellman Calibration turns this testable criterion into a one-dimensional
recalibration procedure that can control calibration error without Bellman
completeness or value-function realizability. Appendix~\ref{app:concl}
discusses the main takeaways, practical limitations from coverage and nuisance
estimation, and future directions including stationary weighting, smoother
calibrators, and larger offline RL evaluations.

 

\bibliographystyle{plainnat}
\bibliography{../../common/ref}


 

 




 
\appendix

\section{Table for experiment}
\label{app:tablemain}
\begin{table*}[h]
\centering
\caption{\textbf{Value-scale correction across offline RL failure modes.}
Ratios are relative to the matched uncalibrated baseline. Win rate is the
seed-matched true-\(V\) MSE win rate.}
\label{tab:main-simulation-results}
\scriptsize
\setlength{\tabcolsep}{4pt}
\begin{tabular}{@{}p{0.45\linewidth}lrrr@{}}
\toprule
Regime and learner & Calibrator & Rel. \(V\) MSE & Rel. cal. err. & Win rate \\
\midrule
Affine misspec., linear FQE & Linear & 0.750 & 0.631 & 0.94 \\
Affine misspec., linear FQE & Isotonic & 0.734 & 0.764 & 0.94 \\
Finite iteration, RF FQE \(K=4\) & Linear & 0.780 & 0.703 & 0.94 \\
Finite iteration, RF FQE \(K=4\) & Isotonic & 0.776 & 0.736 & 0.94 \\
Temporal reward shift, RF FQE & Linear recent & 0.630 & 0.717 & 0.97 \\
Temporal reward shift, RF FQE & Isotonic recent & 0.627 & 0.916 & 0.98 \\
Monotone misspec., linear FQE & Linear & 0.972 & 0.986 & 0.64 \\
Monotone misspec., linear FQE & Isotonic & 0.916 & 0.537 & 0.79 \\
Monotone misspec., linear FQE & Histogram & 0.917 & 0.544 & 0.77 \\
\bottomrule
\end{tabular}
\end{table*}


\section{Conclusion and Future Work}
\label{app:concl}
We introduced Bellman calibration as a post-hoc reliability criterion for
long-horizon value prediction. The criterion asks that predictions agree, on
average, with their target-policy Bellman backups within groups of states
assigned similar predicted values. This is weaker than statewise Bellman
consistency, but strong enough to diagnose systematic value-scale errors and to
support model-agnostic recalibration of an already learned value predictor.

Iterated Bellman Calibration turns this criterion into a practical correction:
fit a one-dimensional calibration map using off-policy Bellman targets, then
iterate the induced calibrated Bellman update. Because the recalibration is
one-dimensional, histogram and isotonic variants can achieve finite-sample
control of Bellman calibration error without Bellman completeness or
value-function realizability. The resulting value-error bounds separate
statistical calibration error, finite-iteration error, nuisance-estimation
error, and the approximation error of the best predictor expressible as a
function of the original value estimate. Thus calibration is most useful when
the base predictor is informative but mis-scaled, and its limits are explicit
when the original prediction discards information needed for accurate value
prediction.

The main practical challenges are coverage, nuisance estimation, and
extrapolation. Off-policy Bellman calibration requires enough overlap to estimate
target-policy Bellman backups on the calibration sample, and poor coverage can
make importance-weighted or doubly robust targets unstable. The method also
corrects the value scale through the original predictor, so one-dimensional
calibration cannot recover state information absent from that predictor. Future
work includes combining Bellman calibration with stationary-density-ratio or
emphatic-style weighting \citep{vdLaanKallus2025FQE,mahmood2015emphatic,patterson2022generalized}, developing smoother calibrators for extrapolation
under distribution shift \citep{platt1999probabilistic,jiang2011smooth,guo2017calibration}, and studying calibration as both a diagnostic and
post-hoc correction in larger offline RL benchmarks.
 



\section{Additional Experimental Details}
\label{app:experiments}

For the non-temporal-shift experiments, we use 5-fold cross-calibration
\citep{van2023causal}. In each fold, the value learner and calibrator are
trained out of fold, and calibrated predictions are evaluated on the held-out
fold. Fold-specific calibrated predictors are aggregated by pointwise medians.
The uncalibrated comparator uses the same base learner refit on the full
training sample. Calibration is fit by importance-weighted fitted value
iteration, using the target-to-behavior action-probability ratio as weights. In
the temporal-shift experiment, we train FQE on a large old reward-regime batch,
fit the calibration map on a small held-out current-regime batch, and evaluate
only on current-regime states. All main results use 100 independent replications
and independent evaluation data.



\paragraph{Data-generating processes.}
All four simulation panels use discounted MDPs with continuous states
\(S_t\in\mathbb{R}^d\), actions \(A_t\in\{1,2,3\}\), and softmax behavior and
target policies. The target policy has logits \(s^\top w_a^\pi/\tau_\pi\). The
behavior policy uses shifted logits \(s^\top w_a^b/\tau_b\), where
\(w_a^b=w_a^\pi-\Delta u_a\). The main experiments use the moderate-shift
setting. Initial states are Gaussian, \(S_0\sim N(\mu_0,I_d)\). Except where
noted below, transitions and rewards are
\[
  S_{t+1}
  =
  \tanh\{M S_t+\alpha_{A_t}+0.22\sin(S_t)+0.08\cos(S_t^{\mathrm{rev}})\}
  + \sigma_\varepsilon \varepsilon_t ,
\]
and
\[
  r_0(s,a)
  =
  s^\top\beta
  +0.65\sin(s_1)
  +0.35\cos(s_2)
  +s^\top\eta_a
  +\rho_a
  -0.06\|s\|_2^2 .
\]
Observed rewards add mean-zero Gaussian noise. The ground-truth value \(V^\pi(s)\) is
computed by independent target-policy Monte Carlo rollouts using the conditional
mean reward. Training transitions, calibration folds, Bellman-evaluation
transitions, initial evaluation states, and Monte Carlo rollouts use independent
seeds.

The four panels isolate different value-scale errors. In the affine
misspecification panel, the same nonlinear MDP is used, but the reward is
transformed to
\[
  r_{\mathrm{aff}}(s,a)=1.25+0.7\,r_0(s,a).
\]
The learner is linear FQE with a deliberately restricted feature class. This
produces informative value predictions with an affine slope/intercept error, so
linear calibration should be sufficient.

In the finite-iteration panel, the dynamics are made persistent,
\[
  S_{t+1}=0.92S_t+\alpha_{A_t}+0.03\sin(S_t)+\sigma_\varepsilon\varepsilon_t,
\]
with \(\gamma=0.985\) and horizon \(120\). Rewards are smooth but long-horizon:
\[
  r_{\mathrm{fin}}(s,a)
  =
  0.8+0.9\tanh(s_1)+0.25s_2+\rho_a
  +0.15\tanh(s_1+0.5s_2)+s^\top\eta_a .
\]
Random-feature FQE is stopped after \(K=4\) Bellman iterations, so its values
retain useful ranking information but are under-scaled relative to the
long-horizon target.

In the temporal reward-shift panel, the old and current environments share
dynamics, policies, and state distribution. Only the reward changes:
\[
  r_{\mathrm{current}}(s,a)=2.0+1.5\,r_{\mathrm{old}}(s,a).
\]
We train random-feature FQE on a large old-regime behavior batch
(\(n_{\mathrm{old}}=2000\)), fit the calibration map on a small recent
current-regime batch (\(n_{\mathrm{recent}}=100\)), and evaluate on independent
current-regime ground-truth values and Bellman-calibration transitions. A same-size
current-regime retrain comparator uses the same learner.

In the monotone-misspecification panel, the nonlinear reward is transformed by
a monotone saturation,
\[
  r_{\mathrm{mono}}(s,a)
  =
  \operatorname{sign}\{r_0(s,a)\}\sqrt{|r_0(s,a)|+1}.
\]
The linear FQE value prediction remains sufficiently ordered to be useful but has a
nonlinear value-scale distortion. This is the setting where isotonic and
histogram calibration should improve over an affine correction.

\paragraph{Baseline value predictors.}
Each method first fits a \(Q\)-function from the offline behavior-policy data
and then forms
\[
  \widehat V(s)=\sum_a \pi(a\mid s)\widehat Q(s,a).
\]
The experiments use random-feature FQE and linear FQE because they give
transparent error mechanisms: finite Bellman-iteration bias and affine or
monotone value-scale misspecification.

\paragraph{Cross-calibration.}
For fold \(k\), the \(Q\)-learner is fit on the other folds. The held-out fold
supplies \((\widehat V^{(-k)}(S_i),\widehat V^{(-k)}(S_i'),R_i)\). These
out-of-fold tuples are pooled to fit a single value-space calibration map by
importance-weighted fitted value iteration. At evaluation time, calibrated
values are aggregated as
\[
  \widehat V_{\mathrm{cal}}(s)
  =
  \mathrm{median}_{k=1,\ldots,K_{\mathrm{fold}}}
  g\!\left(\widehat V^{(-k)}(s)\right).
\]
The matched uncalibrated comparator is a \(Q\)-learner refit on all available
training data.

\paragraph{Value-space fitted value iteration.}
Let \(\omega_i=\pi(A_i\mid S_i)/b(A_i\mid S_i)\). In all reported runs the
action-ratio weights are clipped at 20 and normalized to mean one within each
calibration fit. Starting from the identity map, calibration iterates
\[
  g_{t+1}\in\arg\min_g\sum_i \omega_i
  \left[g(\widehat V(S_i))-\{R_i+\gamma g_t(\widehat V(S_i'))\}\right]^2.
\]
The implemented calibrators are weighted affine regression, weighted histogram
binning, weighted isotonic regression, and a weighted isotonic-histogram
hybrid.

\paragraph{Metrics.}
The scalar OPE estimate is the mean of \(g(\widehat V(S_0))\) over independent
initial states. The main prediction metric is held-out value-function MSE,
\[
  \mathbb{E}\!\left[\{g(\widehat V(S))-V^\pi(S)\}^2\right],
\]
estimated using independent target-policy ground-truth values. Bellman calibration
error is estimated on independent Bellman-evaluation transitions by binning
\(g(\widehat V(S))\) into quantile bins and averaging the weighted squared gaps
between bin-mean predictions and bin-mean Bellman outcomes
\[
  r_0(S,A)+\gamma g(\widehat V(S')) .
\]
We use independent evaluation batches  of size 50,000 and 50 bins.  


 
 



\section{Additional details and extensions}
 



\subsection{Hybrid Isotonic--Histogram Calibration}
\label{appendix::hybrid}

Isotonic iterated Bellman calibration adaptively pools nearby predicted values
through the monotonicity constraint. This adaptivity is useful in practice, but
it complicates the estimation-error analysis because applying
Algorithm~\ref{alg::cal-gen} with \(\mathcal G=\mathcal G_{\mathrm{iso}}\) may
produce a different data-dependent partition at each iteration. The iterates
therefore need not correspond to a single aggregated Bellman operator.

The hybrid isotonic--histogram procedure separates partition selection from
Bellman iteration. It first performs one isotonic regression of the initial
Bellman target as a function of \(\hat v(S)\), extracts the flat regions of the
fitted isotonic function as bins, and then runs histogram-based Bellman
calibration on this fixed partition. Thus the method retains the adaptivity of
isotonic calibration while restoring the fixed aggregated-operator structure
used in the histogram analysis.

\begin{algorithm}[t]
\caption{\textbf{Iso--Hist Iterated Bellman Calibration}}
\label{alg::iso-hist}
\begin{algorithmic}[1]

\BCINPUT Value predictor $\hat{v}$, calibration data $\mathcal C_n$,
Bellman target-estimate map $v\mapsto\widehat Y_v$, iterations $K$

\STATE {\small\textbf{Stage 1: fit outcome-adaptive partition}}
\STATE \hspace{1em} $\widehat Y_i^{(0)} \gets \widehat Y_{\hat v}(O_i)$
\STATE \hspace{1em} $g_{\mathrm{iso}} \gets 
       \arg\min_{g \in \mathcal G_{\mathrm{iso}}}
       \sum_{i=1}^n (\widehat Y_i^{(0)} - g(\hat{v}(S_i)))^2$
\STATE \hspace{1em} Extract bins $\hat I_1,\dots,\hat I_{\hat B}$
       from flat regions of $g_{\mathrm{iso}}$
\STATE \hspace{1em} Define $\widehat{\mathcal G}_{\hat B} :=
       \{g:g  \text{ is constant on each }\hat I_b\}$

\STATE {\small\textbf{Stage 2: histogram calibration}}
\STATE \hspace{1em} Apply Algorithm~\ref{alg::cal-gen} with class 
       $\widehat{\mathcal G}_{\hat B}$ to obtain $V_K$

\BCOUTPUT $V_K$

\end{algorithmic}
\end{algorithm}

Once the partition is fixed, the guarantees of Section~\ref{sec:hist} apply
directly, provided the number of bins produced by the initial isotonic step
admits a deterministic high-probability bound. More generally, the same
argument applies to other data-adaptive binning methods that produce a fixed
partition before the Bellman iterations begin.

\begin{theorem}[Calibration and Estimation Error for Alg.~\ref{alg::iso-hist}]
\label{theorem::regretiso}
Let $V_K$ denote the output of Algorithm~\ref{alg::iso-hist}. Suppose
the isotonic regression step produces at most $\hat B$ bins, where 
$\hat B \le B_n$ for some deterministic sequence $B_n$ with probability at least $1 - \delta$. Assume \ref{cond::C1}--\ref{cond::Ccoverage} and \ref{cond::C3} with $B = B_n$.
Then, with probability at least $1 - 2\delta$, the guarantees in
Theorems~\ref{theorem::Calerrorhist} and~\ref{theorem::regretHist} hold for
$V_K$ with $B := B_n$.
\end{theorem}

Existing isotonic regression theory suggests that the effective number of
induced bins satisfies \(\widehat B=O(n^{1/3})\) when the calibration function
is continuously differentiable \citep{deng2021confidence}; this rate can also
be enforced by directly constraining the number of constant segments
\citep{van2024self}. For the approximation term in
Theorem~\ref{theorem::regretHist} to remain comparable to the initial error,
the induced histogram class must also approximate the identity map on the range
of \(\hat v\). This is plausible when \(\hat v\) is already close to \(v_0\);
when \(\hat v\) is substantially miscalibrated, isotonic regression may instead
pool poorly predicted regions into a common calibrated value.

\subsection{Alternative Calibration Targets}
\label{app:alternative-calibration-targets}

The Bellman calibration criterion studied in the main text calibrates a
candidate value predictor \(v\) with respect to its own Bellman target,
\[
E_\pi\!\left[R+\gamma v(S')\mid v(S)\right].
\]
A stronger alternative would instead calibrate \(v\) directly against the true
value function \(v_0\). Specifically, we say that \(v\) is
\emph{strongly Bellman calibrated} if
\[
v(S)
=
E\!\left[v_0(S)\mid v(S)\right]
\quad\text{almost surely}.
\]
Equivalently, among states assigned the same predicted value, the prediction
agrees on average with the true long-run return under \(\pi\).

This target is conceptually appealing, but substantially less practical. It
requires estimating the conditional mean of the unknown quantity \(v_0(S)\)
given the predicted value \(v(S)\). In general, this conditional mean is not directly
available from one-step transitions alone and is typically no easier to estimate
than the underlying value function itself
\citep{kallus2020double,kallus2022efficiently}. We therefore focus on the
weaker but directly attainable Bellman-target calibration criterion developed
in the main text.



 

\section{Experimental results and info}
\label{appendx::exp}
\subsection{Synthetic CRM environment}
 

We evaluate all methods in a synthetic customer-relationship-management (CRM)
Markov decision process (MDP) that mimics monthly retention and revenue dynamics
for a subscription service. The state $s_t \in \mathbb{R}^6$ encodes tenure in
months, engagement, fatigue, value segment, price sensitivity, and an indicator
of whether the customer is active. At each time $t$, the agent selects one of
three actions: no promotion, light promotion, or strong promotion. Rewards
correspond to monthly revenue after discounts and are zero after churn.

Transitions capture key CRM effects. Churn probability depends on tenure,
engagement, fatigue, and action via a logistic model; visit probability depends
on engagement, fatigue, and action; and revenue conditional on visit scales a
baseline value segment by an action-dependent uplift and a price-sensitivity
effect, with log-normal noise. Engagement decays over time but is boosted by
successful promotions, while fatigue increases with promotion intensity and
decays slowly otherwise. Once a customer churns or reaches a maximum tenure of
60 months, the process enters an absorbing state.

We simulate offline datasets of $n_{\text{cust}} = 50{,}000$ customers over a
horizon of $T = 24$ months with discount factor $\gamma = 0.99$. Behavior data are
generated by a fixed heuristic policy that sends light promotions by default,
suppresses promotions for highly engaged and fatigued customers, and
occasionally sends strong promotions to low-engagement, high-value customers. We
evaluate off-policy value estimation for a deterministic \emph{aggressive}
revenue-seeking target policy that sends strong promotions to low-engagement,
high-sensitivity customers and avoids promotions for highly engaged, high-value
customers.

Ground-truth values $V^\pi(s_0)$ are estimated via Monte Carlo rollouts in the
same environment, and we report $(1-\gamma)$-scaled RMSE between the estimated and
Monte Carlo values at initial states, averaged across 50 independent runs.


\paragraph{Base estimators.}
We consider three base estimators: boosted FQE, linear FVI, and neural FVI.
The linear model uses a hand-crafted semi-structured feature map consisting of
the six raw state variables, normalized and log-transformed tenure, squared
terms for engagement, fatigue, value segment, and price sensitivity, three
state-only interaction terms, and additive action one-hot indicators, with no
state-action interactions. The neural model takes the raw state concatenated
with the action as input and uses a feedforward ReLU network with three hidden
layers of width 256. Both FVI estimators are trained by iterating one-step
Bellman regressions: the linear model uses ridge-regularized least squares,
while the neural model is trained for 100 outer iterations using Adam with
learning rate \(0.005\), one epoch per iteration, and mini-batch size 2000.
We do not tune the number of neural iterations, since part of the goal is to
illustrate a realistic setting in which a base estimator remains poorly
calibrated. Selected neural iterates are saved for the snapshot comparisons.
Both FVI estimators use an internal 80/20 train-validation split. The boosted
baseline is implemented as a LightGBM FQE model trained on the concatenated
state-action input. Under sample splitting, base estimators are fit on the
training subset and Bellman calibrators on the held-out calibration subset; we
also retrain each base estimator on the full dataset to obtain the
\emph{Full Raw} baselines. 



For the calibration methods, we use 1000 fitted-value-iteration updates.


 

 \subsubsection{Neural snapshots under sample splitting}

Tables~\ref{tab:neural-snap-split-50k-p05} and \ref{tab:neural-snapshots-50k-ptrain-075} report neural FVI snapshot performance under sample splitting for the additional training proportions \(p_{\mathrm{train}}=0.5\) and \(p_{\mathrm{train}}=0.75\). Together, they show how the effect of calibration evolves over training and how it depends on the amount of data reserved for fitting versus calibration.

\begin{table}[h]
\centering
\small
\begin{tabular}{r cccc}
\toprule
Iter & Full Raw & Iso & Histogram & Hybrid Iso \\
\midrule
10  & 1.309$\pm$0.06 & \textbf{0.581$\pm$0.05} & 0.582$\pm$0.05 & 0.581$\pm$0.05 \\
25  & 0.491$\pm$0.13 & \textbf{0.456$\pm$0.06} & 0.456$\pm$0.06 & \textbf{0.456$\pm$0.06} \\
50  & \textbf{0.320$\pm$0.07} & 0.381$\pm$0.05 & 0.382$\pm$0.05 & 0.382$\pm$0.05 \\
100 & \textbf{0.316$\pm$0.05} & 0.350$\pm$0.05 & 0.351$\pm$0.05 & 0.350$\pm$0.05 \\
\bottomrule
\end{tabular}
\caption{Neural snapshot performance with sample splitting for \(n=50{,}000\) and training proportion \(p_{\mathrm{train}}=0.5\). Lower is better.}
\label{tab:neural-snap-split-50k-p05}
\end{table}

\begin{table}[h]
\centering
\small
\begin{tabular}{c cccc}
\toprule
Iteration & Full Raw & Iso & Histogram & Hybrid Iso \\
\midrule
10  & 1.28$\pm$0.05 & 0.52$\pm$0.05 & \textbf{0.52$\pm$0.05} & 0.52$\pm$0.05 \\
25  & 0.44$\pm$0.12 & \textbf{0.38$\pm$0.03} & 0.38$\pm$0.04 & 0.38$\pm$0.04 \\
50  & 0.34$\pm$0.06 & \textbf{0.33$\pm$0.04} & 0.33$\pm$0.04 & 0.33$\pm$0.04 \\
100 & 0.33$\pm$0.07 & \textbf{0.31$\pm$0.04} & 0.31$\pm$0.04 & 0.31$\pm$0.04 \\
\bottomrule
\end{tabular}
\caption{Neural FVI snapshot results for \(n=50{,}000\) with \(p_{\mathrm{train}}=0.75\). Lower is better.}
\label{tab:neural-snapshots-50k-ptrain-075}
\end{table}


 


\subsubsection{Comparing to the uncalibrated (raw) model without retraining on the full data}
Table~\ref{tab:main-split} compares calibrated estimators to the corresponding raw model fit only on the training split, so any improvements reflect calibration rather than retraining on the full dataset. This comparison is relevant when retraining on all available data is costly or when a gold-standard calibration dataset is available, such as recent trajectories or on-policy data, alongside a larger historical dataset used for initial fitting. Table~\ref{tab:neural-snap-split} further examines this comparison for the neural method across training iterations, illustrating how the effect of calibration evolves over the course of fitting.


\begin{table}[h]
\centering
\small
\begin{tabular}{r l ccc}
\toprule
$n$ & Model & Raw & Iso & Hybrid Iso \\
\midrule
\multirow{3}{*}{10{,}000}
& Boost & 0.681$\pm$0.1 & \textbf{0.671$\pm$0.1} & 0.697$\pm$0.2 \\
& Lin  & 0.640$\pm$0.04 & \textbf{0.612$\pm$0.05} & 0.641$\pm$0.05 \\
& Neur  & 0.582$\pm$0.2 & 0.550$\pm$0.2 & \textbf{0.520$\pm$0.1} \\
\midrule
\multirow{3}{*}{50{,}000}
& Boost & 0.433$\pm$0.06 & \textbf{0.429$\pm$0.06} & 0.438$\pm$0.08 \\
& Lin  & 0.625$\pm$0.02 & \textbf{0.614$\pm$0.03} & 0.616$\pm$0.03 \\
& Neur  & 0.419$\pm$0.1 & 0.383$\pm$0.08 & \textbf{0.374$\pm$0.08} \\
\midrule
\multirow{3}{*}{100{,}000}
& Boost & 0.360$\pm$0.05 & 0.358$\pm$0.05 & \textbf{0.353$\pm$0.07} \\
& Lin  & 0.616$\pm$0.02 & \textbf{0.605$\pm$0.02} & 0.606$\pm$0.02 \\
& Neur  & 0.379$\pm$0.09 & 0.351$\pm$0.05 & \textbf{0.342$\pm$0.05} \\
\bottomrule
\end{tabular}
\caption{Main results with sample splitting at training proportion \(p_{\mathrm{train}}=0.5\). The raw estimator is fit only on the training split to isolate the effect of calibration.}
\label{tab:main-split}
\end{table}




\begin{table}[h]
\centering
\small
\begin{tabular}{r cccc}
\toprule
Iter & Raw & Iso & Quantile & Hybrid Iso \\
\midrule
10  & 1.424$\pm$0.1 & 1.330$\pm$0.1 & \textbf{0.735$\pm$0.2} & 0.737$\pm$0.2 \\
25  & 0.665$\pm$0.1 & 0.646$\pm$0.1 & 0.577$\pm$0.1 & \textbf{0.571$\pm$0.1} \\
50  & 0.612$\pm$0.2 & 0.585$\pm$0.2 & 0.555$\pm$0.1 & \textbf{0.547$\pm$0.1} \\
100 & 0.582$\pm$0.2 & 0.550$\pm$0.2 & 0.524$\pm$0.1 & \textbf{0.520$\pm$0.1} \\
\bottomrule
\end{tabular}
\caption{Neural FVI snapshot performance with sample splitting at training proportion \(p_{\mathrm{train}}=0.5\). The raw estimator is fit using only the training split, so differences reflect the effect of calibration rather than access to additional data.}\label{tab:neural-snap-split}
\end{table}






 


 

 




\section{Stationary Measures and Contraction Results}
\label{appendix::stationary}

\begin{lemma}[Bellman contraction under a stationary measure]
\label{lemma::contraction}
Let $P$ be a Markov operator with stationary distribution $\mu$ (i.e., $\mu=\mu P$),
and define the Bellman operator $\mathcal T f := r + \gamma P f$.
Then $P$ is nonexpansive in $L^2(\mu)$:
\[
\|P(f-g)\|_{2,\mu} \;\le\; \|f-g\|_{2,\mu}.
\]
Consequently, $\mathcal T$ is a $\gamma$-contraction:
\[
\|\mathcal T (f-g)\|_{2,\mu}
\;\le\;
\gamma\,\|f-g\|_{2,\mu}.
\]
\end{lemma}
\begin{proof}
Because $\mu$ is stationary for $P$, the operator $P$ is nonexpansive in
$L^2(\mu)$. Jensen's inequality implies
\[
(P h)^2 \le P(h^2) \qquad\text{pointwise}.
\]
Integrating both sides with respect to $\mu$ and using $\mu = \mu P$ yields
\[
\|P h\|_{2,\mu}^2
= \int (P h)^2 \, d\mu
\le \int P(h^2)\, d\mu
= \int h^2\, d\mu
= \|h\|_{2,\mu}^2.
\]
Now apply this with $h = f - g$. Since the reward function $r$ cancels,
\[
\mathcal{T}f - \mathcal{T}g
= \gamma\, P(f-g).
\]
Therefore,
\[
\|\mathcal{T} f - \mathcal{T} g\|_{2,\mu}
= \gamma \|P(f-g)\|_{2,\mu}
\le \gamma \|f-g\|_{2,\mu},
\]
which proves the claim.
\end{proof}

\begin{lemma}[Stationarity under coarsening]
\label{lemma:aggregated-stationary}
Let $P$ be a Markov kernel on $\mathcal S$ and let $\lambda$ be stationary for $P$
(i.e., $\lambda P = \lambda$).
For any measurable coarsening map $g:\mathcal S\to\mathbb R$, define
\[
P_g f(s)
:=
\mathbb{E}[\,f(S') \mid g(S)=g(s)\,],
\]
where $(S,S')$ has joint law $\lambda(ds)\,P(s,ds')$.
Then $\lambda$ is stationary for $P_g$.
\end{lemma}

\begin{proof}
Let $(S,S')\sim \lambda(ds)\,P(s,ds')$, so $S\sim\lambda$ and $S'\mid S\sim P(S,\cdot)$.
Then
\begin{align*}
\int P_g f(s)\,\lambda(ds)
&=
\mathbb{E}_{\lambda,P}\!\left[P_g f(S)\right] \\[0.3em]
&=
\mathbb{E}_{\lambda,P}\!\left[\mathbb{E}_{\lambda,P}\!\left[f(S') \mid g(S)\right]\right] \\[0.3em]
&=
\mathbb{E}_{\lambda,P}[f(S')].
\end{align*}
Since $\lambda$ is stationary for $P$, the marginal of $S'$ is again $\lambda$, hence
$\mathbb{E}_{\lambda,P}[f(S')] = \int f(s)\,\lambda(ds)$.
\end{proof}



\section{Calibration--Refinement Proofs}

 
In this section, set
\[
\mathcal B_{\pi,\hat v}q
:=
\Pi_{\hat v}\mathcal T_\pi q,
\qquad
V_{\hat v}^\star\in\mathcal V_{\hat v},
\qquad
V_{\hat v}^\star
=
\mathcal B_{\pi,\hat v}V_{\hat v}^\star .
\]
 

 


\begin{theorem}[Sharper calibration--refinement bound]
\label{theorem::calrefineOracle}
Under \ref{cond::A1},
\[
\|\hat v-v_0\|_{2,\rho}
\le
\|\hat v-V_{\hat v}^\star\|_{2,\rho}
+
\frac{\sqrt{\mathcal C_{\hat v,\gamma}}}{1-\gamma}
\min_{\theta:\mathbb R\to\mathbb R}
\|\theta\circ\hat v-v_0\|_{2,\rho}.
\]
\end{theorem}

\begin{proof}[Proof of Theorem \ref{theorem::calrefineOracle}]
Define \(P_{\pi,\hat v} := \Pi_{\hat v}P_{\pi}\) and
\(r_{\pi,\hat{v}} := \Pi_{\hat v} r_{\pi}\). Because \(v_0\) and
\(V_{\hat v}^\star\) satisfy the fixed-point equations
\[
v_0 = \mathcal T_\pi v_0,
\qquad
V_{\hat v}^\star = \mathcal B_{\pi,\hat v} V_{\hat v}^\star
= r_{\pi,\hat{v}} + \gamma\,P_{\pi,\hat v}V_{\hat v}^\star,
\]
we begin by writing
\[
V_{\hat v}^\star - v_0
=
\bigl( \mathcal B_{\pi,\hat v}V_{\hat v}^\star
      - \mathcal B_{\pi,\hat v}v_0 \bigr)
+
\bigl( \mathcal B_{\pi,\hat v}v_0 - \mathcal{T}_\pi v_0 \bigr).
\]
Using linearity on differences,
\[
\mathcal B_{\pi,\hat v}V_{\hat v}^\star
-
\mathcal B_{\pi,\hat v}v_0
=
\gamma\,P_{\pi,\hat v}(V_{\hat v}^\star - v_0),
\]
so
\[
V_{\hat v}^\star - v_0
=
\gamma\,P_{\pi,\hat v}(V_{\hat v}^\star - v_0)
+
\bigl(\mathcal B_{\pi,\hat v}v_0 - v_0\bigr),
\]
where we used \(\mathcal T_\pi v_0 = v_0\) for the second term. Since
\(\mathcal B_{\pi,\hat v}v_0=\Pi_{\hat v}v_0\), this gives
\[
V_{\hat v}^\star-v_0
=
\gamma P_{\pi,\hat v}(V_{\hat v}^\star-v_0)
+
(\Pi_{\hat v}-I)v_0 .
\]
Let \(a:=(\Pi_{\hat v}-I)v_0\). Iterating the display yields
\[
V_{\hat v}^\star-v_0
=
\sum_{t=0}^\infty \gamma^t P_{\pi,\hat v}^t a .
\]
For each \(t\ge0\), set
\[
c_t
:=
\left\|
\frac{d(\rho P_{\pi,\hat v}^t)}{d\rho}
\right\|_\infty .
\]
Jensen's inequality implies
\[
\|P_{\pi,\hat v}^t a\|_{2,\rho}
\le
\sqrt{c_t}\|a\|_{2,\rho}.
\]
Therefore,
\[
\|V_{\hat v}^\star-v_0\|_{2,\rho}
\le
\frac{\sqrt{\mathcal C_{\hat v,\gamma}}}{1-\gamma}
\|\Pi_{\hat v}v_0-v_0\|_{2,\rho}.
\tag{$\star$}
\]
Since \(\Pi_{\hat v}\) is the \(L^2(\rho)\)-projection onto
\(\{\theta\circ\hat v:\theta:\mathbb R\to\mathbb R\}\),
\[
\|\Pi_{\hat v}v_0-v_0\|_{2,\rho}
=
\min_{\theta:\mathbb R\to\mathbb R}
\|\theta\circ\hat v-v_0\|_{2,\rho}.
\]
Finally, by the triangle inequality,
\[
\|\hat{v} - v_0\|_{2,\rho}
\le
\|\hat{v} - V_{\hat v}^\star\|_{2,\rho}
+
\|V_{\hat v}^\star - v_0\|_{2,\rho},
\]
Combining this display with \((\star)\) yields the stated result.
\end{proof}

\begin{proof}[Proof of Theorem \ref{theorem::calrefine}]
Let \(P_{\pi,\hat v} := \Pi_{\hat v}P_{\pi}\).
As in the proof of Theorem~\ref{theorem::calrefineOracle},
\[
\|V_{\hat v}^\star-v_0\|_{2,\rho}
\le
\frac{\sqrt{\mathcal C_{\hat v,\gamma}}}{1-\gamma}
\min_{\theta:\mathbb R\to\mathbb R}
\|\theta\circ\hat v-v_0\|_{2,\rho}.
\tag{$\star$}
\]
Moreover,
\[
\hat v - V_{\hat v}^\star
=
\hat v-\mathcal B_{\pi,\hat v}\hat v
+
\gamma P_{\pi,\hat v}(\hat v-V_{\hat v}^\star),
\]
and therefore
\[
\hat v - V_{\hat v}^\star
=
\sum_{t=0}^\infty \gamma^t
P_{\pi,\hat v}^t
\{\hat v-\mathcal B_{\pi,\hat v}\hat v\}.
\]
Since \(\mathcal B_{\pi,\hat v}\hat v=\Gamma_0(\hat v)\), the
concentrability argument above yields
\[
\|\hat v - V_{\hat v}^\star\|_{2,\rho}
\le
\frac{\sqrt{\mathcal C_{\hat v,\gamma}}}{1-\gamma}
\mathrm{Cal}_{\ell^2}(\hat v).
\tag{$\dagger$}
\]
The triangle inequality and displays \((\star)\) and \((\dagger)\) give the
claim.
\end{proof}



\section{Proof of Theorem \ref{theorem::DRpseudo}}

\begin{proof}[Proof of Theorem \ref{theorem::DRpseudo}]
Fix $v$ and treat the nuisance estimates as fixed functions. By definition,
\[
\widehat{\mathcal T}_\pi v(s)
=
\mathbb{E}\Big[
(\pi \hat q_v)(S)
+
\widehat w_{\pi}(A \mid S)
\big\{
  R + \gamma\, v(S') - \hat{q}_v(S,A)
\big\}
\;\Big|\; S = s
\Big].
\]
Since $(\pi \hat q_v)(S)$ depends only on $S$,
\[
\mathbb{E}\big[(\pi \hat q_v)(S)\mid S=s\big] = (\pi \hat q_v)(s).
\]
For the second term, using the definition of $b_0$ and $q_v := r_0 + \gamma P v$,
\begin{align*}
&\mathbb{E}\Big[
\widehat w_{\pi}(A \mid S)
\big\{
  R + \gamma\, v(S') - \hat{q}_v(S,A)
\big\}
\;\Big|\; S = s
\Big] \\
&\quad=
b_0\Big\{
\widehat w_\pi(\cdot \mid s)\,
\mathbb{E}\big[
R + \gamma\, v(S') - \hat q_v(s,\cdot)
\mid S = s, A = \cdot
\big]
\Big\}(s) \\
&\quad=
b_0\big\{\widehat w_\pi (q_v - \hat q_v)\big\}(s).
\end{align*}
Hence
\[
\widehat{\mathcal T}_\pi v(s)
=
(\pi \hat q_v)(s)
+
b_0\big\{\widehat w_\pi (q_v - \hat q_v)\big\}(s).
\]

Subtracting $(\mathcal T_\pi v)(s) = (\pi q_v)(s)$ yields
\begin{align*}
\widehat{\mathcal T}_\pi v(s) - (\mathcal T_\pi v)(s)
&=
(\pi \hat q_v - \pi q_v)(s)
+
b_0\big\{\widehat w_\pi (q_v - \hat q_v)\big\}(s) \\
&=
(\pi (\hat q_v - q_v))(s)
-
b_0\big\{\widehat w_\pi (\hat q_v - q_v)\big\}(s).
\end{align*}
Using $(\pi f)(s) = b_0\{w_\pi f\}(s)$ with $f = \hat q_v - q_v$ gives
\[
(\pi (\hat q_v - q_v))(s)
=
b_0\{w_\pi (\hat q_v - q_v)\}(s),
\]
so
\[
\widehat{\mathcal T}_\pi v(s) - (\mathcal T_\pi v)(s)
=
b_0\big\{(w_\pi - \widehat w_\pi)(\hat q_v - q_v)\big\}(s).
\]
Taking conditional expectations given \(v(S)=t\) proves the claim.
\end{proof}



\section{Notation and Maximal Inequalities}

Let $P_0$ denote the joint distribution of $(S, A, R, S')$ induced by the
behavior policy, and let $P_n$ denote the empirical measure of the calibration
sample $\mathcal{C}_n$.



We define empirical \(L^2\) norms with respect to the state and next-state samples.  For any state function \(f\), let
\[
\|f\|_{n,S}
:= 
\Big( \tfrac{1}{n} \sum_{i=1}^{n} f(S_i)^2 \Big)^{1/2},
\quad
\|f\|_{n,S'}
:= 
\Big( \tfrac{1}{n} \sum_{i=1}^{n} f(S_i')^2 \Big)^{1/2},
\]
where \(S_i\) and \(S_i'\) denote the observed states and next states, respectively, in the calibration sample \(\mathcal{C}_n\).



 





Further define
\[
\widehat{\mathcal{G}}
:= 
\Big\{ (f_1 - f_2)\big(\widehat Y_{f_2} - f_2\big)
:\;
f_1, f_2 \in \mathcal H_B(\hat v) \Big\}.
\]
By assumption, both \(\hat{v}\) and the Bellman target-estimate map
\(v\mapsto\widehat Y_v\) are fixed conditional on the training data, which is
independent of the calibration sample \(\mathcal{C}_n\). Consequently, the
classes \(\mathcal H_B(\hat v)\) and \(\widehat{\mathcal{G}}\) are
non-random conditional on the training dataset.

For any distribution \(Q\) and any uniformly bounded function class \(\mathcal{F}\), let  
\(N(\varepsilon, \mathcal{F}, L^2(Q))\) denote the \(\varepsilon\)-covering number of \(\mathcal{F}\) under the \(L^2(Q)\) norm \citep{van1996weak}.  
Define the uniform entropy integral of \(\mathcal{F}\) by
\begin{equation*}
\mathcal{J}(\delta, \mathcal{F})
:=
\int_{0}^{\delta}
\sup_{Q}
\sqrt{\log N(\epsilon, \mathcal{F}, L^2(Q))}\, d\epsilon ,
\end{equation*}
where the supremum is taken over all discrete probability distributions \(Q\).

Finally, for two quantities \(x\) and \(y\), we write \(x \lessapprox y\) to mean that \(x\) is bounded above by \(y\) up to a universal constant that depends only on global constants appearing in our conditions.


\subsection{Local maximal inequality}

 Let $O_1,\ldots,O_n \in \mathcal{O}$ be independent random variables. For any function $f:\mathcal{O} \to \mathbb{R}$, define
\begin{align}
   \|f\|_{2,\bar P} := \sqrt{\frac{1}{n} \sum_{i=1}^n\mathbb{E}[f(O_i)^2]}.
\end{align}
 

 


For a star-shaped class of functions $\mathcal{F}$ and a radius $\delta \in (0,\infty)$,
define the localized Rademacher complexity
\[
\mathcal{R}_n(\mathcal{F}, \delta)
:=
\mathbb{E}\left[
\sup_{\substack{f \in \mathcal{F} \\\|f\|_{2,\bar P} \le \delta}}
\frac{1}{n} \sum_{i=1}^n \epsilon_i f(O_i)
\right],
\]
where $\epsilon_i$ are i.i.d.\ Rademacher random variables.

The following lemma provides a local maximal inequality and restates Lemma~11 of \cite{foster2023orthogonal}.


 


\begin{lemma}[Local maximal inequality]\label{lemma:loc_max_ineq}
Let $\mathcal{F}$ be a star-shaped class of functions satisfying
$\sup_{f\in\mathcal{F}} \|f\|_{\infty} \le M$.
Let $\delta > 0$ satisfy the critical radius condition
$\mathcal{R}_n(\mathcal{F},\delta) \le \delta^2$.
Suppose further that
$n^{-1/2}\sqrt{\log\log(1/\delta)} = o(\delta)$. Then there exists a universal constant $C>0$ such that, for all $u \ge 1$, with
probability at least $1 - e^{-u^2}$, every $f \in \mathcal{F}$ satisfies
\[
\frac{1}{n} \sum_{i=1}^n
\bigl(f(O_i) - \mathbb{E}[f(O_i)]\bigr)
\;\le\;
C\Bigl(
\delta^2
+
\delta\|f\|_{2,\bar P}
+
\frac{u\|f\|_{2,\bar P}}{\sqrt{n}}
+
\frac{M u^2}{n}
\Bigr).
\]
\end{lemma}
\begin{proof}
This is Lemma~11 of \citet{foster2023orthogonal}, specialized to the notation
above.
\end{proof}

The following lemma bounds the localized Rademacher complexity in terms of the uniform entropy integral and is a direct consequence of Theorem~2.1 of \citet{van2011local}.

 

\begin{lemma}\label{lemma:local_rademacher_entropy}
Let \(\mathcal{F}\) be a star-shaped class of functions such that
\(\sup_{f\in\mathcal F}\|f\|_\infty \le M\). Then, for every \(\delta>0\),
\[
\mathcal{R}_n(\mathcal{F}, \delta)
\;\lesssim\;
\frac{1}{\sqrt n}\,\mathcal{J}(\delta,\mathcal{F})
\left(1+\frac{\mathcal{J}(\delta,\mathcal{F})}{\delta \sqrt n}\right),
\]
where the implicit constant depends only on \(M\).
\end{lemma}
\begin{proof}
This bound follows directly from the argument in the proof of
Theorem~2.1 of \citet{van2011local}; see in particular the step where
the local Rademacher complexity is controlled by the uniform entropy
integral for star-shaped classes.
\end{proof}
 

 

\section{Proof of Theorem \ref{theorem::Calerrorhist}}
 
 

\subsection{Technical lemmas}

 



\begin{lemma}
\label{lemma::entropynumbers}
Under our conditions,
\[
\mathcal{J}\left(\delta, \widehat{\mathcal{G}}\right)
\;\lesssim\;
\delta \sqrt{B \log(1/\delta)},
\]
where the implicit constant is independent of \(B\).
\end{lemma}
\begin{proof}
 

By Condition~\ref{cond::C2},
\[
\|\hat q_f-\hat q_g\|_{n,S,A}
\le
L\|f-g\|_{n,S'}
\]
almost surely.
For each $s$, by Jensen's inequality,
\[
|\pi h(s)|^2=\Big|\sum_a \pi(a\mid s)h(a,s)\Big|^2
\le \sum_a \pi(a\mid s)|h(a,s)|^2,
\]
so $\|\pi h\|_{n,S}\le \|h\|_{n,S,A}$. Hence, we also have
\[
\|\pi(\hat q_f-\hat q_g)\|_{n,S}
\le \|\hat q_f-\hat q_g\|_{n,S,A}
\le L \|f - g\|_{n, S'} .
\]
Hence, by the definition of \(\widehat Y_f\) and boundedness of
\(\widehat w_\pi\),
\[
\|\widehat Y_f - \widehat Y_g\|_{L^2(P_n)} 
\lesssim  
\|f - g\|_{n,S'}.
\]
Take $f, g \in \widehat{\mathcal{G}}$ with 
$f = (f_1 - f_2)\big(\widehat Y_{f_2} - f_2\big)$ 
and 
$g = (g_1 - g_2)\big(\widehat Y_{g_2} - g_2\big)$ 
for $f_1, f_2, g_1, g_2 \in \mathcal H_B(\hat v)$. 
By Lipschitz continuity of multiplication and boundedness of nuisances, we have
\[
\|f - g\|_{L^2(P_n)} 
\lesssim  
\|f_1 - g_1\|_{n, S} + \|f_2 - g_2\|_{n, S}+ \|f_2 - g_2\|_{n, S'}.
\]
Taking the supremum over all discrete distributions $Q$, we obtain
\[
\|f - g\|_{L^2(P_n)} 
\lesssim  
\sup_Q \|f_1 - g_1\|_{L^2(Q)} 
+ \sup_Q \|f_2 - g_2\|_{L^2(Q)} .
\]
Hence, by preservation of entropy integrals in \citet{van1996weak},
\[
\log N(\varepsilon, \widehat{\mathcal{G}}, L^2(P_n)) 
\lesssim 
\sup_Q \log N(\varepsilon, \mathcal H_B(\hat v), L^2(Q)) .
\]
Taking the supremum over discrete distributions $Q$ on both sides yields the uniform covering number bound
\[
\sup_Q \log N(\varepsilon, \widehat{\mathcal{G}}, L^2(Q)) 
\lesssim 
\sup_Q \log N(\varepsilon, \mathcal H_B(\hat v), L^2(Q)) .
\]


The class \(\mathcal H_B(\hat v)\) satisfies
\[
\sup_Q \log N\big(\varepsilon, \mathcal H_B(\hat v), L^2(Q)\big)
\;\le\;
\sup_{Q} \log N\big(\varepsilon, \widetilde{\mathcal{F}}_{B},
L^2(Q \circ \hat{v}^{-1})\big),
\]
where \(Q \circ \hat{v}^{-1}\) denotes the pushforward of \(Q\) under
\(\hat{v}\). Hence,
\[
\sup_Q \log N\big(\varepsilon, \mathcal H_B(\hat v), L^2(Q)\big)
\;\le\;
\sup_{Q} \log N\big(\varepsilon, \widetilde{\mathcal{F}}_{B}, L^2(Q)\big),
\]
where, by a slight abuse of notation, the supremum on the right-hand side is
taken over all discrete probability distributions \(Q\) on \(\mathbb{R}\).
The class \(\widetilde{\mathcal{F}}_{B}\) consists of all piecewise-constant
functions on \(\mathbb{R}\) taking at most \(B\) values. Therefore,
\(\widetilde{\mathcal{F}}_{B}\) has VC--subgraph dimension \(O(B)\), and
\citet{van1996weak} implies
\begin{align*}
    \sup_Q \log N(\varepsilon, \mathcal H_B(\hat v), L^2(Q))
&\lesssim
\sup_{Q} \log N\big(\varepsilon, \widetilde{\mathcal{F}}_{B}, L^2(Q)\big)\\
&\lesssim B \log(1/\varepsilon).
\end{align*}
 Consequently, $\mathcal{J}(\delta, \mathcal H_B(\hat v))
\lesssim \delta \sqrt{B \log(1/\delta)}$.
\end{proof}

 


\subsection{Proof of Theorem \ref{theorem::Calerrorhist}}
\begin{proof}[Proof of Theorem \ref{theorem::Calerrorhist}]
Denote $O_i := (S_i, A_i, R_i, S_i')$. The first-order optimality conditions for
$\theta_n^{(K)}$ imply that, for each bin $b \in [B]$,
\[
\frac{1}{n}\sum_{i=1}^n 
\mathbf{1}_{I_b}(\hat{v}(S_i))
\big\{
\widehat Y_{V_{K-1}}(O_i)
- V_K(S_i)
\big\}
= 0.
\]
Hence, for any $f \in \mathcal G_B$, which is a linear combination of these
indicators,
\begin{equation}
\label{eqn::scores}
\frac{1}{n}\sum_{i=1}^n 
f(\hat{v}(S_i))
\big\{
\widehat Y_{V_{K-1}}(O_i)
- V_K(S_i)
\big\}
= 0.
\end{equation}
Moreover, for any function $g:\mathbb{R}\to\mathbb{R}$, we may write
$g \circ V_K = f \circ \hat{v}$ for some $f$, and therefore,
for all such $g$,
\[
\frac{1}{n}\sum_{i=1}^n 
g(V_K(S_i))
\big\{
\widehat Y_{V_{K-1}}(O_i)
- V_K(S_i)
\big\}
= 0.
\]
Noting that \(\Gamma_0(V_K) - V_K\) is of the form
\(g \circ V_K\) for some function \(g\), we have
\begin{align*}
0
&=
\frac{1}{n}\sum_{i=1}^n 
\Big\{
\Gamma_0(V_K)(S_i) - V_K(S_i)
\Big\} \\
&\qquad\qquad\times
\Big\{
\widehat Y_{V_{K-1}}(O_i)
- V_K(S_i)
\Big\}.
\end{align*}
Adding and subtracting \(P_0\) yields
\begin{equation}
\label{eqn::basicineq}
\begin{aligned}
&P_0
\Big\{
\Gamma_0(V_K) - V_K
\Big\}
\Big\{
\widehat Y_{V_K}
- V_K
\Big\}\\
& = (P_0 - P_n)
\Big\{
\Gamma_0(V_K) - V_K
\Big\}
\Big\{
\widehat Y_{V_{K-1}}
- V_K
\Big\}\\
&\quad  + P_0
\Big\{
\Gamma_0(V_K) - V_K
\Big\}
\Big\{
\widehat Y_{V_K}
- \widehat Y_{V_{K-1}}
\Big\}.
\end{aligned}
\end{equation}


We rewrite the left-hand side of \eqref{eqn::basicineq}.
Adding and subtracting and applying the law of total expectation, the calibration error decomposes as
\begin{equation}
\label{eqn::calerrorbound}
\begin{aligned}
&\| \Gamma_0(V_K) - V_K\|^2\\
&=
P_0
\Big\{
\Gamma_0(V_K) - V_K
\Big\}
\Big\{
\widehat Y_{V_K}
- V_K
\Big\} \\
&\quad+
P_0
\Big\{
\Gamma_0(V_K) - V_K
\Big\}
\Big\{
\Gamma_0(V_K)
- \widehat Y_{V_K}
\Big\}.
\end{aligned}
\end{equation}

The first term on the right-hand side can be decomposed as in \eqref{eqn::basicineq}.
By the law of total expectation, the proof of Theorem~\ref{theorem::DRpseudo},
and the Cauchy--Schwarz inequality, the second term satisfies
\begin{align*}
&\left|
P_0
\Big\{
\Gamma_0(V_K) - V_K
\Big\}
\Big\{
\widehat Y_{V_K}
- \Gamma_0(V_K)
\Big\}
\right|
\\[0.4em]
&=
\left|
P_0
\Big\{
\Gamma_0(V_K) - V_K
\Big\}
\Big\{
\widehat Y_{V_K}
- \mathcal{T}_{\pi}(V_K)
\Big\}
\right|
\\[0.4em]
&\le\;
\|\Gamma_0(V_K) - V_K\|\;
\big\|
b_0\{(\widehat w_{\pi} - w_{\pi})\,
(\widehat q_{V_K} - q_{V_K})\}
\big\|_{2,\rho}.
\end{align*}

Hence, \eqref{eqn::basicineq} and \eqref{eqn::calerrorbound} together imply that
\begin{equation}
\label{eqn::basicineq2}
\begin{aligned}
&\|\Gamma_0(V_K) - V_K\|^2\\
&\le\;
(P_0 - P_n)
\Big\{
\Gamma_0(V_K) - V_K
\Big\}
\Big\{
\widehat Y_{V_{K-1}}
- V_K
\Big\}
\\[0.4em]
&\quad+\;
\|\Gamma_0(V_K) - V_K\|\;
\big\|
\widehat{\mathcal T}_\pi V_K
- \widehat{\mathcal T}_\pi V_{K-1}
\big\|
\\[0.4em]
&\quad+\;
\|\Gamma_0(V_K) - V_K\|\;
\big\|
b_0\{(\widehat w_{\pi} - w_{\pi})\,
(\widehat q_{V_K} - q_{V_K})\}
\big\|_{2,\rho}.
\end{aligned}
\end{equation}
where $\widehat{\mathcal T}_\pi v(s)
:= \mathbb{E}[\widehat Y_v(O)\mid S=s]$.
Here, the second term on the right-hand side follows from \eqref{eqn::basicineq}, noting that
\begin{align*} &\left| P_0 \Big\{ \Gamma_0(V_K) - V_K \Big\} \Big\{ \widehat Y_{V_K} - \widehat Y_{V_{K-1}} \Big\} \right| \\[0.4em] &= \left| P_0 \Big\{ \Gamma_0(V_K) - V_K \Big\} \,\mathbb{E}\left[ \widehat Y_{V_K} - \widehat Y_{V_{K-1}} \mid S \right] \right| \\[0.4em] &\le \|\Gamma_0(V_K) - V_K\|\; \Big\| \mathbb{E}\left[ \widehat Y_{V_K} - \widehat Y_{V_{K-1}} \mid S \right] \Big\| \\[0.4em] &= \|\Gamma_0(V_K) - V_K\|\; \big\| \widehat{\mathcal T}_\pi V_K - \widehat{\mathcal T}_\pi V_{K-1} \big\|_{2,\rho}. \end{align*}
 




We now turn to bounding the empirical process term on the right-hand side of \eqref{eqn::basicineq2}. Observe that 
\[
\big(\Gamma_0(V_K) - V_K\big)
\big(\widehat Y_{V_{K-1}} - V_K\big)
\]
lies in a uniformly bounded subset of the class 
\[
\widehat{\mathcal{G}} := \big\{(f_1 - f_2)\big(\widehat Y_{f_2} - f_2\big) : 
f_1, f_2 \in \mathcal H_B(\hat v) \big\},
\]
By Lemma~\ref{lemma::entropynumbers}, it holds that $\mathcal{J}\left(\delta, \widehat{\mathcal{G}}\right)
\;\lesssim\;
\delta \sqrt{B \log(1/\delta)}.$
By assumption, \(\widehat{\mathcal{G}}\) is a fixed, nonrandom function class
conditional on the training data, which is independent of the calibration
sample \(\mathcal{C}_n\). Hence, applying Lemma~\ref{lemma:local_rademacher_entropy}
conditional on the training data, we obtain
\begin{align*}
\mathcal{R}_n(\widehat{\mathcal{G}}, \delta)
&\;\lesssim\;
\frac{1}{\sqrt n}\,
\mathcal{J}(\delta,\widehat{\mathcal{G}})
\left(1+\frac{\mathcal{J}(\delta,\widehat{\mathcal{G}})}{\delta \sqrt n}\right) \\[0.4em]
&\;\lesssim\;
\frac{1}{\sqrt n}\,
\delta \sqrt{B \log(1/\delta)}
\left(
1 + 
\frac{\delta \sqrt{B \log(1/\delta)}}{\delta \sqrt n}
\right) \\[0.4em]
&\;\lesssim\;
\frac{\delta \sqrt{B \log(1/\delta)}}{\sqrt n}
\;+\;
\frac{B \log(1/\delta)}{n}.
\end{align*}
The critical radius
\[
\delta_n := \sqrt{\frac{B}{n}\,\log\Big(\frac{n}{B}\Big)},
\]
satisfies
\[
\delta_n
= \inf\{\delta > 0 : \mathcal{R}_n(\widehat{\mathcal{G}}, \delta)
\;\lesssim\; \delta^2\}.
\]

Applying Lemma~\ref{lemma:loc_max_ineq} conditional on the training data
with \(\mathcal{F} := \widehat{\mathcal{G}}\), we conclude that the following holds with probability at least
\(1 - e^{-u^2}\) for every \(f \in \widehat{\mathcal{G}}\):
\[
\frac{1}{n} \sum_{i=1}^n 
\big(f(O_i) - \mathbb{E}[f(O_i)]\big)
\;\lesssim\;
\delta_n^2
\;+\;
\delta_n\|f\|_{2,P_0}
\;+\;
\frac{u\|f\|_{2,P_0}}{\sqrt{n}}
\;+\;
\frac{u^2}{n}.
\]
Choosing \(u=\sqrt{\log(1/\eta)}\) gives \(1-e^{-u^2}=1-\eta\). Hence, with probability at least \(1-\eta\),
\begin{align*}
&\frac{1}{n} \sum_{i=1}^n 
\big(f(O_i) - \mathbb{E}[f(O_i)]\big)\\
&\lesssim
\delta_n^2
+ 
\delta_n\|f\|_{2,P_0}  +
\frac{\sqrt{\log(1/\eta)}\,\|f\|_{2,P_0}}{\sqrt{n}}
 + 
\frac{\log(1/\eta)}{n}.
\end{align*}





By boundedness of our nuisances,
\[
\big\|
\{\Gamma_0(V_K) - V_K\}
\{\widehat Y_{V_{K-1}} - V_K\}
\big\|
\;\lesssim\;
\|\Gamma_0(V_K) - V_K\|_{2,\rho}.
\]
Hence, with probability at least \(1 - \eta\),
\[
\begin{aligned}
(P_n - P_0)
&\Big\{
\Gamma_0(V_K) - V_K
\Big\}
\Big\{
\widehat Y_{V_{K-1}}
- V_K
\Big\} \\
&\lesssim\;
\delta_n^2
\;+\;
\delta_n \|\Gamma_0(V_K) - V_K\| \\[0.35em]
&\quad+\;
\frac{\sqrt{\log(1/\eta)}\,\|\Gamma_0(V_K) - V_K\|}{\sqrt{n}}
\;+\;
\frac{\log(1/\eta)}{n},
\end{aligned}
\]
where the implicit constants do not depend on \(B\).



Combining the above with \eqref{eqn::basicineq2}, we find with probability at least \(1 - \eta\),
\begin{equation}
\label{eqn::basicineq3}
\begin{aligned}
&\|\Gamma_0(V_K) - V_K\|^2\\
&\lesssim\;
\delta_n^2
\;+\;
\delta_n \|\Gamma_0(V_K) - V_K\|\\
&\quad+\;
\frac{\sqrt{\log(1/\eta)}}{\sqrt{n}}\,
\|\Gamma_0(V_K) - V_K\|
\;+\;
\frac{\log(1/\eta)}{n}
\\[0.4em]
&\quad+\;
\|\Gamma_0(V_K) - V_K\|\;
\big\|
\widehat{\mathcal T}_\pi V_K
- \widehat{\mathcal T}_\pi V_{K-1}
\big\|
\\[0.4em]
&\quad+\;
\|\Gamma_0(V_K) - V_K\|\;
\big\|
b_0\{(\widehat w_{\pi} - w_{\pi})\,
(\widehat q_{V_K} - q_{V_K})\}
\big\|_{2,\rho}.
\end{aligned}
\end{equation}
The inequality in \eqref{eqn::basicineq3} implies that, 
with probability at least $1-\eta$, the calibration error satisfies
\begin{align*}
\|\Gamma_0(V_K) - V_K\|
\;\lesssim\;&\;
\delta_n
\;+\;
\sqrt{\frac{\log(1/\eta)}{n}}
\\[0.4em]
&\;+\;
\big\|
\widehat{\mathcal T}_\pi V_K
- \widehat{\mathcal T}_\pi V_{K-1}
\big\|
\\[0.4em]
&\;+\;
\big\|
b_0\{(\widehat w_{\pi} - w_{\pi})\,
(\widehat q_{V_K} - q_{V_K})\}
\big\|_{2,\rho}.
\end{align*}
Recall that $\delta_n := \sqrt{\frac{B}{n}\,\log\Big(\frac{n}{B}\Big)}$. Then,  with probability at least $1-\eta$,
\begin{align*}
\|\Gamma_0(V_K) - V_K\|
\;\lesssim\;&\;
\sqrt{\frac{B}{n}\,\log\Big(\frac{n}{B}\Big)}
\;+\;
\sqrt{\frac{\log(1/\eta)}{n}}
\\[0.4em]
&\;+\;
\big\|
\widehat{\mathcal T}_\pi V_K
- \widehat{\mathcal T}_\pi V_{K-1}
\big\|
\\[0.4em]
&\;+\;
\big\|
b_0\{(\widehat w_{\pi} - w_{\pi})\,
(\widehat q_{V_K} - q_{V_K})\}
\big\|_{2,\rho}.
\end{align*}
\end{proof}


\section{Proof of Theorem \ref{theorem::regretHist}}

\subsection{Additional notation}
Define the bin–index map $\hat{b}(s)$ by setting $\hat{b}(s) = b$ whenever
$\hat{v}(s) \in I_b$.  Let
$\Pi_B$ denote the binning projection
\[
(\Pi_B g)(s)
:= E\!\left[g(S)\,\middle|\,\hat v(S)\in I_{\hat b(s)}\right].
\] Denote 
\[
\mathcal B_{\pi,B}(g)(s)
:= (\Pi_B \mathcal{T}_{\pi}g)(s).
\]
Define $V_B^\star$ as the fixed point satisfying
\[
\mathcal B_{\pi,B}(V_B^\star) = V_B^\star.
\]
Such a fixed point exists by Banach's fixed point theorem because 
$\mathcal B_{\pi,B}(g) \in \mathcal H_B(\hat v)$ whenever 
$g \in \mathcal H_B(\hat v)$, and the operator is a $\gamma$-contraction in the
sup-norm. 


Recall the projected
transition operator 
\[
(P_{\pi,B} h)(s)
=
\mathbb{E}_{\pi}\left[h(S') \,\middle|\, \hat{v}(S) \in I_{\hat{b}(s)} \right].
\]
It holds that
\[
P_{\pi,B}
\;:=\;
\Pi_B P_\pi,
\]
and
\[
\Pi_B\mathcal{T}_\pi(f)
-
\Pi_B\mathcal{T}_\pi(g)
=
\gamma\, P_{\pi,B}(f-g).
\]

\subsection{Supporting lemmas}
We begin with the deterministic ingredients for the histogram value bound.

Throughout, histogram bins mean the retained cells of the fitted partition:
empty empirical cells are merged with a neighboring cell or discarded, and
\(B\) denotes the number of retained cells. Population projections are defined
on the corresponding non-null cells; \(\rho\)-null cells are identified in
\(L^2(\rho)\) and omitted from the index set.




\begin{lemma}[Sup-norm contraction of the aggregated Bellman operator]
\label{lemma::agg-sup-contraction}
For any bounded \(f,g:\mathcal S\to\mathbb R\),
\[
\|\mathcal B_{\pi,B}f-\mathcal B_{\pi,B}g\|_\infty
\le
\gamma\|f-g\|_\infty .
\]
\end{lemma}

\begin{proof}
Since the rewards cancel on differences,
\[
\mathcal B_{\pi,B}f-\mathcal B_{\pi,B}g
=
\gamma P_{\pi,B}(f-g).
\]
The operator \(P_{\pi,B}\) is a Markov averaging
operator. Hence, for any bounded \(h\),
\[
|P_{\pi,B}h(s)|
\le
E_\pi[|h(S')|\mid \hat v(S)\in I_{\hat b(s)}]
\le
\|h\|_\infty .
\]
Taking \(h=f-g\) proves the claim.
\end{proof}

\begin{lemma}[Approximation error under aggregated concentrability]
\label{lemma::approxerror}
Assume \ref{cond::Ccoverage}. Then
\[
\|V_B^\star - v_0\|_{2,\rho}
\le
\frac{\sqrt{\mathcal C_{B,\gamma}}}{1-\gamma}\,
\|(I-\Pi_B)v_0\|_{2,\rho}.
\]
\end{lemma}

\begin{proof}
Because \(v_0=\mathcal T_\pi v_0\) and
\(V_B^\star=\Pi_B\mathcal T_\pi V_B^\star\),
\[
V_B^\star-v_0
=
\gamma P_{\pi,B}(V_B^\star-v_0)
+
(\Pi_B-I)v_0 .
\]
Let \(a:=(\Pi_B-I)v_0\). Iterating the display and using boundedness gives the
resolvent expansion
\[
V_B^\star-v_0
=
\sum_{t=0}^\infty \gamma^t P_{\pi,B}^t a .
\]
For any \(t\ge0\), set
\[
c_t
:=
\left\|
\frac{d(\rho P_{\pi,B}^t)}{d\rho}
\right\|_\infty .
\]
Jensen's inequality for the Markov operator \(P_{\pi,B}^t\) gives
\[
\|P_{\pi,B}^t a\|_{2,\rho}^2
\le
\int P_{\pi,B}^t(a^2)\,d\rho
=
\int a^2\,d(\rho P_{\pi,B}^t)
\le
c_t\|a\|_{2,\rho}^2 .
\]
Therefore,
\[
\|V_B^\star-v_0\|_{2,\rho}
\le
\sum_{t=0}^\infty \gamma^t\sqrt{c_t}\,\|a\|_{2,\rho}
=
\frac{\sqrt{\mathcal C_{B,\gamma}}}{1-\gamma}
\|(I-\Pi_B)v_0\|_{2,\rho}.
\]
\end{proof}


\begin{lemma}[Inexact iterations under aggregated concentrability]
\label{lemma::inexacterrors}
Assume \ref{cond::Ccoverage}. Let \(\{\eta_k\}_{k\ge1}\) be any sequence such
that
\[
\|V_k - \mathcal B_{\pi,B}(V_{k-1})\|_{2,\rho}
\le
\eta_k
\qquad\text{for all } k.
\]
For \(m\ge0\), set
\[
c_m
:=
\left\|
\frac{d(\rho P_{\pi,B}^m)}{d\rho}
\right\|_\infty .
\]
Then, for any \(K\ge1\),
\[
\|V_K - V_B^\star\|_{2,\rho}
\le
\gamma^K\|\hat v - V_B^\star\|_\infty
+
\sum_{j=1}^K
\gamma^{K-j}\sqrt{c_{K-j}}\,\eta_j .
\]
Consequently,
\[
\|V_K - V_B^\star\|_{2,\rho}
\le
\gamma^K\|\hat v - V_B^\star\|_\infty
+
\frac{\sqrt{\mathcal C_{B,\gamma}}}{1-\gamma}
\max_{1\le j\le K}\eta_j .
\]
\end{lemma}

\begin{proof}
Let
\[
e_k:=V_k-\mathcal B_{\pi,B}(V_{k-1}).
\]
Since \(V_B^\star=\mathcal B_{\pi,B}V_B^\star\),
\[
V_k-V_B^\star
=
\gamma P_{\pi,B}(V_{k-1}-V_B^\star)+e_k.
\]

Iterating this recursion yields
\[
V_K-V_B^\star
=
\gamma^K P_{\pi,B}^K(\hat v-V_B^\star)
+
\sum_{j=1}^K\gamma^{K-j}P_{\pi,B}^{K-j}e_j .
\]
The first term is bounded by
\(\gamma^K\|\hat v-V_B^\star\|_\infty\) because \(P_{\pi,B}\) is a Markov
operator. For the remaining terms, the same Jensen and change-of-measure
argument used in Lemma~\ref{lemma::approxerror} gives
\[
\|P_{\pi,B}^{K-j}e_j\|_{2,\rho}
\le
\sqrt{c_{K-j}}\,\|e_j\|_{2,\rho}
\le
\sqrt{c_{K-j}}\,\eta_j .
\]
The displayed bounds follow by the triangle inequality and the definition of
\(\mathcal C_{B,\gamma}\).
\end{proof}


\begin{lemma}[Error bound for inexact updates]
\label{lemma::errorperiter}
There exists \(C<\infty\) such that, with probability at least \(1-\delta\),
uniformly over \(k=1,\ldots,K\),
 \begin{align*}
&\|V_k -
  \mathcal B_{\pi,B}(V_{k-1})\|_{2,\rho}\\
 & \qquad \leq
C\left\{
\sqrt{\frac{B}{n}\,\log\Big(\frac{n}{B}\Big)}
\;+\;
\sqrt{\frac{\log(K/\delta)}{n}}
+
\big\|
b_0\{(\widehat w_{\pi} - w_{\pi})\,
(\widehat q_{V_{k-1}} - q_{V_{k-1}})\}
\big\|_{2,\rho}
\right\}
\end{align*}
 \end{lemma}
\begin{proof}
Equation~\ref{eqn::scores} in the proof of Theorem~\ref{theorem::Calerrorhist} shows that, for any transformation $f \in \mathcal G_B$ and any $k$,
\begin{equation*}
\frac{1}{n}\sum_{i=1}^n 
f(\hat{v}(S_i))
\big\{
\widehat Y_{V_{k-1}}(O_i)
- V_k(S_i)
\big\}
= 0.
\end{equation*}
Applying the above with 
$f \circ \hat{v} = V_B^{\star(k)} - V_k$, 
where 
\[
V_B^{\star(k)} := \mathcal B_{\pi,B}(V_{k-1}),
\]
we obtain
\begin{equation*}
\frac{1}{n}\sum_{i=1}^n 
\big(V_B^{\star(k)} - V_k\big)(S_i)
\Big\{
\widehat Y_{V_{k-1}}(O_i)
-
V_k(S_i)
\Big\}
= 0.
\end{equation*}
Adding and subtracting $P_0$, we find
\begin{equation}
\begin{aligned}
&P_0\left[
  (V_B^{\star(k)} - V_k)
  \Big\{
    \widehat Y_{V_{k-1}}
    - V_k
  \Big\}
\right]\\
&=
(P_0 - P_n)\left[
  (V_B^{\star(k)} - V_k)
  \Big\{
    \widehat Y_{V_{k-1}}
    - V_k
  \Big\}
\right].
\end{aligned}
\label{eq:pn_p0_identity}
\end{equation}

First, we study the left–hand side of
\eqref{eq:pn_p0_identity}. We have
\begin{align*}
    & P_0\left[
  (V_B^{\star(k)} - V_k)
  \Big\{
    \widehat Y_{V_{k-1}}
    - V_k
  \Big\}
\right]\\
& = \|V_B^{\star(k)} - V_k\|_{2,\rho}^2\\
& \quad + P_0\left[
  (V_B^{\star(k)} - V_k)
  \Big\{
    \widehat Y_{V_{k-1}}
    - V_B^{\star(k)}
  \Big\}
\right].
\end{align*}
Since \(V_B^{\star(k)}=\Pi_B\mathcal T_\pi V_{k-1}\) and
\(V_B^{\star(k)}-V_k\in\mathcal H_B(\hat v)\), the projection identity
gives
\[
P_0\left[
  (V_B^{\star(k)}-V_k)
  \{V_B^{\star(k)}-\mathcal T_\pi V_{k-1}\}
\right]=0.
\]
Thus
\begin{align*}
& P_0\left[
  (V_B^{\star(k)} - V_k)
  \Big\{
    \widehat Y_{V_{k-1}}
    - V_B^{\star(k)}
  \Big\}
\right] \\[0.4em]
&=
P_0\left[
  (V_B^{\star(k)} - V_k)
  \Big\{
    \widehat Y_{V_{k-1}}
    - \mathcal T_{\pi}V_{k-1}
  \Big\}
\right].
\end{align*}
By Theorem~\ref{theorem::DRpseudo} and arguing as in the proof of
Theorem~\ref{theorem::Calerrorhist}, it follows that
\begin{align*}
&\left|
P_0\left[
  (V_B^{\star(k)} - V_k)
  \Big\{
    \widehat Y_{V_{k-1}}
    - V_B^{\star(k)}
  \Big\}
\right]
\right|\\
&\;\le\;
\|V_B^{\star(k)} - V_k\|_{2,\rho}\;
\big\|
b_0\{(\widehat w_{\pi} - w_{\pi})\,
(\widehat q_{V_{k-1}} - q_{V_{k-1}})\}
\big\|_{2,\rho}.
\end{align*}
Putting it all together,
\begin{align*}
&\|V_B^{\star(k)} - V_k\|_{2,\rho}^2
 \leq
P_0\left[
  (V_B^{\star(k)} - V_k)
  \Big\{
    \widehat Y_{V_{k-1}}
    - V_k
  \Big\}
\right] \\
&+ 
\|V_B^{\star(k)} - V_k\|_{2,\rho}\;
\big\|
b_0\{(\widehat w_{\pi} - w_{\pi})\,
(\widehat q_{V_{k-1}} - q_{V_{k-1}})\}
\big\|_{2,\rho}.
\end{align*}
Hence, by \eqref{eq:pn_p0_identity},
\begin{equation}
\begin{aligned}
&\|V_B^{\star(k)} - V_k\|_{2,\rho}^2 \\
& \quad \leq
(P_0 - P_n)\left[
  (V_B^{\star(k)} - V_k)
  \Big\{
    \widehat Y_{V_{k-1}}
    - V_k
  \Big\}
\right]\\
& \quad + 
\|V_B^{\star(k)} - V_k\|_{2,\rho}\;
\big\|
b_0\{(\widehat w_{\pi} - w_{\pi})\,
(\widehat q_{V_{k-1}} - q_{V_{k-1}})\}
\big\|_{2,\rho}.
\end{aligned}
\label{eq:pn_p0_identity2}
\end{equation}


Next, we obtain a high–probability bound for the first term right–hand side of
\eqref{eq:pn_p0_identity2}.
Applying Lemmas~\ref{lemma:loc_max_ineq}, \ref{lemma:local_rademacher_entropy},
and \ref{lemma::entropynumbers}, and arguing as in the proof of
Theorem~\ref{theorem::Calerrorhist}, we obtain that, with probability at least
\(1 - \delta\),
\[
\begin{aligned}
\Big|(P_0 - P_n)
&\Big[
(V_B^{\star(k)} - V_k)
\big\{
\widehat Y_{V_{k-1}}
-
V_k
\big\}
\Big]\Big|
\\[0.4em]
&\lesssim\;
\delta_n^2
\;+\;
\delta_n\,\|V_B^{\star(k)} - V_k\|_{2,\rho}
\\[0.4em]
&\quad+\;
\frac{\sqrt{\log(1/\delta)}\,\|V_B^{\star(k)} - V_k\|_{2,\rho}}{\sqrt{n}}
\;+\;
\frac{\log(1/\delta)}{n},
\end{aligned}
\]
where \(\delta_n := \sqrt{(B/n)\,\log(n/B)}\), and the implicit constants do
not depend on \(B\).


Plugging this high probability bound into \eqref{eq:pn_p0_identity2},  we obtain that, with probability at least
\(1 - \delta\),
\begin{equation*}
\begin{aligned}
&\|V_B^{\star(k)} - V_k\|_{2,\rho}^2 \\
&\lesssim\;
\delta_n^2
\;+\;
\delta_n\,\|V_B^{\star(k)} - V_k\|_{2,\rho}
\\[0.4em]
&\quad+\;
\frac{\sqrt{\log(1/\delta)}\,\|V_B^{\star(k)} - V_k\|_{2,\rho}}{\sqrt{n}}
\;+\;
\frac{\log(1/\delta)}{n}\\
& \quad + 
\|V_B^{\star(k)} - V_k\|_{2,\rho}\;
\big\|
b_0\{(\widehat w_{\pi} - w_{\pi})\,
(\widehat q_{V_{k-1}} - q_{V_{k-1}})\}
\big\|_{2,\rho}.
\end{aligned}
\end{equation*}
Recalling that $\delta_n := \sqrt{\frac{B}{n}\,\log\Big(\frac{n}{B}\Big)}$, the inequality above implies that, 
with probability at least $1-\delta$,  
 \begin{align*}
\|V_B^{\star(k)} - V_k\|_{2,\rho}
\;\lesssim\;&\;
\sqrt{\frac{B}{n}\,\log\Big(\frac{n}{B}\Big)}
\;+\;
\sqrt{\frac{\log(1/\delta)}{n}}
\\[0.4em]
&  + 
\big\|
b_0\{(\widehat w_{\pi} - w_{\pi})\,
(\widehat q_{V_{k-1}} - q_{V_{k-1}})\}
\big\|_{2,\rho}.
\end{align*}
Taking the failure probability to be \(\delta/K\) and applying a union bound over
\(k=1,\ldots,K\) proves the result.
\end{proof}

\subsection{Proof of Theorem \ref{theorem::regretHist}}

\begin{proof}[Proof of Theorem \ref{theorem::regretHist}]
For \(C<\infty\) large enough, define
\begin{align*}
  \eta_k
&=
C \Bigg\{\sqrt{\frac{B}{n}\,\log\!\Big(\frac{n}{B}\Big)}
\;+\;
\sqrt{\frac{\log(K/\delta)}{n}}
\\
& \quad
\;+\;
\big\|
b_0\{(\widehat w_{\pi} - w_{\pi})\,
(\widehat q_{V_{k-1}} - q_{V_{k-1}})\}
\big\|_{2,\rho}
\Bigg\}
\end{align*}

By Lemma~\ref{lemma::errorperiter}, with probability at least \(1-\delta\),
uniformly over \(k=1,\ldots,K\),
\[
\|V_k-\mathcal B_{\pi,B}(V_{k-1})\|_{2,\rho}
\le
\eta_k .
\]
Lemma~\ref{lemma::inexacterrors} therefore gives
\begin{align*}
\|V_K - V_B^\star\|_{2,\rho}
\le\;&
\gamma^K\|\hat v - V_B^\star\|_\infty\\
&+
\frac{\sqrt{\mathcal C_{B,\gamma}}}{1-\gamma}
\max_{1\le j\le K}\eta_j .
\end{align*}
Combining this with Lemma~\ref{lemma::approxerror} and the triangle inequality
yields
\begin{align*}
\|V_K-v_0\|_{2,\rho}
\le\;&
\gamma^K\|\hat v - V_B^\star\|_\infty\\
&+
\frac{\sqrt{\mathcal C_{B,\gamma}}}{1-\gamma}
\left[
\|(I-\Pi_B)v_0\|_{2,\rho}
+
\max_{1\le j\le K}\eta_j
\right].
\end{align*}
Since \(\Pi_B\) is the \(L^2(\rho)\) projection onto
\(\mathcal H_B(\hat v)\),
\[
\|(I-\Pi_B)v_0\|_{2,\rho}
=
\min_{h\in\mathcal H_B(\hat v)}\|h-v_0\|_{2,\rho}.
\]
Substituting the definition of \(\eta_k\) and absorbing constants completes the
proof.
\end{proof}






\subsection{Bound on Successive Iteration Errors}
\label{appendix::convergence}
\begin{lemma}[Bound on Successive Iteration Errors]
Under the conditions of Theorem~\ref{theorem::regretHist},
\begin{align*}
&\|V_K - V_{K-1}\|_{2,\rho}\\
&\lesssim
\gamma^{K-1}\,
\|\hat v - V_B^\star\|_{\infty}
\\[0.8em]
&\quad+\;
\frac{\sqrt{\mathcal C_{B,\gamma}}}{1-\gamma}
C\Bigg(
\sqrt{\frac{B}{n}\,\log\!\Big(\frac{n}{B}\Big)}
\;+\;
\sqrt{\frac{\log(K/\delta)}{n}}
\\[0.8em]
&\qquad\qquad+
\max_{1\le j\le K}
\Big\|
b_0\{
(\widehat w_{\pi} - w_{\pi})\,
(\widehat q_{V_{j-1}} - q_{V_{j-1}})
\}
\big\|_{2,\rho}
\Bigg).
\end{align*}
\end{lemma}

\begin{proof}
Lemma~\ref{lemma::inexacterrors} and Lemma~\ref{lemma::errorperiter} give the
same bound for \(\|V_m-V_B^\star\|_{2,\rho}\), for \(m=K-1\) and \(m=K\), with
\(\gamma^m\|\hat v-V_B^\star\|_\infty\) as the finite-iteration term. By the
triangle inequality,
\begin{align*}
\|V_K - V_{K-1}\|_{2,\rho}
&\le\;
\|V_{K-1} - V_B^\star\|_{2,\rho}
\\[0.25em]
&\quad+\;
\|V_K - V_B^\star\|_{2,\rho}.
\end{align*}
Substituting the two bounds and absorbing constants proves the claim.
\end{proof}



\section{Proofs of Theorems \ref{theorem::Calerroriso} and \ref{theorem::regretiso}}


\begin{proof}[Proof of Theorem \ref{theorem::Calerroriso}]
The first-order optimality conditions for isotonic regression (equivalently, its
interpretation as a histogram estimator) imply that for any function
$f:\mathbb{R}\rightarrow\mathbb{R}$ that is a linear combination of the
indicator functions defining the isotonic partition,
\begin{equation*}
\frac{1}{n}\sum_{i=1}^n
f\!\big(V_K(S_i)\big)
\Big\{
\widehat Y_{V_{K-1}}(O_i)
-
V_K(S_i)
\Big\}
= 0.
\end{equation*}
See, for example, the proof of Lemma~C.1 in \citet{van2023causal} and
\citet{van2024stabilized}. Choosing $f$ appropriately yields
\begin{align*}
0
&=
\frac{1}{n}\sum_{i=1}^n
\Big\{
\Gamma_0(V_K)(S_i) - V_K(S_i)
\Big\} \\
&\qquad\qquad\times
\Big\{
\widehat Y_{V_{K-1}}(O_i)
- V_K(S_i)
\Big\},
\end{align*}
which is the same basic equality as \eqref{eqn::basicineq} in the proof of
Theorem~\ref{theorem::Calerrorhist}. The remainder of the argument proceeds
along the same lines with minor modifications.

Specifically, let $\mathcal{F}_{TV}$ denote the union of $\mathcal G_{\mathrm{iso}}$,
which is uniformly bounded by $2M$ under Condition~\ref{cond::C1}, with all
functions of bounded total variation bounded by the constant $C$ in
Condition~\ref{cond::C3}. By \citet{van1996weak}, this class satisfies the
uniform entropy integral bound $\mathcal{J}(\delta,\mathcal{F}_{TV})
\lesssim \sqrt{\delta}$.

Let
\[
m_K(t)
:=
E\!\left[(\mathcal T_\pi V_K)(S)\mid \hat v(S)=t,\mathcal C_n\right].
\]
Condition~\ref{cond::C3} states that \(m_K\) has bounded total variation. Since
\[
\Gamma_0(V_K)(S)=E\{m_K(\hat v(S))\mid V_K(S)\},
\]
and \(V_K\) is a nondecreasing function of \(\hat v\), Lemma 6 of
\citet{van2024stabilized} implies that the coarsened curve
\(\Gamma_0(V_K)\) also has bounded total variation and lies in
\(\mathcal{F}_{TV, \hat{v}} := \{f \circ \hat{v}: f \in
\mathcal{F}_{TV}\}\). Thus, it holds that
\[
\big(\Gamma_0(V_K) - V_K\big)
\big(\widehat Y_{V_{K-1}} - V_K\big)
\]
lies in a uniformly bounded subset of the class
\[
\widehat{\mathcal{G}}
:=
\Big\{
(f_1 - f_2)\big(\widehat Y_{f_2} - f_2\big)
:\;
f_1, f_2 \in \mathcal{F}_{TV, \hat{v}}
\Big\}.
\]
Arguing as in the proof of Theorem~\ref{theorem::Calerrorhist}, we have
\[
\mathcal{J}(\delta, \widehat{\mathcal{G}})
\lesssim
\mathcal{J}(\delta, \mathcal{F}_{TV})
\lesssim
\sqrt{\delta}.
\]
The proof now follows directly from Theorem~\ref{theorem::Calerrorhist} with this
new choice of $\widehat{\mathcal{G}}$ and its associated critical radius
$\delta_n = n^{-1/3}$ for monotone functions.
\end{proof}

 \begin{proof}[Proof of Theorem \ref{theorem::regretiso}]
The result follows directly from a union bound and the proofs of
Theorems~\ref{theorem::Calerrorhist} and \ref{theorem::regretHist}, with
$\mathcal G_B$ replaced by $\mathcal G_{B_n}$, the class of all
piecewise-constant functions with at most $B_n$ constant segments. In
particular, the entropy bound in Lemma~\ref{lemma::entropynumbers} continues to
apply to this class. Notably, the proofs of these theorems allow for
data-adaptive partitions and require only a deterministic upper bound on the
number of constant segments.
\end{proof}
 

\clearpage

\section*{NeurIPS Paper Checklist}

\begin{enumerate}[leftmargin=5mm]
\item \textbf{Claims}

Question: Do the main claims made in the abstract and introduction accurately
reflect the paper's contributions and scope?

Answer: \answerYes{}.

Justification: The abstract and introduction state the methodological claims,
including Bellman calibration, post-hoc value recalibration, finite-sample
guarantees, and synthetic empirical validation; see Sections~\ref{sec:experiments}
and the appendix for the corresponding theoretical and experimental support.

\item \textbf{Limitations}

Question: Does the paper discuss the limitations of the work performed by the
authors?

Answer: \answerYes{}.

Justification: The theory states the relevant assumptions, including
off-policy coverage, sample splitting, boundedness, and restrictions induced by
the one-dimensional calibration map. The conclusion and appendix also identify
distribution shift, extrapolation, and broader empirical validation as
limitations and future work.

\item \textbf{Theory assumptions and proofs}

Question: For each theoretical result, does the paper provide the full set of
assumptions and a complete (and correct) proof?

Answer: \answerYes{}.

Justification: The main results state their assumptions, and the appendix gives
proofs for the calibration-refinement results, doubly robust Bellman targets,
histogram and isotonic calibration guarantees, and supporting maximal
inequalities.

\item \textbf{Experimental result reproducibility}

Question: Does the paper fully disclose all the information needed to reproduce
the main experimental results of the paper to the extent that it affects the
main claims and/or conclusions of the paper (regardless of whether the code and
data are provided or not)?

Answer: \answerYes{}.

Justification: The experiments are synthetic. The paper specifies the data
generating processes, policies, rewards, sample sizes, cross-calibration
protocol, evaluation batches, base learners, calibrators, and metrics in
Section~\ref{sec:experiments} and Appendix~\ref{app:experiments}.

\item \textbf{Open access to data and code}

Question: Does the paper provide open access to the data and code, with
sufficient instructions to faithfully reproduce the main experimental results,
as described in supplemental material?

Answer: \answerYes{}.

Justification: The experiments do not rely on private data; all data are
synthetically generated from the paper's specified MDPs. An anonymous
supplementary code archive accompanies the submission with scripts,
configurations, seeds, and figure-generation code for the reported synthetic
experiments.

\item \textbf{Experimental setting/details}

Question: Does the paper specify all the training and test details (e.g., data
splits, hyperparameters, how they were chosen, type of optimizer) necessary to
understand the results?

Answer: \answerYes{}.

Justification: The appendix describes the offline data generation, train and
calibration splits, target and behavior policies, FQE/FVI learners, random
feature and neural specifications, iteration counts, weight clipping, and
evaluation protocol.

\item \textbf{Experiment statistical significance}

Question: Does the paper report error bars suitably and correctly defined or
other appropriate information about the statistical significance of the
experiments?

Answer: \answerYes{}.

Justification: The main results are based on independent replications and
independent evaluation data. Tables and figures report replication-level
variation where applicable, and the appendix states the number of replications
and evaluation batches used for the reported metrics.

\item \textbf{Experiments compute resources}

Question: For each experiment, does the paper provide sufficient information on
the computer resources (type of compute workers, memory, time of execution)
needed to reproduce the experiments?

Answer: \answerYes{}.

Justification: The experiments are synthetic offline-RL simulations using
linear, random-feature, boosted-tree, and modest neural value learners; they do
not require pretrained foundation models or large-scale distributed training.
The paper gives the sample sizes, model sizes, iterations, and replication
counts needed to assess computational scale, and the accompanying code archive
documents the executable configurations.

\item \textbf{Code of ethics}

Question: Does the research conducted in the paper conform, in every respect,
with the NeurIPS Code of Ethics
\url{https://neurips.cc/public/EthicsGuidelines}?

Answer: \answerYes{}.

Justification: The work is a methodological study of calibration for offline RL
value prediction using synthetic experiments. It does not involve human
subjects, private data, scraped content, or the release of high-risk models or
datasets.

\item \textbf{Broader impacts}

Question: Does the paper discuss both potential positive societal impacts and
negative societal impacts of the work performed?

Answer: \answerYes{}.

Justification: The paper motivates the method through value prediction and
off-policy evaluation, where better calibration can improve reliability. It also
notes that practical deployment remains limited by distribution shift and
coverage; in high-stakes applications, recovered or recalibrated values should
not be used without domain validation and appropriate safeguards.

\item \textbf{Safeguards}

Question: Does the paper describe safeguards that have been put in place for
responsible release of data or models that have a high risk for misuse (e.g.,
pre-trained language models, image generators, or scraped datasets)?

Answer: \answerNA{}.

Justification: The paper does not release pretrained language models, image
generators, scraped datasets, or other high-risk artifacts. The released
materials are synthetic experiment code and outputs for reproducing the paper's
methodological results.

\item \textbf{Licenses for existing assets}

Question: Are the creators or original owners of assets (e.g., code, data,
models), used in the paper, properly credited and are the license and terms of
use explicitly mentioned and properly respected?

Answer: \answerNA{}.

Justification: The experiments use synthetic data generated by the authors and
do not rely on existing datasets, pretrained models, or externally released
research assets. Prior methods and software ecosystems are cited where they are
used for methodological context.

\item \textbf{New assets}

Question: Are new assets introduced in the paper well documented and is the
documentation provided alongside the assets?

Answer: \answerYes{}.

Justification: The new assets are experiment scripts, synthetic-data
generators, configurations, seeds, and plotting code in the anonymous
supplementary archive. These assets are documented to support reproduction of
the reported synthetic experiments.

\item \textbf{Crowdsourcing and research with human subjects}

Question: For crowdsourcing experiments and research with human subjects, does
the paper include the full text of instructions given to participants and
screenshots, if applicable, as well as details about compensation (if any)?

Answer: \answerNA{}.

Justification: The paper does not involve crowdsourcing, user studies, or
research with human subjects.

\item \textbf{Institutional review board (IRB) approvals or equivalent for
research with human subjects}

Question: Does the paper describe potential risks incurred by study
participants, whether such risks were disclosed to the subjects, and whether
Institutional Review Board (IRB) approvals (or an equivalent approval/review
based on the requirements of your country or institution) were obtained?

Answer: \answerNA{}.

Justification: The paper does not involve crowdsourcing or human-subject
research, so IRB or equivalent approval is not applicable.

\item \textbf{Declaration of LLM usage}

Question: Does the paper describe the usage of LLMs if it is an important,
original, or non-standard component of the core methods in this research? Note
that if the LLM is used only for writing, editing, or formatting purposes and
does not impact the core methodology, scientific rigor, or originality of the
research, declaration is not required.

Answer: \answerNA{}.

Justification: LLMs are not used as part of the method, theory, experiments,
data generation, or evaluation in this research. Any ordinary writing, editing,
or formatting assistance would not affect the core methodology, scientific
rigor, or originality of the work.
\end{enumerate}

\end{document}
Off-policy fitting: a high-capacity value model trained on large
behavior-policy data can be informative but mis-scaled for the target-policy
state--action distribution \citep{fujimoto2019off}.
&
Use the large off-policy sample to learn the value predictor, then recalibrate its
scale on a smaller on-policy or target-policy reference sample.
\\
\midrule

Finite Bellman iteration: computational limits can force early stopping or a
shorter effective horizon, mis-scaling values for the intended return target
\citep{munos2008finite}.
&
Re-align predictions to the intended horizon without rerunning the full Bellman
iteration.
\\
\midrule

Approximation and optimization error: misspecified value classes, restricted
critics, or unstable saddle-point training can yield biased value predictions
\citep{uehara2021finite}.
&
Remap the biased but informative prediction to observed or Bellman-consistent values,
rather than requiring exact specification.
\\

\bottomrule
\end{tabularx}
\end{table}

\textbf{Contributions.}
We connect two largely separate literatures: post-hoc calibration for
supervised prediction and long-horizon value learning in RL. Existing work gives
rich methods for learning value functions, but little theory for calibrating the
outputs of an already learned value predictor. Our contributions are threefold.

\begin{enumerate}[leftmargin=1.5em]
\item We formalize \textbf{Bellman calibration}, a dynamic analogue of
supervised calibration for infinite-horizon value prediction. Bellman consistency
implies calibration, but calibration is weaker: predictions need only match
average Bellman backups among states assigned similar values. \textit{To diagnose calibration}, we propose
debiased estimators of $\ell^2$-Bellman calibration error.

\item We introduce \textbf{post-hoc Bellman calibration} via histogram binning
and isotonic regression. The methods apply to any value predictor and use
doubly robust Bellman targets for off-policy data.

\item We prove \textbf{finite-sample guarantees} for histogram calibration and
value error. The bounds require neither Bellman completeness nor realizability,
separate statistical, auxiliary-model, and iteration errors, and
relate value error to the best predictor based only on the original model output.
\end{enumerate}

\subsection{Related work}

\textbf{Post-hoc calibration.}
Our work builds on post-hoc calibration for supervised prediction
\citep{zadrozny2001obtaining,zadrozny2002transforming,niculescu2005predicting,
bella2010calibration,guo2017calibration}; see \citet{roth2022uncertain} for an
overview. This literature calibrates static predictions against observed labels,
empirical frequencies, or conditional means. We study long-horizon value
predictions in MDPs, where returns are not directly observed and must be inferred
through Bellman structure. We are also related to causal calibration for static
treatment-effect prediction \citep{van2023causal,whitehouse2024orthogonal}, but
our setting is dynamic and long-horizon, with calibration targets defined by
value functions rather than one-step conditional effects. The closest work is \citep{van2025automaticDRL}, who use isotonic post-hoc
calibration of \(Q\)-functions inside Double RL policy-value estimators
\citep{kallus2020double,kallus2022efficiently}. Their calibration step is used
as an auxiliary-model correction for a global value estimator. We instead study calibrated value prediction directly and give finite-sample
guarantees for both Bellman calibration error and value error. The same framework also applies to \(Q\)-function calibration by replacing
importance-weighted FVI with unweighted FQE and treating each state-action pair
as the state.
 

\textbf{Relation to Bellman learning and fine-tuning.}
Bellman calibration complements methods that change value learning itself,
including minimax or adversarial Bellman objectives
\citep{uehara2020minimax,uehara2023offline,xie2021batch,di2023pessimistic},
conservative updates \citep{kumar2020conservative,an2021uncertainty}, and
structural assumptions such as linearity, low rank, representations, or
discretization
\citep{shah2020sample,chang2022learning,peng1993convergence,van2006performance}.
It can be applied after any such method because it only needs scalar value
predictions and a reference sample. Unlike offline-to-online transfer,
critic/policy fine-tuning, or safe policy improvement
\citep{rusu2016progressive,nair2020awac,lee2022offline,laroche2019safe},
Bellman calibration fixes the target policy, value model, representation, and
Bellman learner, and fits only a one-dimensional post-hoc map to the target
value scale. Thus it is an output-scale correction rather than RL fine-tuning,
with finite-sample guarantees for calibration and value error. Related work on
calibrated dynamics uncertainty, conformal sets, and distributional RL targets
different objects
\citep{malik2019calibrated,zhang2023conformal,sun2023conformal,
bellemare2017distributional,dabney2018implicit}.

\section{Calibration for Value Functions}
\subsection{Setup and Notation}

We consider an MDP with state space \(\mathcal S\), discrete action space
\(\mathcal A\), initial distribution \(\rho\), transition kernel
\(P(s'\mid s,a)\), reward function \(r_0(s,a)\), and discount factor
\(\gamma\in[0,1)\). Data are collected under a behavior policy
\(b_0(a\mid s)\). Each calibration sample is an i.i.d. transition
\((S,A,R,S')\), where \((S,A)\sim \rho\times b_0\),
\(S'\sim P(\cdot\mid S,A)\), and
\(R=r_0(S,A)+\varepsilon\) with \(E[\varepsilon\mid S,A]=0\). We observe
\(\mathcal C_n:=\{(S_i,A_i,R_i,S_i')\}_{i=1}^n\).

Let \(\pi(a\mid s)\) be the target policy. The value function under \(\pi\) is
\[
v_0(s)
=
E_{\pi}\left[\sum_{t=0}^\infty \gamma^t R_t
\,\middle|\, S_0=s\right],
\]
where \(E_\pi\) denotes the law induced by
\(A_t\sim\pi(\cdot\mid S_t)\) and
\(S_{t+1}\sim P(\cdot\mid S_t,A_t)\).

For a state-action function \(f\) and state function \(v\), define
\[
(\pi f)(s):=\sum_{a\in\mathcal A}\pi(a\mid s)f(s,a),
\qquad
(Pv)(s,a):=E[v(S')\mid S=s,A=a].
\]
Write \(w_\pi:=\pi/b_0\), \(r_{0,\pi}:=\pi r_0\), \(P_\pi:=\pi P\), and
\[
\mathcal T_\pi v := r_{0,\pi}+\gamma P_\pi v
\]
for the policy Bellman operator. The value function is the unique bounded fixed
point \citep{bellman1952theory,bellman1966dynamic}
\begin{equation}
\label{eqn::bellman}
v_0=\mathcal T_\pi v_0 .
\end{equation}
A measure \(\mu\) is stationary for a Markov operator \(P\) if \(\mu P=\mu\).

Let \(L^2(\rho)\) denote square-integrable state functions with norm
\(\|f\|_{2,\rho}:=\{E_\rho[f(S)^2]\}^{1/2}\). For a measure \(\mu\) on
\(\mathcal S\), write
\[
\|f\|_{2,\mu}
:=
\left\{\int f(s)^2\,\mu(ds)\right\}^{1/2}.
\]
Let \(\|f\|_{n,S,A}\) and \(\|g\|_{n,S'}\) denote empirical \(L^2\) norms over
\(\{(S_i,A_i)\}_{i=1}^n\) and \(\{S_i'\}_{i=1}^n\), respectively. We write
\(\lesssim\) for inequalities up to absolute constants and
\([B]:=\{1,\ldots,B\}\).

 

\subsection{Bellman Calibration: A Minimal Requirement}

Let \(\hat v\) be an estimated value function for the target policy \(\pi\),
fit on data independent of the calibration sample \(\mathcal C_n\). Our goal is
not to replace the underlying value-learning procedure, but to impose a weak
post-hoc Bellman consistency requirement on its output. Throughout, we call the
scalar \(\hat v(S)\) the predicted value; in the supervised calibration
literature, this same scalar is often called a prediction score.

A natural starting point is \emph{statewise Bellman consistency}:
\[
\hat v(S)
\approx
E_\pi\!\left[R+\gamma \hat v(S')\mid S\right].
\]
Equivalently, the Bellman residual
\(R+\gamma \hat v(S')-\hat v(S)\) should have conditional mean zero given the
full state. In high-dimensional or continuous state spaces, this is a strong
requirement: checking or enforcing it requires estimating a state-dependent
conditional mean and is closely tied to structural assumptions such as Bellman
completeness \citep{chen2019information}.

We instead require Bellman consistency only after grouping states by their
predicted value:
\[
\hat v(S)
\approx
E_\pi\!\left[R+\gamma \hat v(S')\mid \hat v(S)\right].
\]
Thus, predictions should agree on average with their Bellman targets within
groups of states assigned the same predicted value \(\hat v(S)\). This is the
dynamic analogue of post-hoc calibration in supervised prediction, where
predictions are required to be correct on average conditional on their own
numerical output
\citep{zadrozny2002transforming,niculescu2005predicting,guo2017calibration}.
The benefit is that a high-dimensional Bellman moment is reduced to a
one-dimensional conditional moment, enabling post-hoc correction without
requiring the full fitted Bellman update to be correctly specified.

\textbf{Bellman calibration.}
For any candidate value function \(v\), define the Bellman-calibration map
\begin{equation}
\label{eqn::calibrationmap}
\Gamma_0(v)(s)
:=
E_\pi\!\left[R+\gamma v(S') \mid v(S)=v(s)\right].
\end{equation}
We say that \(v\) is \emph{Bellman calibrated} if
\[
v(S)=\Gamma_0(v)(S)\quad\text{almost surely},
\]
or equivalently,
\[
E_\pi\!\left[R+\gamma v(S')-v(S)\mid v(S)\right]=0
\quad\text{almost surely}.
\]
 

\textbf{A necessary but weak condition.}
The true value function \(v_0\) is Bellman calibrated. The Bellman equation gives
\(v_0(S)=E_\pi[R+\gamma v_0(S')\mid S]\); conditioning on the coarser predicted value
\(v_0(S)\) yields
\[
v_0(S)
=
E_\pi[R+\gamma v_0(S')\mid v_0(S)]
=
\Gamma_0(v_0)(S).
\]
Thus Bellman consistency implies Bellman calibration. The converse need not
hold: because calibration only conditions on the scalar predicted value \(v(S)\),
state-level errors can average out within predicted-value groups
\citep{gupta2020distribution,van2023causal}. Hence Bellman calibration is a
minimal reliability requirement, not a guarantee of value accuracy.

\textbf{Reward-model interpretation.}
Define the reward implied by a value predictor \(v\) as
\[
r_{\pi,v}(s):=v(s)-\gamma P_\pi v(s).
\]
This is the one-step reward under which \(v\) would satisfy the Bellman
equation for policy \(\pi\). Bellman calibration is equivalent to
\[
E_\pi\!\left[R-r_{\pi,v}(S)\mid v(S)\right]=0.
\]
Thus, among states assigned the same predicted long-run value, the reward
implied by \(v\) matches the true one-step reward on average. When
\(\gamma=0\), this reduces to ordinary regression calibration
\citep{roth2022uncertain}.


 

\subsection{Bellman Calibration and Value Error}

We next relate Bellman calibration to value prediction error. Because
calibration is based only on the predicted value, it cannot recover information
lost by compressing the state \(S\) to the scalar predicted value
\(\hat v(S)\). It can, however, remove systematic Bellman miscalibration among
transformations of this predicted value.

Let \(\mathcal G:=\{g:\mathbb R\to\mathbb R\}\) denote the unrestricted class
of scalar calibration maps. The value predictors generated by the original
prediction \(\hat v(S)\) form the class
\[
\mathcal V_{\hat v}
:=
\{g\circ\hat v:g\in\mathcal G\},
\qquad
\Pi_{\hat v}
:=
\text{\(L^2(\rho)\)-projection onto }\mathcal V_{\hat v} .
\]
Thus, \(\Pi_{\hat v}f\) is the best predictor of \(f\) that
depends on \(S\) only through \(\hat v(S)\). Define the projected transition
and Bellman operators by
\[
P_{\pi,\hat v}
:=
\Pi_{\hat v}P_\pi,
\qquad
\mathcal B_{\pi,\hat v}q
:=
\Pi_{\hat v}\mathcal T_\pi q,
\qquad
P_{\pi,\hat v}f(s)
=
E_\pi[f(S')\mid \hat v(S)=\hat v(s)].
\]
Let \(V_{\hat v}^\star\in\mathcal V_{\hat v}\) be the fixed point
\begin{equation}
\label{eqn::oraclefixed}
V_{\hat v}^\star
=
\mathcal B_{\pi,\hat v}V_{\hat v}^\star .
\end{equation}
This is the population Bellman-calibrated predictor within
\(\mathcal V_{\hat v}\).

\begin{enumerate}[label=\textbf{A\arabic*}, ref={A\arabic*}, leftmargin=1.5em, series=cond]
\item \label{cond::A1}
\(\rho\) is stationary for \(P_{\pi,\hat v}\).
\end{enumerate}

Under \ref{cond::A1}, \(\mathcal B_{\pi,\hat v}\) is a
\(\gamma\)-contraction on \(L^2(\rho)\)
(Appendix~\ref{appendix::stationary};
\citep{patterson2022generalized,vdLaanKallus2025FQE}).

\begin{theorem}[Calibration--refinement bound]
\label{theorem::calrefine}
Under \ref{cond::A1},
\[
\|\hat v-v_0\|_{2,\rho}
\;\le\;
\frac{1}{1-\gamma}
\|\Pi_{\hat v}v_0-v_0\|_{2,\rho}
+
\|\hat v-V_{\hat v}^\star\|_{2,\rho} .
\]
\end{theorem}

Theorem~\ref{theorem::calrefine} separates value error into prediction-only
approximation error and miscalibration. The prediction-only approximation error
\(\|\Pi_{\hat v}v_0-v_0\|_{2,\rho}\)
measures information lost by reducing the state to the predicted value \(\hat v(S)\).
The miscalibration term
\(\|\hat v-V_{\hat v}^\star\|_{2,\rho}\) measures how far
the predictor is from the population Bellman-calibrated predictor using the same original prediction. If
\(\hat v\) is Bellman calibrated within \(\mathcal V_{\hat v}\), then
\(\hat v=V_{\hat v}^\star\), so the second term vanishes. Thus, calibration
cannot make an uninformative prediction informative, but it can remove systematic
Bellman miscalibration without changing the underlying representation. This mirrors
classical calibration--refinement decompositions in
classification and regression
\citep{murphy1973new,degroot1983comparison,van2023causal,whitehouse2024orthogonal}.


If \(\hat v\) is Bellman calibrated, then \(\hat v=V_{\hat v}^\star\), so the
second term vanishes and only the prediction-only approximation error remains, up to the factor
\((1-\gamma)^{-1}\):
\[
\|\hat v-v_0\|_{2,\rho}
\;\le\;
\frac{1}{1-\gamma}
\min_{\theta:\mathbb R\to\mathbb R}
\|\theta\circ \hat v-v_0\|_{2,\rho}
+
\|\hat v-V_{\hat v}^\star\|_{2,\rho}.
\]
Hence, post-hoc calibration competes with the best value predictor expressible
as a function of the original prediction.

\subsection{Evaluating Bellman Calibration Error with Doubly Robust Bellman Targets}
\label{sec:bellman-calibration-error}

We use Bellman calibration error both as a population reliability criterion and
as an empirical diagnostic for comparing value predictors. For a candidate
value function \(v\), define
\begin{equation}
\label{eqn::calerror}
\mathrm{Cal}_{\ell^2}^2(v)
:=
\|v-\Gamma_0(v)\|_{2,\rho}^2
=
E_\rho\!\left[
  \left\{
    E_\pi[R+\gamma v(S')-v(S)\mid v(S)]
  \right\}^2
\right].
\end{equation}
This quantity is zero exactly when \(v\) is Bellman calibrated. It measures
systematic Bellman residuals that remain after grouping states by the predicted
value \(v(S)\), and therefore serves as a scalar diagnostic of numerical
reliability.

Because data are generated under the behavior policy \(b_0\), we estimate the
Bellman target conditional on the predicted value with a doubly robust Bellman target estimate, extending
the calibration target estimates of \citep{xu2022calibration,van2023causal} to
dynamic settings. Let \(\widehat w_\pi\) estimate \(w_\pi:=\pi/b_0\). For a
fixed \(v\), define
\[
q_v(s,a)
:=
E\{R+\gamma v(S')\mid S=s,A=a\},
\qquad
T_\pi v(t)
:=
E_\pi\{R+\gamma v(S')\mid v(S)=t\}.
\]
Let \(\hat q_v\) estimate \(q_v\). For \(O=(S,A,R,S')\), set
\[
\widehat Y_v(O)
:=
(\pi \hat q_v)(S)
+
\widehat w_\pi(A\mid S)
\{R+\gamma v(S')-\hat q_v(S,A)\},
\]
and write
\[
\widehat T_\pi v(t)
:=
E\{\widehat Y_v(O)\mid v(S)=t\}.
\]
This construction includes importance weighting \((\hat q_v=0)\) and
model-based Bellman-target evaluation \((\widehat w_\pi=0)\) as special cases.
If either \(w_\pi\) or \(q_v\) is correctly estimated, then
\(\widehat T_\pi v(t)=T_\pi v(t)\). Thus \(\widehat Y_v-v(S)\) targets the
Bellman residual conditional on the predicted value in \(\mathrm{Cal}_{\ell^2}(v)\).

\begin{theorem}[Doubly robust Bellman target conditional on the predicted value]
\label{theorem::DRpseudo}
For any fixed \(v\),
\[
\widehat T_\pi v(t)-T_\pi v(t)
=
E\!\left[
  b_0\!\left\{(w_\pi-\widehat w_\pi)(\hat q_v-q_v)\right\}(S)
  \mid v(S)=t
\right].
\]
Consequently, \(\widehat Y_v\) is conditionally unbiased for
\(T_\pi v\{v(S)\}\) given \(v(S)\) if either
\(\widehat w_\pi=w_\pi\) or \(\hat q_v=q_v\).
\end{theorem}

\paragraph{Debiased calibration-error estimation.}
Let \(U_i:=v(S_i)\) and \(\widehat Y_i:=\widehat Y_v(O_i)\). Following
\citet{xu2022calibration}, estimate the squared calibration error by the cross-fitted debiased moment
\[
\widehat{\mathrm{Cal}}_{\ell^2,\mathrm{db}}^2(v)
:=
\frac1n\sum_{i=1}^n
\{\widehat Y_i-U_i\}
\{\widehat T_{\pi,-i}v(U_i)-U_i\},
\]
where \(\widehat T_{\pi,-i}v\) is fit without observation \(i\). The population
identity behind this estimator is
\[
E\!\left[
\{\widehat Y_v-v(S)\}
\left\{E(\widehat Y_v\mid v(S))-v(S)\right\}
\right]
=
E\!\left[
\left\{E(\widehat Y_v\mid v(S))-v(S)\right\}^2
\right].
\]
Under the predicted-value conditional unbiasedness established above, this equals
\(\mathrm{Cal}_{\ell^2}^2(v)\), up to error from estimating
\(\widehat T_\pi v\) and the auxiliary estimates in \(\widehat Y_v\).

\section{Algorithms and Theory}

\subsection{Iterated Bellman Calibration}

Algorithm~\ref{alg::cal-gen} gives a post-hoc procedure for transforming an
initial value predictor \(\hat v\) into a calibrated predictor
\(V_{\mathrm{cal}}=V_K=g_K\circ \hat v\). The key idea is to run fitted
Bellman iteration using the one-dimensional input \(\hat v(S)\). Starting from
\(V_0=\hat v\), each iteration forms doubly robust Bellman target estimates
\(\widehat Y_{V_k}(O_i)\) and regresses them on the fixed inputs
\(\hat v(S_i)\) using a calibrator class \(\mathcal G\), yielding
\(V_{k+1}=g_{k+1}\circ \hat v\).

For a scalar map class \(\mathcal G\), write
\(\mathcal V_{\hat v}(\mathcal G):=\{g\circ\hat v:g\in\mathcal G\}\). The
algorithm is an FVI-style update restricted to this class. This restriction avoids
high-dimensional function approximation while correcting Bellman
miscalibration in the original predictor. Under the contraction results below, this projected iteration converges
geometrically, so \(O\{\log(1/\varepsilon)/(1-\gamma)\}\) iterations suffice to
reach optimization error \(\varepsilon\).


\begin{algorithm}[h]
\caption{\textbf{Iterated Bellman Calibration}}
\label{alg::cal-gen}
\begin{algorithmic}[1]
\BCINPUT \(\hat v\), \(\mathcal C_n=\{O_i\}_{i=1}^n\), Bellman target-estimate map \(v\mapsto\widehat Y_v\), class \(\mathcal G\), iterations \(K\)
\STATE Initialize \(V_0=\hat v\)
\FOR{\(k=0,\ldots,K-1\)}
  \STATE \(V_{k+1}=g_{k+1}\circ\hat v \quad \text{where} \quad g_{k+1}\in\arg\min_{g\in\mathcal G}\sum_{i=1}^n
  \{\widehat Y_{V_k}(O_i)-g(\hat v(S_i))\}^2\)
\ENDFOR
\BCOUTPUT \(V_K\)
\end{algorithmic}
\end{algorithm}

Heuristically, \(V_K\) is calibrated because the iteration targets the fixed
point \(V_{\hat v}^\star\) of the Bellman operator projected onto the original
predictions, defined in \eqref{eqn::oraclefixed}. Exact finite-sample recovery of this fixed point would
generally require realizability in the class
\(\mathcal V_{\hat v}(\mathcal G)\) and Bellman completeness of the
fitted update. Our goal is weaker but more robust: to guarantee reliability
through Bellman calibration without requiring the full fitted Bellman iteration
to be correctly specified.

We therefore instantiate Algorithm~\ref{alg::cal-gen} with one-dimensional
calibrators based on histogram regression \citep{stone1977consistent} and
isotonic regression \citep{barlow1972isotonic}. These are dynamic analogues of
histogram binning
\citep{zadrozny2001obtaining,zadrozny2002transforming,gupta2021distribution}
and isotonic calibration
\citep{zadrozny2002transforming,niculescu2005predicting,van2023causal,
van2025generalized}, and they admit finite-sample Bellman calibration
guarantees without Bellman completeness.

 

\subsection{Histogram Calibration}
\label{sec:hist}

Histogram calibration instantiates Algorithm~\ref{alg::cal-gen} with a
piecewise-constant calibrator. Partition the initial values
\(\{\hat v(S_i)\}_{i=1}^n\) into ordered bins \(I_1,\ldots,I_B\), using either
empirical quantiles (\emph{equal-mass}) or uniform discretization
(\emph{equal-width}). We take
\[
\mathcal G_B
:=
\mathrm{span}\{\mathbf 1_{I_b}:b\in[B]\},
\qquad
\mathcal H_B(\hat v)
:=
\{g\circ\hat v:g\in\mathcal G_B\},
\]
so every \(h\in\mathcal H_B(\hat v)\) is a value predictor that is constant
within bins of the original prediction. The bins may be data-adaptive, but
\(B\) is deterministic and may grow with \(n\).

With this class, each iteration simply averages the fitted Bellman targets
within each bin:
\[
g_{k+1}(t)
=
\frac{1}{n_b}
\sum_{i:\hat v(S_i)\in I_b}
\widehat Y_i^{(k)},
\qquad
t\in I_b,
\]
where \(n_b:=\sum_{i=1}^n\mathbf 1\{\hat v(S_i)\in I_b\}\). Thus, histogram
calibration is fitted Bellman iteration on a finite aggregation of the initial
predicted value. At convergence, the output approximately satisfies the empirical aggregated
Bellman equation within each bin:
\[
V_K(s)
\approx
E_n\!\left[
\widehat Y_{V_K}(O)
\,\middle|\,
\hat v(S)\in I_b
\right],
\qquad \hat v(s)\in I_b .
\]
This is the empirical analogue of Bellman calibration and is the key property
used in the next result.

 \textbf{Calibration error.}
We next give finite-sample, finite-iteration bounds for the Bellman calibration
error \eqref{eqn::calerror}. The analysis uses the following regularity
conditions.

\begin{enumerate}[label=\textbf{C\arabic*}, ref={C\arabic*}, leftmargin=1.5em, series=cond]

\item \label{cond::C1}
\textbf{Boundedness.}
\(R\), \(w_\pi\), \(\widehat w_\pi\), \(\hat v\), the Bellman targets \(q_f\),
and their estimates \(\hat q_f\) are uniformly bounded by \(M<\infty\) for all
\(f\) considered.

\item \label{cond::C1b}
\textbf{Sample splitting.}
The auxiliary estimators \(\widehat w_\pi\) and \(v\mapsto\hat q_v\), and the
initial predictor \(\hat v\) are trained on data independent of
\(\mathcal C_n\).

\item \label{cond::C2}
\textbf{Stability of the Bellman-target regression.}
For all \(f,g\in\mathcal H_B(\hat v)\) with \(\|f\|_\infty,\|g\|_\infty\le M\),
\[
\|\hat q_f-\hat q_g\|_{n,S,A}
\le
L\,\|f-g\|_{n,S'} .
\]
\end{enumerate}

These conditions are mainly technical. Boundedness can be enforced by clipping
or truncation, and sample splitting can be replaced by cross-fitting. Condition
\ref{cond::C2} controls the complexity introduced by the map
\(v\mapsto\hat q_v\); it holds trivially for the importance-weighted target
\((\hat q_v=0)\), and holds with \(L=1\) for cellwise estimates of
\(E\{R+\gamma f(S')\mid S,A\}\).

\begin{theorem}[Calibration error for histogram binning]
\label{theorem::Calerrorhist}
Assume \ref{cond::C1}--\ref{cond::C2}. Then there is a constant
\(C<\infty\) such that, for any \(K\in\mathbb N\), with probability at least
\(1-\delta\),
\begin{align*}
\mathrm{Cal}_{\ell^2}(V_K)
\le\;&
C\left\{
\sqrt{\frac{B}{n}\log\!\left(\frac{n}{B}\right)}
+
\sqrt{\frac{\log(1/\delta)}{n}}
\right\} \\
&+
\big\|b_0\{(\widehat w_\pi-w_\pi)
(\widehat q_{V_K}-q_{V_K})\}\big\|_{2,\rho}  +
\big\|
\widehat T_\pi V_K
-
\widehat T_\pi V_{K-1}
\big\|_{2,\rho}.
\end{align*}
\end{theorem}

The bound separates three errors. The first line is the statistical error of
one-dimensional histogram calibration, matching classical rates when the
calibration target is observed directly
\citep{gupta2020distribution,gupta2021distribution,van2025generalized}. The
second line is the doubly robust auxiliary-model error from
Theorem~\ref{theorem::DRpseudo}. The third line is the finite-iteration error,
which decays geometrically as $\gamma^K$ under the contraction conditions in
Appendix~\ref{appendix::convergence}, as in fitted value iteration
\citep{munos2008finite}.

The leading calibration rate does not accumulate over the horizon: \(\gamma\)
enters only through the finite-iteration and auxiliary-model terms. This reflects that
Bellman calibration is a one-step conditional prediction problem, even though
the value function itself is long-horizon. However, Theorem~\ref{theorem::calrefine}
shows that translating calibration into value accuracy still depends on
discounting and coverage.

\textbf{Value error.}
We next show that post-hoc calibration does not add error beyond the
information loss from restricting to the histogram calibrator class, while it can
remove Bellman miscalibration in the original predictor.

Let \(\Pi_B\) be the \(L^2(\rho)\) projection onto \(\mathcal H_B(\hat v)\):
\[
\Pi_B q
:=
\arg\min_{h\in\mathcal H_B(\hat v)}
\|q-h\|_{2,\rho}.
\]
Define the aggregated Bellman operator and fixed point
\[
\mathcal B_{\pi,B}q
:=
\Pi_B\mathcal T_\pi q,
\qquad
V_B^\star
=
\mathcal B_{\pi,B}V_B^\star,
\]
and write \(P_{\pi,B}:=\Pi_B P_\pi\). Equivalently, for
\(\hat b(s)=b\) when \(\hat v(s)\in I_b\),
\[
P_{\pi,B}f(s)
=
E_\pi[f(S')\mid \hat v(S)\in I_{\hat b(s)}].
\]

Define
\[
 
c_{B,t}:=
\left\|
\frac{d(\rho P_{\pi,B}^t)}{d\rho}
\right\|_\infty,
\qquad
\mathcal C_{B,\gamma}
:=
\left\{(1-\gamma)\sum_{t=0}^\infty
\gamma^t\sqrt{c_{B,t}}\right\}^2 .
\]

\begin{enumerate}[label=\textbf{C\arabic*}, ref={C\arabic*}, leftmargin=1.5em, resume=cond]
\item \label{cond::Ccoverage}
\textbf{Finite aggregated concentrability coefficient.}
\(\mathcal C_{B,\gamma} < \infty\).
\end{enumerate}
 

\begin{theorem}[Value error for histogram calibration]
\label{theorem::regretHist}
Assume \ref{cond::C1}--\ref{cond::Ccoverage}. Then there exists a constant
\(C<\infty\) such that, for any \(K\in\mathbb N\), with probability at least
\(1-\delta\),
{\footnotesize
\begin{align*}
\|V_K-v_0\|_{2,\rho}
\le\;&
\gamma^K\|\hat v-V_B^\star\|_{\infty}
+
\frac{\sqrt{\mathcal C_{B,\gamma}}}{1-\gamma}
\min_{h\in\mathcal H_B(\hat v)}\|h-v_0\|_{2,\rho}
\\
&+
C\frac{\sqrt{\mathcal C_{B,\gamma}}}{1-\gamma}
\Bigg[
  \sqrt{\frac{B}{n}\log\!\left(\frac{n}{B}\right)}
  +
  \sqrt{\frac{\log(K/\delta)}{n}}
\\
&\qquad\qquad\qquad\qquad
+
\max_{0\le j\le K}
\big\|
(\widehat w_\pi-w_\pi)
(\widehat q_{V_j}-q_{V_j})
\big\|_{2,\rho}
\Bigg].
\end{align*}
}
\end{theorem}
\textbf{Proof idea.}
The result avoids realizability and Bellman completeness for the original value
class because completeness holds after aggregation. For any
\(f\in\mathcal H_B(\hat v)\), the aggregated transition \(P_{\pi,B}f\) is
binwise constant in \(\hat v\), so \(P_{\pi,B}f\in\mathcal H_B(\hat v)\)
and hence \(\mathcal B_{\pi,B}f\in\mathcal H_B(\hat v)\). The proof uses
sup-norm contraction to obtain the aggregated fixed point, then propagates
\(\rho\)-norm one-step regression errors through a discounted resolvent
controlled by \(\mathcal C_{B,\gamma}\). It also bounds
\(V_B^\star-v_0\) by a histogram analogue of
Theorem~\ref{theorem::calrefine}. Unlike prior discretization arguments for
Bellman consistency
\citep{whitt1978approximations,peng1993convergence,van2006performance,
chen2019information,xie2021batch}, we discretize \(\hat v(S)\) to
target Bellman calibration.
 
\textbf{Discussion.}
Theorems~\ref{theorem::Calerrorhist} and~\ref{theorem::regretHist} separate
calibration from value accuracy. With known importance weights and negligible
iteration error, histogram calibration drives Bellman calibration error to zero
at the one-dimensional rate \(\sqrt{(B/n)\log(n/B)}\). Value error contains the
same stochastic term, but also includes the prediction-only approximation error
\[
\frac{\sqrt{\mathcal C_{B,\gamma}}}{1-\gamma}
\min_{h\in\mathcal H_B(\hat v)}\|h-v_0\|_{2,\rho},
\]
the cost of restricting the calibrated predictor to binned transformations of
the original prediction. Thus \(B\) controls the usual bias--variance tradeoff:
coarse bins may discard information in \(\hat v\), while fine bins increase
variance.

It suffices to choose \(B\) by predictive cross-validation. The prediction-only approximation term
satisfies
\[
\min_{h\in\mathcal H_B(\hat v)}\|h-v_0\|_{2,\rho}
\le
\|\hat v-v_0\|_{2,\rho}
+
\min_{h\in\mathcal H_B(\hat v)}\|h-\hat v\|_{2,\rho}.
\]
Thus binning is no worse than the original predictor, up to identity
approximation on the range of \(\hat v(S)\). For equal-mass bins, the usual histogram
bias--variance balance suggests \(B\asymp n^{1/3}\), yielding
\citep{gyorfi2002distribution}
\[
\|V_K-v_0\|_{2,\rho}
\le
\frac{\sqrt{\mathcal C_{B,\gamma}}}{1-\gamma}
\|\hat v-v_0\|_{2,\rho}
+
O_p\{(\log n/n)^{1/3}\},
\]
with Bellman calibration error vanishing at the same rate.


Finally, \(\mathcal C_{B,\gamma}\) captures coverage: it
measures how well the \(\rho\)-norm controls the discounted future
distributions of the aggregated target-policy chain. This is the histogram
analogue of the concentrability coefficients in fitted value iteration
\citep{munos2008finite,foster2021offline}.


\subsection{Isotonic Calibration}

Histogram calibration requires choosing the number of bins \(B\). As a
low-tuning alternative, we consider \emph{isotonic iterated Bellman
calibration}, obtained by taking
\(\mathcal G=\mathcal G_{\mathrm{iso}}\) in
Algorithm~\ref{alg::cal-gen}, where \(\mathcal G_{\mathrm{iso}}\) is the class
of nondecreasing functions on \(\mathbb R\). Each update isotonic-regresses the
fitted Bellman targets on the initial values of \(\hat v(S_i)\), yielding a
monotone calibrator. This can be computed in near-linear time by the
pool-adjacent-violators algorithm \citep{best1990active}.

Like histogram calibration, isotonic calibration averages Bellman targets over
groups of nearby predicted values, but the partition is chosen adaptively by the
monotonicity constraint \citep{barlow1972isotonic}. This provides
regularization without selecting \(B\). Since the identity map is monotone,
isotonic calibration can preserve an already reliable predictor while correcting
systematic monotone distortions.

\begin{enumerate}[label=\textbf{C\arabic*}, ref={C\arabic*}, leftmargin=1.5em, resume=cond]
\item \label{cond::C3}
\textbf{Finite variation.}
There exists \(V<\infty\) such that, almost surely, the function $t
\mapsto
E\!\left[
\mathcal T_\pi(V_K)(S)
\mid
\hat v(S)=t,\mathcal C_n
\right]$
has total variation at most \(V\).
\end{enumerate}

Condition~\ref{cond::C3} is standard in isotonic calibration
\citep{van2023causal} and controls the complexity of the one-dimensional
calibration curve. For histogram binning, the analogous complexity is governed
by the number of bins \(B\).

\begin{theorem}[Calibration error for isotonic calibration]
\label{theorem::Calerroriso}
Assume \ref{cond::C1}--\ref{cond::C3} with
\(\mathcal G=\mathcal G_{\mathrm{iso}}\). Then there exists a constant
\(C<\infty\) such that, for any \(K\in\mathbb N\), with probability at least
\(1-\delta\),
{\small
\begin{align*}
\mathrm{Cal}_{\ell^2}(V_K)
\le\;&
C\left\{
n^{-1/3}
+
\sqrt{\frac{\log(1/\delta)}{n}}
\right\}  +
\big\|
(\widehat w_\pi-w_\pi)
(\widehat q_{V_K}-q_{V_K})
\big\|_{2,\rho}  +
\big\|
\widehat T_\pi V_K
-
\widehat T_\pi V_{K-1}
\big\|_{2,\rho}.
\end{align*}
}
\end{theorem}
 
The isotonic analysis parallels the histogram case, but uses the adaptive
pooling induced by the pool-adjacent-violators algorithm. We do not assume the
calibration target is monotone; monotonicity is used only as a regularizer for
pooling nearby predicted values. The resulting \(n^{-1/3}\) rate matches the
classical isotonic regression rate \citep{chatterjee2013improved} and the rate
of histogram binning with \(B\asymp n^{1/3}\), but avoids choosing \(B\) in
advance.

 
Our histogram proof does not yield value-error bounds for fully iterated
isotonic calibration without Bellman completeness, because the PAVA partition can
change across Bellman iterations. Appendix~\ref{appendix::hybrid} gives a hybrid
procedure that first learns the partition by isotonic regression and then
applies the histogram analysis conditionally on that fixed partition.

\section{Synthetic experiments}
\label{sec:experiments}

We evaluate Bellman calibration as a post-hoc correction for fitted off-policy
value predictors. Each base learner first estimates \(\widehat Q(s,a)\), inducing
\(\widehat V(s)=\sum_a \pi(a\mid s)\widehat Q(s,a)\). Calibration then fits a
one-dimensional map \(g\) and reports \(g\{\widehat V(s)\}\). The map is learned
by importance-weighted fitted value iteration with weights
\(\omega_i=\pi(A_i\mid S_i)/b(A_i\mid S_i)\). The base learners are linear FQE
and random-feature FQE, where the latter applies linear Bellman regression after
a fixed random Fourier feature expansion. We consider linear, histogram, and
isotonic calibration. See Appendix \ref{app:experiments} for details.

The experiments isolate the failure mode illustrated in
Table~\ref{tab:rl-calibration-motivation}: the learned value predictor is informative but
mis-scaled. In the non-temporal shift experiments, we use 5-fold
cross-calibration \citep{van2023causal}: the value learner and calibrator are
fit out of fold, calibrated predictions are aggregated by pointwise medians, and
the uncalibrated comparator is the same learner refit on the full training
sample. In the temporal shift experiment, we train FQE on a large batch from an
old reward regime, fit \(g\) on a small held-out batch from the current reward
regime, and evaluate only on current-regime states. All main results use 100
independent replications and independent evaluation data. Values below one
indicate improvement.


\begin{figure}[t]
  \centering
  \includegraphics[width=0.98\linewidth]{figures/calibration_story_compact.pdf}
  \caption{\textbf{Post-hoc value correction improves value MSE and Bellman calibration error when
  the raw prediction is informative but mis-scaled.} Each bar is a ratio to the matched
  uncalibrated baseline. Columns show affine misspecification, finite Bellman
  iteration, temporal reward shift, and monotone misspecification.}
  \label{fig:calibration-story}
\end{figure}

\begin{table*}[t]
\centering
\caption{\textbf{Value-scale correction across offline RL failure modes.}
Ratios are relative to the matched uncalibrated baseline. Win rate is the
seed-matched true-\(V\) MSE win rate.}
\label{tab:main-simulation-results}
\scriptsize
\setlength{\tabcolsep}{4pt}
\begin{tabular}{@{}p{0.45\linewidth}lrrr@{}}
\toprule
Regime and learner & Calibrator & Rel. \(V\) MSE & Rel. cal. err. & Win rate \\
\midrule
Affine misspec., linear FQE & Linear & 0.750 & 0.631 & 0.94 \\
Affine misspec., linear FQE & Isotonic & 0.734 & 0.764 & 0.94 \\
Finite iteration, RF FQE \(K=4\) & Linear & 0.780 & 0.703 & 0.94 \\
Finite iteration, RF FQE \(K=4\) & Isotonic & 0.776 & 0.736 & 0.94 \\
Temporal reward shift, RF FQE & Linear recent & 0.630 & 0.717 & 0.97 \\
Temporal reward shift, RF FQE & Isotonic recent & 0.627 & 0.916 & 0.98 \\
Monotone misspec., linear FQE & Linear & 0.972 & 0.986 & 0.64 \\
Monotone misspec., linear FQE & Isotonic & 0.916 & 0.537 & 0.79 \\
Monotone misspec., linear FQE & Histogram & 0.917 & 0.544 & 0.77 \\
\bottomrule
\end{tabular}
\end{table*}

Figure~\ref{fig:calibration-story} and
Table~\ref{tab:main-simulation-results} show consistent gains across three
failure modes. Under affine misspecification, linear calibration suffices,
reducing true-\(V\) MSE by about \(25\%\) and Bellman calibration error by about
\(37\%\). With finite Bellman iteration error, random-feature FQE value predictions retain useful
ranking information but are under-scaled; linear and isotonic calibration reduce
true-\(V\) MSE by about \(22\%\) and Bellman calibration error by
\(26\)--\(34\%\). Under temporal reward shift, FQE trained on old-regime data is
recalibrated on a small recent batch and evaluated in the current regime. This
recent linear calibration reduces current-regime MSE by about \(37\%\) and
Bellman calibration error by about \(28\%\), while also outperforming a
same-size retrain on current-regime data. The monotone-misspecification rows show when richer calibrators help. Linear
calibration only mildly changes Bellman calibration error, whereas isotonic
calibration nearly halves it and also improves true-\(V\) MSE. Histogram calibration is comparable to isotonic in this setting. We emphasize
linear and isotonic calibration because they are low-tuning post-hoc choices:
linear calibration targets affine value-scale errors, while isotonic calibration
allows more general monotone distortions.

 
 


 

 

\bibliographystyle{plainnat}
\bibliography{../../common/ref}


 

 




 
\appendix

\section{Conclusion and Future Work}

We introduced \emph{Iterated Bellman Calibration}, a post-hoc method for
improving the reliability of value predictors in infinite-horizon MDPs, and
established finite-sample guarantees under weak assumptions. Our analysis also
highlights the central role of distribution shift. Future work includes
reducing off-policy mismatch through stationary-density-ratio weighting or
related emphatic-style corrections
\citep{vdLaanKallus2025FQE,mahmood2015emphatic,patterson2022generalized},
developing smoother calibrators that extrapolate better under distribution
shift \citep{platt1999probabilistic,jiang2011smooth,guo2017calibration}, and
broadening the empirical study of calibration in RL.



\section{Additional Experimental Details}
\label{app:experiments}

\paragraph{Data-generating processes.}
All four simulation panels use discounted MDPs with continuous states
\(S_t\in\mathbb{R}^d\), actions \(A_t\in\{1,2,3\}\), and softmax behavior and
target policies. The target policy has logits \(s^\top w_a^\pi/\tau_\pi\). The
behavior policy uses shifted logits \(s^\top w_a^b/\tau_b\), where
\(w_a^b=w_a^\pi-\Delta u_a\). The main experiments use the moderate-shift
setting. Initial states are Gaussian, \(S_0\sim N(\mu_0,I_d)\). Except where
noted below, transitions and rewards are
\[
  S_{t+1}
  =
  \tanh\{M S_t+\alpha_{A_t}+0.22\sin(S_t)+0.08\cos(S_t^{\mathrm{rev}})\}
  + \sigma_\varepsilon \varepsilon_t ,
\]
and
\[
  r_0(s,a)
  =
  s^\top\beta
  +0.65\sin(s_1)
  +0.35\cos(s_2)
  +s^\top\eta_a
  +\rho_a
  -0.06\|s\|_2^2 .
\]
Observed rewards add mean-zero Gaussian noise. The ground-truth value \(V^\pi(s)\) is
computed by independent target-policy Monte Carlo rollouts using the conditional
mean reward. Training transitions, calibration folds, Bellman-evaluation
transitions, initial evaluation states, and Monte Carlo rollouts use independent
seeds.

The four panels isolate different value-scale errors. In the affine
misspecification panel, the same nonlinear MDP is used, but the reward is
transformed to
\[
  r_{\mathrm{aff}}(s,a)=1.25+0.7\,r_0(s,a).
\]
The learner is linear FQE with a deliberately restricted feature class. This
produces informative value predictions with an affine slope/intercept error, so
linear calibration should be sufficient.

In the finite-iteration panel, the dynamics are made persistent,
\[
  S_{t+1}=0.92S_t+\alpha_{A_t}+0.03\sin(S_t)+\sigma_\varepsilon\varepsilon_t,
\]
with \(\gamma=0.985\) and horizon \(120\). Rewards are smooth but long-horizon:
\[
  r_{\mathrm{fin}}(s,a)
  =
  0.8+0.9\tanh(s_1)+0.25s_2+\rho_a
  +0.15\tanh(s_1+0.5s_2)+s^\top\eta_a .
\]
Random-feature FQE is stopped after \(K=4\) Bellman iterations, so its values
retain useful ranking information but are under-scaled relative to the
long-horizon target.

In the temporal reward-shift panel, the old and current environments share
dynamics, policies, and state distribution. Only the reward changes:
\[
  r_{\mathrm{current}}(s,a)=2.0+1.5\,r_{\mathrm{old}}(s,a).
\]
We train random-feature FQE on a large old-regime behavior batch
(\(n_{\mathrm{old}}=2000\)), fit the calibration map on a small recent
current-regime batch (\(n_{\mathrm{recent}}=100\)), and evaluate on independent
current-regime ground-truth values and Bellman-calibration transitions. A same-size
current-regime retrain comparator uses the same learner.

In the monotone-misspecification panel, the nonlinear reward is transformed by
a monotone saturation,
\[
  r_{\mathrm{mono}}(s,a)
  =
  \operatorname{sign}\{r_0(s,a)\}\sqrt{|r_0(s,a)|+1}.
\]
The linear FQE value prediction remains sufficiently ordered to be useful but has a
nonlinear value-scale distortion. This is the setting where isotonic and
histogram calibration should improve over an affine correction.

\paragraph{Baseline value predictors.}
Each method first fits a \(Q\)-function from the offline behavior-policy data
and then forms
\[
  \widehat V(s)=\sum_a \pi(a\mid s)\widehat Q(s,a).
\]
The experiments use random-feature FQE and linear FQE because they give
transparent error mechanisms: finite Bellman-iteration bias and affine or
monotone value-scale misspecification.

\paragraph{Cross-calibration.}
For fold \(k\), the \(Q\)-learner is fit on the other folds. The held-out fold
supplies \((\widehat V^{(-k)}(S_i),\widehat V^{(-k)}(S_i'),R_i)\). These
out-of-fold tuples are pooled to fit a single value-space calibration map by
importance-weighted fitted value iteration. At evaluation time, calibrated
values are aggregated as
\[
  \widehat V_{\mathrm{cal}}(s)
  =
  \mathrm{median}_{k=1,\ldots,K_{\mathrm{fold}}}
  g\!\left(\widehat V^{(-k)}(s)\right).
\]
The matched uncalibrated comparator is a \(Q\)-learner refit on all available
training data.

\paragraph{Value-space fitted value iteration.}
Let \(\omega_i=\pi(A_i\mid S_i)/b(A_i\mid S_i)\). In all reported runs the
action-ratio weights are clipped at 20 and normalized to mean one within each
calibration fit. Starting from the identity map, calibration iterates
\[
  g_{t+1}\in\arg\min_g\sum_i \omega_i
  \left[g(\widehat V(S_i))-\{R_i+\gamma g_t(\widehat V(S_i'))\}\right]^2.
\]
The implemented calibrators are weighted affine regression, weighted histogram
binning, weighted isotonic regression, and a weighted isotonic-histogram
hybrid.

\paragraph{Metrics.}
The scalar OPE estimate is the mean of \(g(\widehat V(S_0))\) over independent
initial states. The main prediction metric is held-out value-function MSE,
\[
  \mathbb{E}\!\left[\{g(\widehat V(S))-V^\pi(S)\}^2\right],
\]
estimated using independent target-policy ground-truth values. Bellman calibration
error is estimated on independent Bellman-evaluation transitions by binning
\(g(\widehat V(S))\) into quantile bins and averaging the weighted squared gaps
between bin-mean predictions and bin-mean Bellman outcomes
\[
  r_0(S,A)+\gamma g(\widehat V(S')) .
\]
We use independent evaluation batches  of size 50,000 and 50 bins.  

\paragraph{Reproducibility.}
All experiments use fixed configurations and 100 independent replications. The
code archive includes the exact seeds, raw replication-level outputs, grouped
summaries, and figure-generation scripts.
 

 



\section{Additional details and extensions}
 



\subsection{Hybrid Isotonic--Histogram Calibration}
\label{appendix::hybrid}

Isotonic iterated Bellman calibration adaptively pools nearby predicted values
through the monotonicity constraint. This adaptivity is useful in practice, but
it complicates the estimation-error analysis because applying
Algorithm~\ref{alg::cal-gen} with \(\mathcal G=\mathcal G_{\mathrm{iso}}\) may
produce a different data-dependent partition at each iteration. The iterates
therefore need not correspond to a single aggregated Bellman operator.

The hybrid isotonic--histogram procedure separates partition selection from
Bellman iteration. It first performs one isotonic regression of the initial
Bellman target as a function of \(\hat v(S)\), extracts the flat regions of the
fitted isotonic function as bins, and then runs histogram-based Bellman
calibration on this fixed partition. Thus the method retains the adaptivity of
isotonic calibration while restoring the fixed aggregated-operator structure
used in the histogram analysis.

\begin{algorithm}[t]
\caption{\textbf{Iso--Hist Iterated Bellman Calibration}}
\label{alg::iso-hist}
\begin{algorithmic}[1]

\BCINPUT Value predictor $\hat{v}$, calibration data $\mathcal C_n$,
Bellman target-estimate map $v\mapsto\widehat Y_v$, iterations $K$

\STATE {\small\textbf{Stage 1: fit outcome-adaptive partition}}
\STATE \hspace{1em} $\widehat Y_i^{(0)} \gets \widehat Y_{\hat v}(O_i)$
\STATE \hspace{1em} $g_{\mathrm{iso}} \gets 
       \arg\min_{g \in \mathcal G_{\mathrm{iso}}}
       \sum_{i=1}^n (\widehat Y_i^{(0)} - g(\hat{v}(S_i)))^2$
\STATE \hspace{1em} Extract bins $\hat I_1,\dots,\hat I_{\hat B}$
       from flat regions of $g_{\mathrm{iso}}$
\STATE \hspace{1em} Define $\widehat{\mathcal G}_{\hat B} :=
       \{g:g  \text{ is constant on each }\hat I_b\}$

\STATE {\small\textbf{Stage 2: histogram calibration}}
\STATE \hspace{1em} Apply Algorithm~\ref{alg::cal-gen} with class 
       $\widehat{\mathcal G}_{\hat B}$ to obtain $V_K$

\BCOUTPUT $V_K$

\end{algorithmic}
\end{algorithm}

Once the partition is fixed, the guarantees of Section~\ref{sec:hist} apply
directly, provided the number of bins produced by the initial isotonic step
admits a deterministic high-probability bound. More generally, the same
argument applies to other data-adaptive binning methods that produce a fixed
partition before the Bellman iterations begin.

\begin{theorem}[Calibration and Estimation Error for Alg.~\ref{alg::iso-hist}]
\label{theorem::regretiso}
Let $V_K$ denote the output of Algorithm~\ref{alg::iso-hist}. Suppose
the isotonic regression step produces at most $\hat B$ bins, where 
$\hat B \le B_n$ for some deterministic sequence $B_n$ with probability at least $1 - \delta$. Assume \ref{cond::C1}-\ref{cond::C3} with $B = B_n$.
Then, with probability at least $1 - 2\delta$, the guarantees in
Theorems~\ref{theorem::Calerrorhist} and~\ref{theorem::regretHist} hold for
$V_K$ with $B := B_n$.
\end{theorem}

Existing isotonic regression theory suggests that the effective number of
induced bins satisfies \(\widehat B=O(n^{1/3})\) when the calibration function
is continuously differentiable \citep{deng2021confidence}; this rate can also
be enforced by directly constraining the number of constant segments
\citep{van2024self}. For the approximation term in
Theorem~\ref{theorem::regretHist} to remain comparable to the initial error,
the induced histogram class must also approximate the identity map on the range
of \(\hat v\). This is plausible when \(\hat v\) is already close to \(v_0\);
when \(\hat v\) is substantially miscalibrated, isotonic regression may instead
pool poorly predicted regions into a common calibrated value.

\subsection{Alternative Calibration Targets}
\label{app:alternative-calibration-targets}

The Bellman calibration criterion studied in the main text calibrates a
candidate value predictor \(v\) with respect to its own Bellman target,
\[
E_\pi\!\left[R+\gamma v(S')\mid v(S)\right].
\]
A stronger alternative would instead calibrate \(v\) directly against the true
value function \(v_0\). Specifically, we say that \(v\) is
\emph{strongly Bellman calibrated} if
\[
v(S)
=
E\!\left[v_0(S)\mid v(S)\right]
\quad\text{almost surely}.
\]
Equivalently, among states assigned the same predicted value, the prediction
agrees on average with the true long-run return under \(\pi\).

This target is conceptually appealing, but substantially less practical. It
requires estimating the conditional mean of the unknown quantity \(v_0(S)\)
given the predicted value \(v(S)\). In general, this conditional mean is not directly
available from one-step transitions alone and is typically no easier to estimate
than the underlying value function itself
\citep{kallus2020double,kallus2022efficiently}. We therefore focus on the
weaker but directly attainable Bellman-target calibration criterion developed
in the main text.



 

\section{Experimental results and info}
\label{appendx::exp}
\subsection{Synthetic CRM environment}
 

We evaluate all methods in a synthetic customer-relationship-management (CRM)
Markov decision process (MDP) that mimics monthly retention and revenue dynamics
for a subscription service. The state $s_t \in \mathbb{R}^6$ encodes tenure in
months, engagement, fatigue, value segment, price sensitivity, and an indicator
of whether the customer is active. At each time $t$, the agent selects one of
three actions: no promotion, light promotion, or strong promotion. Rewards
correspond to monthly revenue after discounts and are zero after churn.

Transitions capture key CRM effects. Churn probability depends on tenure,
engagement, fatigue, and action via a logistic model; visit probability depends
on engagement, fatigue, and action; and revenue conditional on visit scales a
baseline value segment by an action-dependent uplift and a price-sensitivity
effect, with log-normal noise. Engagement decays over time but is boosted by
successful promotions, while fatigue increases with promotion intensity and
decays slowly otherwise. Once a customer churns or reaches a maximum tenure of
60 months, the process enters an absorbing state.

We simulate offline datasets of $n_{\text{cust}} = 50{,}000$ customers over a
horizon of $T = 24$ months with discount factor $\gamma = 0.99$. Behavior data are
generated by a fixed heuristic policy that sends light promotions by default,
suppresses promotions for highly engaged and fatigued customers, and
occasionally sends strong promotions to low-engagement, high-value customers. We
evaluate off-policy value estimation for a deterministic \emph{aggressive}
revenue-seeking target policy that sends strong promotions to low-engagement,
high-sensitivity customers and avoids promotions for highly engaged, high-value
customers.

Ground-truth values $V^\pi(s_0)$ are estimated via Monte Carlo rollouts in the
same environment, and we report $(1-\gamma)$-scaled RMSE between the estimated and
Monte Carlo values at initial states, averaged across 50 independent runs.


\paragraph{Base estimators.}
We consider three base estimators: boosted FQE, linear FVI, and neural FVI.
The linear model uses a hand-crafted semi-structured feature map consisting of
the six raw state variables, normalized and log-transformed tenure, squared
terms for engagement, fatigue, value segment, and price sensitivity, three
state-only interaction terms, and additive action one-hot indicators, with no
state-action interactions. The neural model takes the raw state concatenated
with the action as input and uses a feedforward ReLU network with three hidden
layers of width 256. Both FVI estimators are trained by iterating one-step
Bellman regressions: the linear model uses ridge-regularized least squares,
while the neural model is trained for 100 outer iterations using Adam with
learning rate \(0.005\), one epoch per iteration, and mini-batch size 2000.
We do not tune the number of neural iterations, since part of the goal is to
illustrate a realistic setting in which a base estimator remains poorly
calibrated. Selected neural iterates are saved for the snapshot comparisons.
Both FVI estimators use an internal 80/20 train-validation split. The boosted
baseline is implemented as a LightGBM FQE model trained on the concatenated
state-action input. Under sample splitting, base estimators are fit on the
training subset and Bellman calibrators on the held-out calibration subset; we
also retrain each base estimator on the full dataset to obtain the
\emph{Full Raw} baselines. 



For the calibration methods, we use 1000 fitted-value-iteration updates.


 

 \subsubsection{Neural snapshots under sample splitting}

Tables~\ref{tab:neural-snap-split-50k-p05} and \ref{tab:neural-snapshots-50k-ptrain-075} report neural FVI snapshot performance under sample splitting for the additional training proportions \(p_{\mathrm{train}}=0.5\) and \(p_{\mathrm{train}}=0.75\). Together, they show how the effect of calibration evolves over training and how it depends on the amount of data reserved for fitting versus calibration.

\begin{table}[h]
\centering
\small
\begin{tabular}{r cccc}
\toprule
Iter & Full Raw & Iso & Histogram & Hybrid Iso \\
\midrule
10  & 1.309$\pm$0.06 & \textbf{0.581$\pm$0.05} & 0.582$\pm$0.05 & 0.581$\pm$0.05 \\
25  & 0.491$\pm$0.13 & \textbf{0.456$\pm$0.06} & 0.456$\pm$0.06 & \textbf{0.456$\pm$0.06} \\
50  & \textbf{0.320$\pm$0.07} & 0.381$\pm$0.05 & 0.382$\pm$0.05 & 0.382$\pm$0.05 \\
100 & \textbf{0.316$\pm$0.05} & 0.350$\pm$0.05 & 0.351$\pm$0.05 & 0.350$\pm$0.05 \\
\bottomrule
\end{tabular}
\caption{Neural snapshot performance with sample splitting for \(n=50{,}000\) and training proportion \(p_{\mathrm{train}}=0.5\). Lower is better.}
\label{tab:neural-snap-split-50k-p05}
\end{table}

\begin{table}[h]
\centering
\small
\begin{tabular}{c cccc}
\toprule
Iteration & Full Raw & Iso & Histogram & Hybrid Iso \\
\midrule
10  & 1.28$\pm$0.05 & 0.52$\pm$0.05 & \textbf{0.52$\pm$0.05} & 0.52$\pm$0.05 \\
25  & 0.44$\pm$0.12 & \textbf{0.38$\pm$0.03} & 0.38$\pm$0.04 & 0.38$\pm$0.04 \\
50  & 0.34$\pm$0.06 & \textbf{0.33$\pm$0.04} & 0.33$\pm$0.04 & 0.33$\pm$0.04 \\
100 & 0.33$\pm$0.07 & \textbf{0.31$\pm$0.04} & 0.31$\pm$0.04 & 0.31$\pm$0.04 \\
\bottomrule
\end{tabular}
\caption{Neural FVI snapshot results for \(n=50{,}000\) with \(p_{\mathrm{train}}=0.75\). Lower is better.}
\label{tab:neural-snapshots-50k-ptrain-075}
\end{table}


 


\subsubsection{Comparing to the uncalibrated (raw) model without retraining on the full data}
Table~\ref{tab:main-split} compares calibrated estimators to the corresponding raw model fit only on the training split, so any improvements reflect calibration rather than retraining on the full dataset. This comparison is relevant when retraining on all available data is costly or when a gold-standard calibration dataset is available, such as recent trajectories or on-policy data, alongside a larger historical dataset used for initial fitting. Table~\ref{tab:neural-snap-split} further examines this comparison for the neural method across training iterations, illustrating how the effect of calibration evolves over the course of fitting.


\begin{table}[h]
\centering
\small
\begin{tabular}{r l ccc}
\toprule
$n$ & Model & Raw & Iso & Hybrid Iso \\
\midrule
\multirow{3}{*}{10{,}000}
& Boost & 0.681$\pm$0.1 & \textbf{0.671$\pm$0.1} & 0.697$\pm$0.2 \\
& Lin  & 0.640$\pm$0.04 & \textbf{0.612$\pm$0.05} & 0.641$\pm$0.05 \\
& Neur  & 0.582$\pm$0.2 & 0.550$\pm$0.2 & \textbf{0.520$\pm$0.1} \\
\midrule
\multirow{3}{*}{50{,}000}
& Boost & 0.433$\pm$0.06 & \textbf{0.429$\pm$0.06} & 0.438$\pm$0.08 \\
& Lin  & 0.625$\pm$0.02 & \textbf{0.614$\pm$0.03} & 0.616$\pm$0.03 \\
& Neur  & 0.419$\pm$0.1 & 0.383$\pm$0.08 & \textbf{0.374$\pm$0.08} \\
\midrule
\multirow{3}{*}{100{,}000}
& Boost & 0.360$\pm$0.05 & 0.358$\pm$0.05 & \textbf{0.353$\pm$0.07} \\
& Lin  & 0.616$\pm$0.02 & \textbf{0.605$\pm$0.02} & 0.606$\pm$0.02 \\
& Neur  & 0.379$\pm$0.09 & 0.351$\pm$0.05 & \textbf{0.342$\pm$0.05} \\
\bottomrule
\end{tabular}
\caption{Main results with sample splitting at training proportion \(p_{\mathrm{train}}=0.5\). The raw estimator is fit only on the training split to isolate the effect of calibration.}
\label{tab:main-split}
\end{table}




\begin{table}[h]
\centering
\small
\begin{tabular}{r cccc}
\toprule
Iter & Raw & Iso & Quantile & Hybrid Iso \\
\midrule
10  & 1.424$\pm$0.1 & 1.330$\pm$0.1 & \textbf{0.735$\pm$0.2} & 0.737$\pm$0.2 \\
25  & 0.665$\pm$0.1 & 0.646$\pm$0.1 & 0.577$\pm$0.1 & \textbf{0.571$\pm$0.1} \\
50  & 0.612$\pm$0.2 & 0.585$\pm$0.2 & 0.555$\pm$0.1 & \textbf{0.547$\pm$0.1} \\
100 & 0.582$\pm$0.2 & 0.550$\pm$0.2 & 0.524$\pm$0.1 & \textbf{0.520$\pm$0.1} \\
\bottomrule
\end{tabular}
\caption{Neural FVI snapshot performance with sample splitting at training proportion \(p_{\mathrm{train}}=0.5\). The raw estimator is fit using only the training split, so differences reflect the effect of calibration rather than access to additional data.}\label{tab:neural-snap-split}
\end{table}






 


 

 




\section{Stationary Measures and Contraction Results}
\label{appendix::stationary}

\begin{lemma}[Bellman contraction under a stationary measure]
\label{lemma::contraction}
Let $P$ be a Markov operator with stationary distribution $\mu$ (i.e., $\mu=\mu P$),
and define the Bellman operator $\mathcal T f := r + \gamma P f$.
Then $P$ is nonexpansive in $L^2(\mu)$:
\[
\|P(f-g)\|_{2,\mu} \;\le\; \|f-g\|_{2,\mu}.
\]
Consequently, $\mathcal T$ is a $\gamma$-contraction:
\[
\|\mathcal T (f-g)\|_{2,\mu}
\;\le\;
\gamma\,\|f-g\|_{2,\mu}.
\]
\end{lemma}
\begin{proof}
Because $\mu$ is stationary for $P$, the operator $P$ is nonexpansive in
$L^2(\mu)$. Jensen's inequality implies
\[
(P h)^2 \le P(h^2) \qquad\text{pointwise}.
\]
Integrating both sides with respect to $\mu$ and using $\mu = \mu P$ yields
\[
\|P h\|_{2,\mu}^2
= \int (P h)^2 \, d\mu
\le \int P(h^2)\, d\mu
= \int h^2\, d\mu
= \|h\|_{2,\mu}^2.
\]
Now apply this with $h = f - g$. Since the reward function $r$ cancels,
\[
\mathcal{T}f - \mathcal{T}g
= \gamma\, P(f-g).
\]
Therefore,
\[
\|\mathcal{T} f - \mathcal{T} g\|_{2,\mu}
= \gamma \|P(f-g)\|_{2,\mu}
\le \gamma \|f-g\|_{2,\mu},
\]
which proves the claim.
\end{proof}

\begin{lemma}[Stationarity under coarsening]
\label{lemma:aggregated-stationary}
Let $P$ be a Markov kernel on $\mathcal S$ and let $\lambda$ be stationary for $P$
(i.e., $\lambda P = \lambda$).
For any measurable coarsening map $g:\mathcal S\to\mathbb R$, define
\[
P_g f(s)
:=
\mathbb{E}[\,f(S') \mid g(S)=g(s)\,],
\]
where $(S,S')$ has joint law $\lambda(ds)\,P(s,ds')$.
Then $\lambda$ is stationary for $P_g$.
\end{lemma}

\begin{proof}
Let $(S,S')\sim \lambda(ds)\,P(s,ds')$, so $S\sim\lambda$ and $S'\mid S\sim P(S,\cdot)$.
Then
\begin{align*}
\int P_g f(s)\,\lambda(ds)
&=
\mathbb{E}_{\lambda,P}\!\left[P_g f(S)\right] \\[0.3em]
&=
\mathbb{E}_{\lambda,P}\!\left[\mathbb{E}_{\lambda,P}\!\left[f(S') \mid g(S)\right]\right] \\[0.3em]
&=
\mathbb{E}_{\lambda,P}[f(S')].
\end{align*}
Since $\lambda$ is stationary for $P$, the marginal of $S'$ is again $\lambda$, hence
$\mathbb{E}_{\lambda,P}[f(S')] = \int f(s)\,\lambda(ds)$.
\end{proof}



\section{Proof of Theorem \ref{theorem::calrefine}}

 

 
 


\begin{proof}[Proof of Theorem \ref{theorem::calrefine}]
Define $P_{\pi,\hat v} := \Pi_{\hat v}P_{\pi}$ and $r_{\pi,\hat{v}} := \Pi_{\hat v} r_{\pi}$. Because $v_0$ and $V_{\hat v}^\star$ satisfy the fixed-point equations
\[
v_0 = \mathcal{T}_\pi v_0,
\qquad
V_{\hat v}^\star = \mathcal B_{\pi,\hat v} V_{\hat v}^\star
= r_{\pi,\hat{v}} + \gamma\,P_{\pi,\hat v}V_{\hat v}^\star,
\]
we begin by writing
\[
V_{\hat v}^\star - v_0
=
\bigl( \mathcal B_{\pi,\hat v}V_{\hat v}^\star
      - \mathcal B_{\pi,\hat v}v_0 \bigr)
+
\bigl( \mathcal B_{\pi,\hat v}v_0 - \mathcal{T}_\pi v_0 \bigr).
\]
Using linearity of the conditional expectation in
$\mathcal B_{\pi,\hat v}$ on differences,
\[
\mathcal B_{\pi,\hat v}V_{\hat v}^\star
-
\mathcal B_{\pi,\hat v}v_0
=
\gamma\,P_{\pi,\hat v}(V_{\hat v}^\star - v_0),
\]
so
\[
V_{\hat v}^\star - v_0
=
\gamma\,P_{\pi,\hat v}(V_{\hat v}^\star - v_0)
+
\bigl(\mathcal B_{\pi,\hat v}v_0 - v_0\bigr),
\]
where we used $\mathcal{T}_\pi v_0 = v_0$ for the second term.

Taking $L^2(\rho)$ norms and applying the triangle inequality,
\begin{align*}
\|V_{\hat v}^\star - v_0\|_{2,\rho}
&\le
\gamma\,\|P_{\pi,\hat v}(V_{\hat v}^\star - v_0)\|_{2,\rho}
\\
&\qquad +
\|\mathcal B_{\pi,\hat v}v_0 - v_0\|_{2,\rho}.
\end{align*}
By Lemma \ref{lemma::contraction},
\[
\|P_{\pi,\hat v}h\|_{2,\rho}
\le
\|h\|_{2,\rho}.
\]
Applying this with $h=V_{\hat v}^\star - v_0$ gives
\[
\|V_{\hat v}^\star - v_0\|_{2,\rho}
\le
\gamma\,\|V_{\hat v}^\star - v_0\|_{2,\rho}
+
\|\mathcal B_{\pi,\hat v}v_0 - v_0\|_{2,\rho}.
\]

Note that $\mathcal B_{\pi,\hat v}v = \Pi_{\hat v}\mathcal{T}_\pi v$,
where $\Pi_{\hat v}$ is the $L^2$ projection onto functions of $\hat{v}$.
Thus $\mathcal B_{\pi,\hat v}v_0 = \Pi_{\hat v}v_0$, and therefore
\[
\|V_{\hat v}^\star - v_0\|_{2,\rho}
\le
\gamma\,\|V_{\hat v}^\star - v_0\|_{2,\rho}
+
\|\Pi_{\hat v}v_0 - v_0\|_{2,\rho}.
\]

Rearranging,
\[
\|V_{\hat v}^\star - v_0\|_{2,\rho}
\le
\frac{1}{1-\gamma}\,
\|\Pi_{\hat v}v_0 - v_0\|_{2,\rho}.
\tag{$\star$}
\]
Finally, by the triangle inequality,
\[
\|\hat{v} - v_0\|_{2,\rho}
\le
\|\hat{v} - V_{\hat v}^\star\|_{2,\rho}
+
\|V_{\hat v}^\star - v_0\|_{2,\rho},
\]
which yields the stated result.
\end{proof}



\section{Proof of Theorem \ref{theorem::DRpseudo}}

\begin{proof}[Proof of Theorem \ref{theorem::DRpseudo}]
Fix $v$ and treat the nuisance estimates as fixed functions. By definition,
\[
\widehat T_\pi v(s)
=
\mathbb{E}\Big[
(\pi \hat q_v)(S)
+
\widehat w_{\pi}(A \mid S)
\big\{
  R + \gamma\, v(S') - \hat{q}_v(S,A)
\big\}
\;\Big|\; S = s
\Big].
\]
Since $(\pi \hat q_v)(S)$ depends only on $S$,
\[
\mathbb{E}\big[(\pi \hat q_v)(S)\mid S=s\big] = (\pi \hat q_v)(s).
\]
For the second term, using the definition of $b_0$ and $q_v := r_0 + \gamma P v$,
\begin{align*}
&\mathbb{E}\Big[
\widehat w_{\pi}(A \mid S)
\big\{
  R + \gamma\, v(S') - \hat{q}_v(S,A)
\big\}
\;\Big|\; S = s
\Big] \\
&\quad=
b_0\Big\{
\widehat w_\pi(\cdot \mid s)\,
\mathbb{E}\big[
R + \gamma\, v(S') - \hat q_v(s,\cdot)
\mid S = s, A = \cdot
\big]
\Big\}(s) \\
&\quad=
b_0\big\{\widehat w_\pi (q_v - \hat q_v)\big\}(s).
\end{align*}
Hence
\[
\widehat T_\pi v(s)
=
(\pi \hat q_v)(s)
+
b_0\big\{\widehat w_\pi (q_v - \hat q_v)\big\}(s).
\]

Subtracting $\mathcal T_\pi(v)(s) = (\pi q_v)(s)$ yields
\begin{align*}
\widehat T_\pi v(s) - \mathcal T_\pi(v)(s)
&=
(\pi \hat q_v - \pi q_v)(s)
+
b_0\big\{\widehat w_\pi (q_v - \hat q_v)\big\}(s) \\
&=
(\pi (\hat q_v - q_v))(s)
-
b_0\big\{\widehat w_\pi (\hat q_v - q_v)\big\}(s).
\end{align*}
Using $(\pi f)(s) = b_0\{w_\pi f\}(s)$ with $f = \hat q_v - q_v$ gives
\[
(\pi (\hat q_v - q_v))(s)
=
b_0\{w_\pi (\hat q_v - q_v)\}(s),
\]
so
\[
\widehat T_\pi v(s) - \mathcal T_\pi(v)(s)
=
b_0\big\{(w_\pi - \widehat w_\pi)(\hat q_v - q_v)\big\}(s).
\]
This proves the claim.
\end{proof}



\section{Notation and Maximal Inequalities}

Let $P_0$ denote the joint distribution of $(S, A, R, S')$ induced by the
behavior policy, and let $P_n$ denote the empirical measure of the calibration
sample $\mathcal{C}_n$.



We define empirical \(L^2\) norms with respect to the state and next-state samples.  For any state function \(f\), let
\[
\|f\|_{n,S}
:= 
\Big( \tfrac{1}{n} \sum_{i=1}^{n} f(S_i)^2 \Big)^{1/2},
\quad
\|f\|_{n,S'}
:= 
\Big( \tfrac{1}{n} \sum_{i=1}^{n} f(S_i')^2 \Big)^{1/2},
\]
where \(S_i\) and \(S_i'\) denote the observed states and next states, respectively, in the calibration sample \(\mathcal{C}_n\).



 





Further define
\[
\widehat{\mathcal{G}}
:= 
\Big\{ (f_1 - f_2)\big(\widehat Y_{f_2} - f_2\big)
:\;
f_1, f_2 \in \mathcal H_B(\hat v) \Big\}.
\]
By assumption, both \(\hat{v}\) and the Bellman target-estimate map
\(v\mapsto\widehat Y_v\) are fixed conditional on the training data, which is
independent of the calibration sample \(\mathcal{C}_n\). Consequently, the
classes \(\mathcal H_B(\hat v)\) and \(\widehat{\mathcal{G}}\) are
non-random conditional on the training dataset.

For any distribution \(Q\) and any uniformly bounded function class \(\mathcal{F}\), let  
\(N(\varepsilon, \mathcal{F}, L^2(Q))\) denote the \(\varepsilon\)-covering number of \(\mathcal{F}\) under the \(L^2(Q)\) norm \citep{van1996weak}.  
Define the uniform entropy integral of \(\mathcal{F}\) by
\begin{equation*}
\mathcal{J}(\delta, \mathcal{F})
:=
\int_{0}^{\delta}
\sup_{Q}
\sqrt{\log N(\epsilon, \mathcal{F}, L^2(Q))}\, d\epsilon ,
\end{equation*}
where the supremum is taken over all discrete probability distributions \(Q\).

Finally, for two quantities \(x\) and \(y\), we write \(x \lessapprox y\) to mean that \(x\) is bounded above by \(y\) up to a universal constant that depends only on global constants appearing in our conditions.


\subsection{Local maximal inequality}

 Let $O_1,\ldots,O_n \in \mathcal{O}$ be independent random variables. For any function $f:\mathcal{O} \to \mathbb{R}$, define
\begin{align}
   \|f\|_{2,\bar P} := \sqrt{\frac{1}{n} \sum_{i=1}^n\mathbb{E}[f(O_i)^2]}.
\end{align}
 

 


For a star-shaped class of functions $\mathcal{F}$ and a radius $\delta \in (0,\infty)$,
define the localized Rademacher complexity
\[
\mathcal{R}_n(\mathcal{F}, \delta)
:=
\mathbb{E}\left[
\sup_{\substack{f \in \mathcal{F} \\\|f\|_{2,\bar P} \le \delta}}
\frac{1}{n} \sum_{i=1}^n \epsilon_i f(O_i)
\right],
\]
where $\epsilon_i$ are i.i.d.\ Rademacher random variables.

The following lemma provides a local maximal inequality and restates Lemma~11 of \cite{foster2023orthogonal}.


 


\begin{lemma}[Local maximal inequality]\label{lemma:loc_max_ineq}
Let $\mathcal{F}$ be a star-shaped class of functions satisfying
$\sup_{f\in\mathcal{F}} \|f\|_{\infty} \le M$.
Let $\delta > 0$ satisfy the critical radius condition
$\mathcal{R}_n(\mathcal{F},\delta) \le \delta^2$.
Suppose further that
$n^{-1/2}\sqrt{\log\log(1/\delta)} = o(\delta)$. Then there exists a universal constant $C>0$ such that, for all $u \ge 1$, with
probability at least $1 - e^{-u^2}$, every $f \in \mathcal{F}$ satisfies
\[
\frac{1}{n} \sum_{i=1}^n
\bigl(f(O_i) - \mathbb{E}[f(O_i)]\bigr)
\;\le\;
C\Bigl(
\delta^2
+
\delta\|f\|_{2,\bar P}
+
\frac{u\|f\|_{2,\bar P}}{\sqrt{n}}
+
\frac{M u^2}{n}
\Bigr).
\]
\end{lemma}

The following lemma bounds the localized Rademacher complexity in terms of the uniform entropy integral and is a direct consequence of Theorem~2.1 of \citet{van2011local}.

 

\begin{lemma}\label{lemma:local_rademacher_entropy}
Let \(\mathcal{F}\) be a star-shaped class of functions such that
\(\sup_{f\in\mathcal F}\|f\|_\infty \le M\). Then, for every \(\delta>0\),
\[
\mathcal{R}_n(\mathcal{F}, \delta)
\;\lesssim\;
\frac{1}{\sqrt n}\,\mathcal{J}(\delta,\mathcal{F})
\left(1+\frac{\mathcal{J}(\delta,\mathcal{F})}{\delta \sqrt n}\right),
\]
where the implicit constant depends only on \(M\).
\end{lemma}
\begin{proof}
This bound follows directly from the argument in the proof of
Theorem~2.1 of \citet{van2011local}; see in particular the step where
the local Rademacher complexity is controlled by the uniform entropy
integral for star-shaped classes.
\end{proof}
 

 

\section{Proof of Theorem \ref{theorem::Calerrorhist}}
 
 

\subsection{Technical lemmas}

 



\begin{lemma}
\label{lemma::entropynumbers}
Under our conditions,
\[
\mathcal{J}\left(\delta, \widehat{\mathcal{G}}\right)
\;\lesssim\;
\delta \sqrt{B \log(1/\delta)},
\]
where the implicit constant is independent of \(B\).
\end{lemma}
\begin{proof}
 

By Condition~\ref{cond::C2},
\[
\|\hat q_f-\hat q_g\|_{n,S,A}
\le
L\|f-g\|_{n,S'}
\]
almost surely.
For each $s$, by Jensen's inequality,
\[
|\pi h(s)|^2=\Big|\sum_a \pi(a\mid s)h(a,s)\Big|^2
\le \sum_a \pi(a\mid s)|h(a,s)|^2,
\]
so $\|\pi h\|_{n,S}\le \|h\|_{n,S,A}$. Hence, we also have
\[
\|\pi(\hat q_f-\hat q_g)\|_{n,S}
\le \|\hat q_f-\hat q_g\|_{n,S,A}
\le L \|f - g\|_{n, S'} .
\]
Hence, by the definition of \(\widehat Y_f\) and boundedness of
\(\widehat w_\pi\),
\[
\|\widehat Y_f - \widehat Y_g\|_{L^2(P_n)} 
\lesssim  
\|f - g\|_{n,S'}.
\]
Take $f, g \in \widehat{\mathcal{G}}$ with 
$f = (f_1 - f_2)\big(\widehat Y_{f_2} - f_2\big)$ 
and 
$g = (g_1 - g_2)\big(\widehat Y_{g_2} - g_2\big)$ 
for $f_1, f_2, g_1, g_2 \in \mathcal H_B(\hat v)$. 
By Lipschitz continuity of multiplication and boundedness of nuisances, we have
\[
\|f - g\|_{L^2(P_n)} 
\lesssim  
\|f_1 - g_1\|_{n, S} + \|f_2 - g_2\|_{n, S}+ \|f_2 - g_2\|_{n, S'}.
\]
Taking the supremum over all discrete distributions $Q$, we obtain
\[
\|f - g\|_{L^2(P_n)} 
\lesssim  
\sup_Q \|f_1 - g_1\|_{L^2(Q)} 
+ \sup_Q \|f_2 - g_2\|_{L^2(Q)} .
\]
Hence, by preservation of entropy integrals in \citet{van1996weak},
\[
\log N(\varepsilon, \widehat{\mathcal{G}}, L^2(P_n)) 
\lesssim 
\sup_Q \log N(\varepsilon, \mathcal H_B(\hat v), L^2(Q)) .
\]
Taking the supremum over discrete distributions $Q$ on both sides yields the uniform covering number bound
\[
\sup_Q \log N(\varepsilon, \widehat{\mathcal{G}}, L^2(Q)) 
\lesssim 
\sup_Q \log N(\varepsilon, \mathcal H_B(\hat v), L^2(Q)) .
\]


The class \(\mathcal H_B(\hat v)\) satisfies
\[
\sup_Q \log N\big(\varepsilon, \mathcal H_B(\hat v), L^2(Q)\big)
\;\le\;
\sup_{Q} \log N\big(\varepsilon, \widetilde{\mathcal{F}}_{B},
L^2(Q \circ \hat{v}^{-1})\big),
\]
where \(Q \circ \hat{v}^{-1}\) denotes the pushforward of \(Q\) under
\(\hat{v}\). Hence,
\[
\sup_Q \log N\big(\varepsilon, \mathcal H_B(\hat v), L^2(Q)\big)
\;\le\;
\sup_{Q} \log N\big(\varepsilon, \widetilde{\mathcal{F}}_{B}, L^2(Q)\big),
\]
where, by a slight abuse of notation, the supremum on the right-hand side is
taken over all discrete probability distributions \(Q\) on \(\mathbb{R}\).
The class \(\widetilde{\mathcal{F}}_{B}\) consists of all piecewise-constant
functions on \(\mathbb{R}\) taking at most \(B\) values. Therefore,
\(\widetilde{\mathcal{F}}_{B}\) has VC--subgraph dimension \(O(B)\), and
\citet{van1996weak} implies
\begin{align*}
    \sup_Q \log N(\varepsilon, \mathcal H_B(\hat v), L^2(Q))
&\lesssim
\sup_{Q} \log N\big(\varepsilon, \widetilde{\mathcal{F}}_{B}, L^2(Q)\big)\\
&\lesssim B \log(1/\varepsilon).
\end{align*}
 Consequently, $\mathcal{J}(\delta, \mathcal H_B(\hat v))
\lesssim \delta \sqrt{B \log(1/\delta)}$.
\end{proof}

 


\subsection{Proof of Theorem \ref{theorem::Calerrorhist}}
\begin{proof}[Proof of Theorem \ref{theorem::Calerrorhist}]
Denote $O_i := (S_i, A_i, R_i, S_i')$. The first-order optimality conditions for
$\theta_n^{(K)}$ imply that, for each bin $b \in [B]$,
\[
\frac{1}{n}\sum_{i=1}^n 
\mathbf{1}_{I_b}(\hat{v}(S_i))
\big\{
\widehat Y_{V_{K-1}}(O_i)
- V_K(S_i)
\big\}
= 0.
\]
Hence, for any $f \in \mathcal G_B$, which is a linear combination of these
indicators,
\begin{equation}
\label{eqn::scores}
\frac{1}{n}\sum_{i=1}^n 
f(\hat{v}(S_i))
\big\{
\widehat Y_{V_{K-1}}(O_i)
- V_K(S_i)
\big\}
= 0.
\end{equation}
Moreover, for any function $g:\mathbb{R}\to\mathbb{R}$, we may write
$g \circ V_K = f \circ \hat{v}$ for some $f$, and therefore,
for all such $g$,
\[
\frac{1}{n}\sum_{i=1}^n 
g(V_K(S_i))
\big\{
\widehat Y_{V_{K-1}}(O_i)
- V_K(S_i)
\big\}
= 0.
\]
Noting that \(\Gamma_0(V_K) - V_K\) is of the form
\(g \circ V_K\) for some function \(g\), we have
\begin{align*}
0
&=
\frac{1}{n}\sum_{i=1}^n 
\Big\{
\Gamma_0(V_K)(S_i) - V_K(S_i)
\Big\} \\
&\qquad\qquad\times
\Big\{
\widehat Y_{V_{K-1}}(O_i)
- V_K(S_i)
\Big\}.
\end{align*}
Adding and subtracting \(P_0\) yields
\begin{equation}
\label{eqn::basicineq}
\begin{aligned}
&P_0
\Big\{
\Gamma_0(V_K) - V_K
\Big\}
\Big\{
\widehat Y_{V_K}
- V_K
\Big\}\\
& = (P_0 - P_n)
\Big\{
\Gamma_0(V_K) - V_K
\Big\}
\Big\{
\widehat Y_{V_{K-1}}
- V_K
\Big\}\\
&\quad  + P_0
\Big\{
\Gamma_0(V_K) - V_K
\Big\}
\Big\{
\widehat Y_{V_K}
- \widehat Y_{V_{K-1}}
\Big\}.
\end{aligned}
\end{equation}


We rewrite the left-hand side of \eqref{eqn::basicineq}.
Adding and subtracting and applying the law of total expectation, the calibration error decomposes as
\begin{equation}
\label{eqn::calerrorbound}
\begin{aligned}
&\| \Gamma_0(V_K) - V_K\|^2\\
&=
P_0
\Big\{
\Gamma_0(V_K) - V_K
\Big\}
\Big\{
\widehat Y_{V_K}
- V_K
\Big\} \\
&\quad+
P_0
\Big\{
\Gamma_0(V_K) - V_K
\Big\}
\Big\{
\Gamma_0(V_K)
- \widehat Y_{V_K}
\Big\}.
\end{aligned}
\end{equation}

The first term on the right-hand side can be decomposed as in \eqref{eqn::basicineq}.
By the law of total expectation, the proof of Theorem~\ref{theorem::DRpseudo},
and the Cauchy--Schwarz inequality, the second term satisfies
\begin{align*}
&\left|
P_0
\Big\{
\Gamma_0(V_K) - V_K
\Big\}
\Big\{
\widehat Y_{V_K}
- \Gamma_0(V_K)
\Big\}
\right|
\\[0.4em]
&=
\left|
P_0
\Big\{
\Gamma_0(V_K) - V_K
\Big\}
\Big\{
\widehat Y_{V_K}
- \mathcal{T}_{\pi}(V_K)
\Big\}
\right|
\\[0.4em]
&\le\;
\|\Gamma_0(V_K) - V_K\|\;
\big\|
b_0\{(\widehat w_{\pi} - w_{\pi})\,
(\widehat q_{V_K} - q_{V_K})\}
\big\|_{2,\rho}.
\end{align*}

Hence, \eqref{eqn::basicineq} and \eqref{eqn::calerrorbound} together imply that
\begin{equation}
\label{eqn::basicineq2}
\begin{aligned}
&\|\Gamma_0(V_K) - V_K\|^2\\
&\le\;
(P_0 - P_n)
\Big\{
\Gamma_0(V_K) - V_K
\Big\}
\Big\{
\widehat Y_{V_{K-1}}
- V_K
\Big\}
\\[0.4em]
&\quad+\;
\|\Gamma_0(V_K) - V_K\|\;
\big\|
\widehat T_\pi V_K
- \widehat T_\pi V_{K-1}
\big\|
\\[0.4em]
&\quad+\;
\|\Gamma_0(V_K) - V_K\|\;
\big\|
b_0\{(\widehat w_{\pi} - w_{\pi})\,
(\widehat q_{V_K} - q_{V_K})\}
\big\|_{2,\rho}.
\end{aligned}
\end{equation}
where $\widehat T_\pi v(s)
:= \mathbb{E}[\widehat Y_v(O)\mid S=s]$.
Here, the second term on the right-hand side follows from \eqref{eqn::basicineq}, noting that
\begin{align*} &\left| P_0 \Big\{ \Gamma_0(V_K) - V_K \Big\} \Big\{ \widehat Y_{V_K} - \widehat Y_{V_{K-1}} \Big\} \right| \\[0.4em] &= \left| P_0 \Big\{ \Gamma_0(V_K) - V_K \Big\} \,\mathbb{E}\left[ \widehat Y_{V_K} - \widehat Y_{V_{K-1}} \mid S \right] \right| \\[0.4em] &\le \|\Gamma_0(V_K) - V_K\|\; \Big\| \mathbb{E}\left[ \widehat Y_{V_K} - \widehat Y_{V_{K-1}} \mid S \right] \Big\| \\[0.4em] &= \|\Gamma_0(V_K) - V_K\|\; \big\| \widehat T_\pi V_K - \widehat T_\pi V_{K-1} \big\|_{2,\rho}. \end{align*}
 




We now turn to bounding the empirical process term on the right-hand side of \eqref{eqn::basicineq2}. Observe that 
\[
\big(\Gamma_0(V_K) - V_K\big)
\big(\widehat Y_{V_{K-1}} - V_K\big)
\]
lies in a uniformly bounded subset of the class 
\[
\widehat{\mathcal{G}} := \big\{(f_1 - f_2)\big(\widehat Y_{f_2} - f_2\big) : 
f_1, f_2 \in \mathcal H_B(\hat v) \big\},
\]
By Lemma~\ref{lemma::entropynumbers}, it holds that $\mathcal{J}\left(\delta, \widehat{\mathcal{G}}\right)
\;\lesssim\;
\delta \sqrt{B \log(1/\delta)}.$
By assumption, \(\widehat{\mathcal{G}}\) is a fixed, nonrandom function class
conditional on the training data, which is independent of the calibration
sample \(\mathcal{C}_n\). Hence, applying Lemma~\ref{lemma:local_rademacher_entropy}
conditional on the training data, we obtain
\begin{align*}
\mathcal{R}_n(\widehat{\mathcal{G}}, \delta)
&\;\lesssim\;
\frac{1}{\sqrt n}\,
\mathcal{J}(\delta,\widehat{\mathcal{G}})
\left(1+\frac{\mathcal{J}(\delta,\widehat{\mathcal{G}})}{\delta \sqrt n}\right) \\[0.4em]
&\;\lesssim\;
\frac{1}{\sqrt n}\,
\delta \sqrt{B \log(1/\delta)}
\left(
1 + 
\frac{\delta \sqrt{B \log(1/\delta)}}{\delta \sqrt n}
\right) \\[0.4em]
&\;\lesssim\;
\frac{\delta \sqrt{B \log(1/\delta)}}{\sqrt n}
\;+\;
\frac{B \log(1/\delta)}{n}.
\end{align*}
The critical radius
\[
\delta_n := \sqrt{\frac{B}{n}\,\log\Big(\frac{n}{B}\Big)},
\]
satisfies
\[
\delta_n
= \inf\{\delta > 0 : \mathcal{R}_n(\widehat{\mathcal{G}}, \delta)
\;\lesssim\; \delta^2\}.
\]

Applying Lemma~\ref{lemma:loc_max_ineq} conditional on the training data
with \(\mathcal{F} := \widehat{\mathcal{G}}\), we conclude that the following holds with probability at least
\(1 - e^{-u^2}\) for every \(f \in \widehat{\mathcal{G}}\):
\[
\frac{1}{n} \sum_{i=1}^n 
\big(f(O_i) - \mathbb{E}[f(O_i)]\big)
\;\lesssim\;
\delta_n^2
\;+\;
\delta_n\|f\|_{2,P_0}
\;+\;
\frac{u\|f\|_{2,P_0}}{\sqrt{n}}
\;+\;
\frac{u^2}{n}.
\]
Choosing \(u=\sqrt{\log(1/\eta)}\) gives \(1-e^{-u^2}=1-\eta\). Hence, with probability at least \(1-\eta\),
\begin{align*}
&\frac{1}{n} \sum_{i=1}^n 
\big(f(O_i) - \mathbb{E}[f(O_i)]\big)\\
&\lesssim
\delta_n^2
+ 
\delta_n\|f\|_{2,P_0}  +
\frac{\sqrt{\log(1/\eta)}\,\|f\|_{2,P_0}}{\sqrt{n}}
 + 
\frac{\log(1/\eta)}{n}.
\end{align*}





By boundedness of our nuisances,
\[
\big\|
\{\Gamma_0(V_K) - V_K\}
\{\widehat Y_{V_{K-1}} - V_K\}
\big\|
\;\lesssim\;
\|\Gamma_0(V_K) - V_K\|_{2,\rho}.
\]
Hence, with probability at least \(1 - \eta\),
\[
\begin{aligned}
(P_n - P_0)
&\Big\{
\Gamma_0(V_K) - V_K
\Big\}
\Big\{
\widehat Y_{V_{K-1}}
- V_K
\Big\} \\
&\lesssim\;
\delta_n^2
\;+\;
\delta_n \|\Gamma_0(V_K) - V_K\| \\[0.35em]
&\quad+\;
\frac{\sqrt{\log(1/\eta)}\,\|\Gamma_0(V_K) - V_K\|}{\sqrt{n}}
\;+\;
\frac{\log(1/\eta)}{n},
\end{aligned}
\]
where the implicit constants do not depend on \(B\).



Combining the above with \eqref{eqn::basicineq2}, we find with probability at least \(1 - \eta\),
\begin{equation}
\label{eqn::basicineq3}
\begin{aligned}
&\|\Gamma_0(V_K) - V_K\|^2\\
&\lesssim\;
\delta_n^2
\;+\;
\delta_n \|\Gamma_0(V_K) - V_K\|\\
&\quad+\;
\frac{\sqrt{\log(1/\eta)}}{\sqrt{n}}\,
\|\Gamma_0(V_K) - V_K\|
\;+\;
\frac{\log(1/\eta)}{n}
\\[0.4em]
&\quad+\;
\|\Gamma_0(V_K) - V_K\|\;
\big\|
\widehat T_\pi V_K
- \widehat T_\pi V_{K-1}
\big\|
\\[0.4em]
&\quad+\;
\|\Gamma_0(V_K) - V_K\|\;
\big\|
b_0\{(\widehat w_{\pi} - w_{\pi})\,
(\widehat q_{V_K} - q_{V_K})\}
\big\|_{2,\rho}.
\end{aligned}
\end{equation}
The inequality in \eqref{eqn::basicineq3} implies that, 
with probability at least $1-\eta$, the calibration error satisfies
\begin{align*}
\|\Gamma_0(V_K) - V_K\|
\;\lesssim\;&\;
\delta_n
\;+\;
\sqrt{\frac{\log(1/\eta)}{n}}
\\[0.4em]
&\;+\;
\big\|
\widehat T_\pi V_K
- \widehat T_\pi V_{K-1}
\big\|
\\[0.4em]
&\;+\;
\big\|
b_0\{(\widehat w_{\pi} - w_{\pi})\,
(\widehat q_{V_K} - q_{V_K})\}
\big\|_{2,\rho}.
\end{align*}
Recall that $\delta_n := \sqrt{\frac{B}{n}\,\log\Big(\frac{n}{B}\Big)}$. Then,  with probability at least $1-\eta$,
\begin{align*}
\|\Gamma_0(V_K) - V_K\|
\;\lesssim\;&\;
\sqrt{\frac{B}{n}\,\log\Big(\frac{n}{B}\Big)}
\;+\;
\sqrt{\frac{\log(1/\eta)}{n}}
\\[0.4em]
&\;+\;
\big\|
\widehat T_\pi V_K
- \widehat T_\pi V_{K-1}
\big\|
\\[0.4em]
&\;+\;
\big\|
b_0\{(\widehat w_{\pi} - w_{\pi})\,
(\widehat q_{V_K} - q_{V_K})\}
\big\|_{2,\rho}.
\end{align*}
\end{proof}


\section{Proof of Theorem \ref{theorem::regretHist}}

\subsection{Additional notation}
Define the bin–index map $\hat{b}(s)$ by setting $\hat{b}(s) = b$ whenever
$\hat{v}(s) \in I_b$.  Let
$\Pi_B$ denote the binning projection
\[
(\Pi_B g)(s)
:= E\!\left[g(S)\,\middle|\,\hat v(S)\in I_{\hat b(s)}\right].
\] Denote 
\[
\mathcal B_{\pi,B}(g)(s)
:= (\Pi_B \mathcal{T}_{\pi}g)(s).
\]
Define $V_B^\star$ as the fixed point satisfying
\[
\mathcal B_{\pi,B}(V_B^\star) = V_B^\star.
\]
Such a fixed point exists by Banach's fixed point theorem because 
$\mathcal B_{\pi,B}(g) \in \mathcal H_B(\hat v)$ whenever 
$g \in \mathcal H_B(\hat v)$, and the operator is a $\gamma$-contraction in the
sup-norm. 


Recall the projected
transition operator 
\[
(P_{\pi,B} h)(s)
=
\mathbb{E}_{\pi}\left[h(S') \,\middle|\, \hat{v}(S) \in I_{\hat{b}(s)} \right].
\]
It holds that
\[
P_{\pi,B}
\;:=\;
\Pi_B P_\pi,
\]
and
\[
\Pi_B\mathcal{T}_\pi(f)
-
\Pi_B\mathcal{T}_\pi(g)
=
\gamma\, P_{\pi,B}(f-g).
\]

\subsection{Supporting lemmas}
We begin with the deterministic ingredients for the histogram value bound.

Throughout, histogram bins mean the retained cells of the fitted partition:
empty empirical cells are merged with a neighboring cell or discarded, and
\(B\) denotes the number of retained cells. Population projections are defined
on the corresponding non-null cells; \(\rho\)-null cells are identified in
\(L^2(\rho)\) and omitted from the index set.




\begin{lemma}[Sup-norm contraction of the aggregated Bellman operator]
\label{lemma::agg-sup-contraction}
For any bounded \(f,g:\mathcal S\to\mathbb R\),
\[
\|\mathcal B_{\pi,B}f-\mathcal B_{\pi,B}g\|_\infty
\le
\gamma\|f-g\|_\infty .
\]
\end{lemma}

\begin{proof}
Since the rewards cancel on differences,
\[
\mathcal B_{\pi,B}f-\mathcal B_{\pi,B}g
=
\gamma P_{\pi,B}(f-g).
\]
The operator \(P_{\pi,B}\) is a Markov averaging
operator. Hence, for any bounded \(h\),
\[
|P_{\pi,B}h(s)|
\le
E_\pi[|h(S')|\mid \hat v(S)\in I_{\hat b(s)}]
\le
\|h\|_\infty .
\]
Taking \(h=f-g\) proves the claim.
\end{proof}

\begin{lemma}[Approximation error under aggregated concentrability]
\label{lemma::approxerror}
Assume \ref{cond::Ccoverage}. Then
\[
\|V_B^\star - v_0\|_{2,\rho}
\le
\frac{\sqrt{\mathcal C_{B,\gamma}}}{1-\gamma}\,
\|(I-\Pi_B)v_0\|_{2,\rho}.
\]
\end{lemma}

\begin{proof}
Because \(v_0=\mathcal T_\pi v_0\) and
\(V_B^\star=\Pi_B\mathcal T_\pi V_B^\star\),
\[
V_B^\star-v_0
=
\gamma P_{\pi,B}(V_B^\star-v_0)
+
(\Pi_B-I)v_0 .
\]
Let \(a:=(\Pi_B-I)v_0\). Iterating the display and using boundedness gives the
resolvent expansion
\[
V_B^\star-v_0
=
\sum_{t=0}^\infty \gamma^t P_{\pi,B}^t a .
\]
For any \(t\ge0\), Jensen's inequality for the Markov operator \(P_{\pi,B}^t\)
gives
\[
\|P_{\pi,B}^t a\|_{2,\rho}^2
\le
\int P_{\pi,B}^t(a^2)\,d\rho
=
\int a^2\,d(\rho P_{\pi,B}^t)
\le
c_{B,t}\|a\|_{2,\rho}^2 .
\]
Therefore,
\[
\|V_B^\star-v_0\|_{2,\rho}
\le
\sum_{t=0}^\infty \gamma^t\sqrt{c_{B,t}}\,\|a\|_{2,\rho}
=
\frac{\sqrt{\mathcal C_{B,\gamma}}}{1-\gamma}
\|(I-\Pi_B)v_0\|_{2,\rho}.
\]
\end{proof}


\begin{lemma}[Inexact iterations under aggregated concentrability]
\label{lemma::inexacterrors}
Assume \ref{cond::Ccoverage}. Let \(\{\eta_k\}_{k\ge1}\) be any sequence such
that
\[
\|V_k - \mathcal B_{\pi,B}(V_{k-1})\|_{2,\rho}
\le
\eta_k
\qquad\text{for all } k.
\]
Then, for any \(K\ge1\),
\[
\|V_K - V_B^\star\|_{2,\rho}
\le
\gamma^K\|\hat v - V_B^\star\|_\infty
+
\sum_{j=1}^K
\gamma^{K-j}\sqrt{c_{B,K-j}}\,\eta_j .
\]
Consequently,
\[
\|V_K - V_B^\star\|_{2,\rho}
\le
\gamma^K\|\hat v - V_B^\star\|_\infty
+
\frac{\sqrt{\mathcal C_{B,\gamma}}}{1-\gamma}
\max_{1\le j\le K}\eta_j .
\]
\end{lemma}

\begin{proof}
Let
\[
e_k:=V_k-\mathcal B_{\pi,B}(V_{k-1}).
\]
Since \(V_B^\star=\mathcal B_{\pi,B}V_B^\star\),
\[
V_k-V_B^\star
=
\gamma P_{\pi,B}(V_{k-1}-V_B^\star)+e_k.
\]

Iterating this recursion yields
\[
V_K-V_B^\star
=
\gamma^K P_{\pi,B}^K(\hat v-V_B^\star)
+
\sum_{j=1}^K\gamma^{K-j}P_{\pi,B}^{K-j}e_j .
\]
The first term is bounded by
\(\gamma^K\|\hat v-V_B^\star\|_\infty\) because \(P_{\pi,B}\) is a Markov
operator. For the remaining terms, the same Jensen and change-of-measure
argument used in Lemma~\ref{lemma::approxerror} gives
\[
\|P_{\pi,B}^{K-j}e_j\|_{2,\rho}
\le
\sqrt{c_{B,K-j}}\,\|e_j\|_{2,\rho}
\le
\sqrt{c_{B,K-j}}\,\eta_j .
\]
The displayed bounds follow by the triangle inequality and the definition of
\(\mathcal C_{B,\gamma}\).
\end{proof}


\begin{lemma}[Error bound for inexact updates]
\label{lemma::errorperiter}
There exists \(C<\infty\) such that, with probability at least \(1-\delta\),
uniformly over \(k=1,\ldots,K\),
 \begin{align*}
&\|V_k -
  \mathcal B_{\pi,B}(V_{k-1})\|_{2,\rho}\\
 & \qquad \leq
C\left\{
\sqrt{\frac{B}{n}\,\log\Big(\frac{n}{B}\Big)}
\;+\;
\sqrt{\frac{\log(K/\delta)}{n}}
+
\big\|
b_0\{(\widehat w_{\pi} - w_{\pi})\,
(\widehat q_{V_{k-1}} - q_{V_{k-1}})\}
\big\|_{2,\rho}
\right\}
\end{align*}
 \end{lemma}
\begin{proof}
Equation~\ref{eqn::scores} in the proof of Theorem~\ref{theorem::Calerrorhist} shows that, for any transformation $f \in \mathcal G_B$ and any $k$,
\begin{equation*}
\frac{1}{n}\sum_{i=1}^n 
f(\hat{v}(S_i))
\big\{
\widehat Y_{V_{k-1}}(O_i)
- V_k(S_i)
\big\}
= 0.
\end{equation*}
Applying the above with 
$f \circ \hat{v} = V_B^{\star(k)} - V_k$, 
where 
\[
V_B^{\star(k)} := \mathcal B_{\pi,B}(V_{k-1}),
\]
we obtain
\begin{equation*}
\frac{1}{n}\sum_{i=1}^n 
\big(V_B^{\star(k)} - V_k\big)(S_i)
\Big\{
\widehat Y_{V_{k-1}}(O_i)
-
V_k(S_i)
\Big\}
= 0.
\end{equation*}
Adding and subtracting $P_0$, we find
\begin{equation}
\begin{aligned}
&P_0\left[
  (V_B^{\star(k)} - V_k)
  \Big\{
    \widehat Y_{V_{k-1}}
    - V_k
  \Big\}
\right]\\
&=
(P_0 - P_n)\left[
  (V_B^{\star(k)} - V_k)
  \Big\{
    \widehat Y_{V_{k-1}}
    - V_k
  \Big\}
\right].
\end{aligned}
\label{eq:pn_p0_identity}
\end{equation}

First, we study the left–hand side of
\eqref{eq:pn_p0_identity}. We have
\begin{align*}
    & P_0\left[
  (V_B^{\star(k)} - V_k)
  \Big\{
    \widehat Y_{V_{k-1}}
    - V_k
  \Big\}
\right]\\
& = \|V_B^{\star(k)} - V_k\|_{2,\rho}^2\\
& \quad + P_0\left[
  (V_B^{\star(k)} - V_k)
  \Big\{
    \widehat Y_{V_{k-1}}
    - V_B^{\star(k)}
  \Big\}
\right].
\end{align*}
Since \(V_B^{\star(k)}=\Pi_B\mathcal T_\pi(V_{k-1})\) and
\(V_B^{\star(k)}-V_k\in\mathcal H_B(\hat v)\), the projection identity
gives
\[
P_0\left[
  (V_B^{\star(k)}-V_k)
  \{V_B^{\star(k)}-\mathcal T_\pi(V_{k-1})\}
\right]=0.
\]
Thus
\begin{align*}
& P_0\left[
  (V_B^{\star(k)} - V_k)
  \Big\{
    \widehat Y_{V_{k-1}}
    - V_B^{\star(k)}
  \Big\}
\right] \\[0.4em]
&=
P_0\left[
  (V_B^{\star(k)} - V_k)
  \Big\{
    \widehat Y_{V_{k-1}}
    - \mathcal T_{\pi}\!\left(V_{k-1}\right)
  \Big\}
\right].
\end{align*}
By Theorem~\ref{theorem::DRpseudo} and arguing as in the proof of
Theorem~\ref{theorem::Calerrorhist}, it follows that
\begin{align*}
&\left|
P_0\left[
  (V_B^{\star(k)} - V_k)
  \Big\{
    \widehat Y_{V_{k-1}}
    - V_B^{\star(k)}
  \Big\}
\right]
\right|\\
&\;\le\;
\|V_B^{\star(k)} - V_k\|_{2,\rho}\;
\big\|
b_0\{(\widehat w_{\pi} - w_{\pi})\,
(\widehat q_{V_{k-1}} - q_{V_{k-1}})\}
\big\|_{2,\rho}.
\end{align*}
Putting it all together,
\begin{align*}
&\|V_B^{\star(k)} - V_k\|_{2,\rho}^2
 \leq
P_0\left[
  (V_B^{\star(k)} - V_k)
  \Big\{
    \widehat Y_{V_{k-1}}
    - V_k
  \Big\}
\right] \\
&+ 
\|V_B^{\star(k)} - V_k\|_{2,\rho}\;
\big\|
b_0\{(\widehat w_{\pi} - w_{\pi})\,
(\widehat q_{V_{k-1}} - q_{V_{k-1}})\}
\big\|_{2,\rho}.
\end{align*}
Hence, by \eqref{eq:pn_p0_identity},
\begin{equation}
\begin{aligned}
&\|V_B^{\star(k)} - V_k\|_{2,\rho}^2 \\
& \quad \leq
(P_0 - P_n)\left[
  (V_B^{\star(k)} - V_k)
  \Big\{
    \widehat Y_{V_{k-1}}
    - V_k
  \Big\}
\right]\\
& \quad + 
\|V_B^{\star(k)} - V_k\|_{2,\rho}\;
\big\|
b_0\{(\widehat w_{\pi} - w_{\pi})\,
(\widehat q_{V_{k-1}} - q_{V_{k-1}})\}
\big\|_{2,\rho}.
\end{aligned}
\label{eq:pn_p0_identity2}
\end{equation}


Next, we obtain a high–probability bound for the first term right–hand side of
\eqref{eq:pn_p0_identity2}.
Applying Lemmas~\ref{lemma:loc_max_ineq}, \ref{lemma:local_rademacher_entropy},
and \ref{lemma::entropynumbers}, and arguing as in the proof of
Theorem~\ref{theorem::Calerrorhist}, we obtain that, with probability at least
\(1 - \delta\),
\[
\begin{aligned}
\Big|(P_0 - P_n)
&\Big[
(V_B^{\star(k)} - V_k)
\big\{
\widehat Y_{V_{k-1}}
-
V_k
\big\}
\Big]\Big|
\\[0.4em]
&\lesssim\;
\delta_n^2
\;+\;
\delta_n\,\|V_B^{\star(k)} - V_k\|_{2,\rho}
\\[0.4em]
&\quad+\;
\frac{\sqrt{\log(1/\delta)}\,\|V_B^{\star(k)} - V_k\|_{2,\rho}}{\sqrt{n}}
\;+\;
\frac{\log(1/\delta)}{n},
\end{aligned}
\]
where \(\delta_n := \sqrt{(B/n)\,\log(n/B)}\), and the implicit constants do
not depend on \(B\).


Plugging this high probability bound into \eqref{eq:pn_p0_identity2},  we obtain that, with probability at least
\(1 - \delta\),
\begin{equation*}
\begin{aligned}
&\|V_B^{\star(k)} - V_k\|_{2,\rho}^2 \\
&\lesssim\;
\delta_n^2
\;+\;
\delta_n\,\|V_B^{\star(k)} - V_k\|_{2,\rho}
\\[0.4em]
&\quad+\;
\frac{\sqrt{\log(1/\delta)}\,\|V_B^{\star(k)} - V_k\|_{2,\rho}}{\sqrt{n}}
\;+\;
\frac{\log(1/\delta)}{n}\\
& \quad + 
\|V_B^{\star(k)} - V_k\|_{2,\rho}\;
\big\|
b_0\{(\widehat w_{\pi} - w_{\pi})\,
(\widehat q_{V_{k-1}} - q_{V_{k-1}})\}
\big\|_{2,\rho}.
\end{aligned}
\end{equation*}
Recalling that $\delta_n := \sqrt{\frac{B}{n}\,\log\Big(\frac{n}{B}\Big)}$, the inequality above implies that, 
with probability at least $1-\delta$,  
 \begin{align*}
\|V_B^{\star(k)} - V_k\|_{2,\rho}
\;\lesssim\;&\;
\sqrt{\frac{B}{n}\,\log\Big(\frac{n}{B}\Big)}
\;+\;
\sqrt{\frac{\log(1/\delta)}{n}}
\\[0.4em]
&  + 
\big\|
b_0\{(\widehat w_{\pi} - w_{\pi})\,
(\widehat q_{V_{k-1}} - q_{V_{k-1}})\}
\big\|_{2,\rho}.
\end{align*}
Taking the failure probability to be \(\delta/K\) and applying a union bound over
\(k=1,\ldots,K\) proves the result.
\end{proof}

\subsection{Proof of Theorem \ref{theorem::regretHist}}

\begin{proof}[Proof of Theorem \ref{theorem::regretHist}]
For \(C<\infty\) large enough, define
\begin{align*}
  \eta_k
&=
C \Bigg\{\sqrt{\frac{B}{n}\,\log\!\Big(\frac{n}{B}\Big)}
\;+\;
\sqrt{\frac{\log(K/\delta)}{n}}
\\
& \quad
\;+\;
\big\|
b_0\{(\widehat w_{\pi} - w_{\pi})\,
(\widehat q_{V_{k-1}} - q_{V_{k-1}})\}
\big\|_{2,\rho}
\Bigg\}
\end{align*}

By Lemma~\ref{lemma::errorperiter}, with probability at least \(1-\delta\),
uniformly over \(k=1,\ldots,K\),
\[
\|V_k-\mathcal B_{\pi,B}(V_{k-1})\|_{2,\rho}
\le
\eta_k .
\]
Lemma~\ref{lemma::inexacterrors} therefore gives
\begin{align*}
\|V_K - V_B^\star\|_{2,\rho}
\le\;&
\gamma^K\|\hat v - V_B^\star\|_\infty\\
&+
\frac{\sqrt{\mathcal C_{B,\gamma}}}{1-\gamma}
\max_{1\le j\le K}\eta_j .
\end{align*}
Combining this with Lemma~\ref{lemma::approxerror} and the triangle inequality
yields
\begin{align*}
\|V_K-v_0\|_{2,\rho}
\le\;&
\gamma^K\|\hat v - V_B^\star\|_\infty\\
&+
\frac{\sqrt{\mathcal C_{B,\gamma}}}{1-\gamma}
\left[
\|(I-\Pi_B)v_0\|_{2,\rho}
+
\max_{1\le j\le K}\eta_j
\right].
\end{align*}
Since \(\Pi_B\) is the \(L^2(\rho)\) projection onto
\(\mathcal H_B(\hat v)\),
\[
\|(I-\Pi_B)v_0\|_{2,\rho}
=
\min_{h\in\mathcal H_B(\hat v)}\|h-v_0\|_{2,\rho}.
\]
Substituting the definition of \(\eta_k\) and absorbing constants completes the
proof.
\end{proof}






\subsection{Bound on Successive Iteration Errors}
\label{appendix::convergence}
\begin{lemma}[Bound on Successive Iteration Errors]
Under the conditions of Theorem~\ref{theorem::regretHist},
\begin{align*}
&\|V_K - V_{K-1}\|_{2,\rho}\\
&\lesssim
\gamma^{K-1}\,
\|\hat v - V_B^\star\|_{\infty}
\\[0.8em]
&\quad+\;
\frac{\sqrt{\mathcal C_{B,\gamma}}}{1-\gamma}
C\Bigg(
\sqrt{\frac{B}{n}\,\log\!\Big(\frac{n}{B}\Big)}
\;+\;
\sqrt{\frac{\log(K/\delta)}{n}}
\\[0.8em]
&\qquad\qquad+
\max_{1\le j\le K}
\Big\|
(\widehat w_{\pi} - w_{\pi})\,
(\widehat q_{V_{j-1}} - q_{V_{j-1}})
\big\|_{2,\rho}
\Bigg).
\end{align*}
\end{lemma}

\begin{proof}
Lemma~\ref{lemma::inexacterrors} and Lemma~\ref{lemma::errorperiter} give the
same bound for \(\|V_m-V_B^\star\|_{2,\rho}\), for \(m=K-1\) and \(m=K\), with
\(\gamma^m\|\hat v-V_B^\star\|_\infty\) as the finite-iteration term. By the
triangle inequality,
\begin{align*}
\|V_K - V_{K-1}\|_{2,\rho}
&\le\;
\|V_{K-1} - V_B^\star\|_{2,\rho}
\\[0.25em]
&\quad+\;
\|V_K - V_B^\star\|_{2,\rho}.
\end{align*}
Substituting the two bounds and absorbing constants proves the claim.
\end{proof}



\section{Proofs of Theorems \ref{theorem::Calerroriso} and \ref{theorem::regretiso}}


\begin{proof}[Proof of Theorem \ref{theorem::Calerroriso}]
The first-order optimality conditions for isotonic regression (equivalently, its
interpretation as a histogram estimator) imply that for any function
$f:\mathbb{R}\rightarrow\mathbb{R}$ that is a linear combination of the
indicator functions defining the isotonic partition,
\begin{equation*}
\frac{1}{n}\sum_{i=1}^n
f\!\big(V_K(S_i)\big)
\Big\{
\widehat Y_{V_{K-1}}(O_i)
-
V_K(S_i)
\Big\}
= 0.
\end{equation*}
See, for example, the proof of Lemma~C.1 in \citet{van2023causal} and
\citet{van2024stabilized}. Choosing $f$ appropriately yields
\begin{align*}
0
&=
\frac{1}{n}\sum_{i=1}^n
\Big\{
\Gamma_0(V_K)(S_i) - V_K(S_i)
\Big\} \\
&\qquad\qquad\times
\Big\{
\widehat Y_{V_{K-1}}(O_i)
- V_K(S_i)
\Big\},
\end{align*}
which is the same basic equality as \eqref{eqn::basicineq} in the proof of
Theorem~\ref{theorem::Calerrorhist}. The remainder of the argument proceeds
along the same lines with minor modifications.

Specifically, let $\mathcal{F}_{TV}$ denote the union of $\mathcal G_{\mathrm{iso}}$,
which is uniformly bounded by $2M$ under Condition~\ref{cond::C1}, with all
functions of bounded total variation bounded by the constant $C$ in
Condition~\ref{cond::C3}. By \citet{van1996weak}, this class satisfies the
uniform entropy integral bound $\mathcal{J}(\delta,\mathcal{F}_{TV})
\lesssim \sqrt{\delta}$.

Let
\[
m_K(t)
:=
E\!\left[(\mathcal T_\pi V_K)(S)\mid \hat v(S)=t,\mathcal C_n\right].
\]
Condition~\ref{cond::C3} states that \(m_K\) has bounded total variation. Since
\[
\Gamma_0(V_K)(S)=E\{m_K(\hat v(S))\mid V_K(S)\},
\]
and \(V_K\) is a nondecreasing function of \(\hat v\), Lemma 6 of
\citet{van2024stabilized} implies that the coarsened curve
\(\Gamma_0(V_K)\) also has bounded total variation and lies in
\(\mathcal{F}_{TV, \hat{v}} := \{f \circ \hat{v}: f \in
\mathcal{F}_{TV}\}\). Thus, it holds that
\[
\big(\Gamma_0(V_K) - V_K\big)
\big(\widehat Y_{V_{K-1}} - V_K\big)
\]
lies in a uniformly bounded subset of the class
\[
\widehat{\mathcal{G}}
:=
\Big\{
(f_1 - f_2)\big(\widehat Y_{f_2} - f_2\big)
:\;
f_1, f_2 \in \mathcal{F}_{TV, \hat{v}}
\Big\}.
\]
Arguing as in the proof of Theorem~\ref{theorem::Calerrorhist}, we have
\[
\mathcal{J}(\delta, \widehat{\mathcal{G}})
\lesssim
\mathcal{J}(\delta, \mathcal{F}_{TV})
\lesssim
\sqrt{\delta}.
\]
The proof now follows directly from Theorem~\ref{theorem::Calerrorhist} with this
new choice of $\widehat{\mathcal{G}}$ and its associated critical radius
$\delta_n = n^{-1/3}$ for monotone functions.
\end{proof}

 \begin{proof}[Proof of Theorem \ref{theorem::regretiso}]
The result follows directly from a union bound and the proofs of
Theorems~\ref{theorem::Calerrorhist} and \ref{theorem::regretHist}, with
$\mathcal G_B$ replaced by $\mathcal G_{B_n}$, the class of all
piecewise-constant functions with at most $B_n$ constant segments. In
particular, the entropy bound in Lemma~\ref{lemma::entropynumbers} continues to
apply to this class. Notably, the proofs of these theorems allow for
data-adaptive partitions and require only a deterministic upper bound on the
number of constant segments.
\end{proof}
 

\end{document}
